% Réactions d'oxydoréduction spontanées et nombres d'oxydation — Chimie Tle TI — Cours signé OnBuch+
% Programme officiel MINESEC, Terminale TI. Compiler avec : tectonic L09-reactions-oxydoreduction-spontanees-nombres-oxydation.tex
\newcommand{\DOCMATIERE}{Chimie}
\newcommand{\DOCNIVEAU}{Tle TI}
\newcommand{\DOCMODULE}{Module 3 · Oxydoréduction}
\newcommand{\DOCLECON}{09}
\newcommand{\DOCTITRE}{Réactions d'oxydoréduction spontanées et nombres d'oxydation}
% OnBuch+ — préambule « cours signés » ENRICHI (compilé avec tectonic / XeLaTeX)
% Système visuel « Pop » d'OnBuch+ : fond crème (jamais blanc), bordures encre
% épaisses, ombres « dures » décalées, titres Archivo Black en capitales.
% Version enrichie pour les cours de chimie : chimie (mhchem, chemfig),
% graphes (pgfplots), boîtes pédagogiques supplémentaires (définition,
% propriété, méthode, attention, le savais-tu, exercice résolu, exercice,
% TP, bilan), page de garde refaite et signature OnBuch+ ancrée en bas.
%
% Macros attendues dans main.tex AVANT \input{preamble} :
%   \DOCMATIERE  \DOCNIVEAU  \DOCMODULE  \DOCLECON (numéro)  \DOCTITRE
\documentclass[11pt]{article}

\usepackage{fontspec}
\usepackage[french]{babel}
\usepackage[a4paper,margin=2cm,top=3.4cm,bottom=2.8cm,headheight=1.4cm,footskip=1.3cm]{geometry}
\usepackage{amsmath,amssymb,amsfonts}
\usepackage[table]{xcolor} % [table] : fournit aussi \rowcolors (alternance de lignes)
\usepackage{pagecolor}
\usepackage{tikz}
\usetikzlibrary{arrows.meta,positioning,calc,shapes.geometric,shapes.misc,shapes.arrows,decorations.pathmorphing,decorations.markings,patterns,shadows,mindmap,trees,fit,backgrounds,matrix,intersections}
\usepackage{pgfplots}
\pgfplotsset{compat=1.18}
\usepackage[version=4]{mhchem}
\usepackage{chemfig}
\usepackage{enumitem}
\usepackage{fancyhdr}
% [most] charge toutes les bibliothèques tcolorbox (listings, minted, raster,
% documentation...) et gonfle inutilement la mémoire TeX sur un document long
% (~300 boîtes) : on ne charge que ce qui est réellement utilisé.
\usepackage[skins,breakable]{tcolorbox}
\usepackage{graphicx}
\usepackage{colortbl}
\usepackage{booktabs}
\usepackage{tabularx}
\usepackage{array}
\usepackage{multicol}
\usepackage{siunitx}
% parse-numbers=false : accepte aussi \SI{2,5 \times 10^{-3}}{...}
\sisetup{locale=FR,output-decimal-marker={,},group-minimum-digits=4,parse-numbers=false}
\DeclareSIUnit\ohm{\ensuremath{\symup{\Omega}}} % U+2126 absent de la police
\usepackage{qrcode}
% \degree : commande fréquemment utilisée par les modèles pour un angle ou une
% température en dehors de \SI{}{} (non fournie par les paquets chargés).
\providecommand{\degree}{^{\circ}}
\usepackage{lastpage}
\usepackage{hyperref}
\hypersetup{hidelinks}

% ============================================================
% Palette « Pop » OnBuch+ — mêmes hex que la classe P (plus_ui.dart)
% ============================================================
\definecolor{poporange}{HTML}{F59321}
\definecolor{popdark}{HTML}{E07A0C}
\definecolor{popcream}{HTML}{FFF6E8}
\definecolor{popink}{HTML}{1C1714}
\definecolor{poppurple}{HTML}{7A5AE0}
\definecolor{popgreen}{HTML}{2E9E6A}
\definecolor{poppink}{HTML}{E0457B}
\definecolor{popblue}{HTML}{2F7DE1}
\definecolor{popgold}{HTML}{C98A12}
\definecolor{popmuted}{HTML}{8A7F76}
\definecolor{popline}{HTML}{EADFD2}
\definecolor{popgray}{HTML}{8A7F76}
\colorlet{popgrey}{popgray}
% Alias tolérants (le modèle nomme parfois les couleurs différemment)
\colorlet{popwhite}{white}
\colorlet{popblack}{popink}
\colorlet{popred}{poppink}
\colorlet{popyellow}{popgold}
% Teintes claires pour fonds de tableaux / zones
\colorlet{popblueL}{popblue!10!white}
\colorlet{poporangeL}{poporange!14!white}
\colorlet{popgreenL}{popgreen!12!white}
\colorlet{poppurpleL}{poppurple!12!white}
\colorlet{poppinkL}{poppink!12!white}
\colorlet{popgoldL}{popgold!14!white}

\pagecolor{popcream}

% ============================================================
% Typographie
% ============================================================
\defaultfontfeatures{Ligatures=TeX}
\setmainfont{PlusJakartaSans-Medium.ttf}[Path=../../../fonts/,BoldFont=PlusJakartaSans-Bold.ttf]
% Maths en sans-serif (Fira Math) avec chiffres et lettres droites de Plus
% Jakarta, pour que les formules se fondent dans le texte.
\usepackage[math-style=ISO,bold-style=ISO]{unicode-math}
\setmathfont{FiraMath-Regular.otf}[Path=../../../fonts/]
\setmathfont{PlusJakartaSans-Medium.ttf}[Path=../../../fonts/,range={up/num,up/{Latin,latin},\mathup}]
\usepackage{icomma} % virgule décimale sans espace en mode math
% Caractères absents de la police de texte (sinon carré vide) : rendus par
% leur équivalent mathématique, en texte comme en formule.
\usepackage{newunicodechar}
\newunicodechar{α}{\ensuremath{\alpha}}
\newunicodechar{Α}{\ensuremath{\alpha}}
\newunicodechar{β}{\ensuremath{\beta}}
\newunicodechar{δ}{\ensuremath{\delta}}
\newunicodechar{✓}{\popcheck{popgreen}}
\newfontfamily\popdisplay{ArchivoBlack-Regular.ttf}[Path=../../../fonts/]
\newfontfamily\poplogo{SpaceGrotesk-Bold.ttf}[Path=../../../fonts/]
\newfontfamily\popsemi{PlusJakartaSans-SemiBold.ttf}[Path=../../../fonts/]
\newfontfamily\popxbold{PlusJakartaSans-ExtraBold.ttf}[Path=../../../fonts/]
\newfontfamily\popxboldls{PlusJakartaSans-ExtraBold.ttf}[Path=../../../fonts/,LetterSpace=3.0]

\setlist[enumerate]{leftmargin=1.7em,itemsep=5pt,topsep=4pt}
\setlist[itemize]{leftmargin=1.7em,itemsep=4pt,topsep=4pt,label={\color{poporange}\textbullet}}
\setlength{\parindent}{0pt}
\setlength{\parskip}{5pt}
\renewcommand{\arraystretch}{1.25}

% chemfig : liaisons un peu plus épaisses, encre Pop
\setchemfig{atom sep=2em,bond style={line width=0.9pt,popink},cram width=3pt}

% pgfplots : style Pop par défaut
\pgfplotsset{
  every axis/.append style={axis line style={popink,line width=1pt},tick style={popink},
    grid style={popline,line width=0.5pt},label style={font=\small},tick label style={font=\footnotesize}},
  popcurve/.style={poporange,line width=1.8pt,no markers,smooth},
}

% ============================================================
% Pill / squiggle
% ============================================================
\newcommand{\poppill}[3]{%
  \tikz[baseline=-0.55ex]\node[draw=popink,line width=1.2pt,rounded corners=2.6pt,%
    fill=#2,inner xsep=6.5pt,inner ysep=3pt]{%
    \popxboldls\fontsize{7}{7}\selectfont\color{#3}\MakeUppercase{#1}};%
}
\newcommand{\popsquiggle}[1]{%
  \tikz[baseline=-0.4ex]{
    \draw[#1,line width=2.6pt,line cap=round]
      plot [smooth] coordinates {(0,0.09) (0.34,0.01) (0.68,0.21) (1.02,0.01) (1.36,0.10)};
  }%
}
% Mot-clé surligné (marqueur orange sous le texte)
\newcommand{\cle}[1]{{\popxbold\color{popdark}#1}}

% ============================================================
% En-tête / pied
% ============================================================
\pagestyle{fancy}
\fancyhf{}
\renewcommand{\headrulewidth}{0pt}
\renewcommand{\footrulewidth}{0pt}
\lhead{\raisebox{-0.15ex}{\poplogo\fontsize{13}{13}\selectfont\color{popink}OnBuch\color{poporange}+}}
\rhead{\poppill{\DOCMATIERE\ \textbullet\ \DOCNIVEAU\ \textbullet\ Leçon \DOCLECON}{white}{popink}}
\lfoot{\popxbold\fontsize{7.6}{7.6}\selectfont\color{popmuted}\MakeUppercase{Cours signé OnBuch+}\ \textbullet\ {\popsemi\DOCTITRE}}
\rfoot{\tikz[baseline=-0.6ex]\node[rounded corners=8.5pt,fill=popink,minimum height=17pt,minimum width=17pt,inner xsep=5pt,inner sep=0]{\popxbold\fontsize{8}{8}\selectfont\color{popcream}\thepage\,/\,\pageref*{LastPage}};}

% ============================================================
% Titres
% ============================================================
\newcommand{\chaptitre}[2]{%
  \begin{tcolorbox}[enhanced,colback=poporange,colframe=popink,boxrule=2.6pt,arc=11pt,
      drop shadow=popink,left=15pt,right=15pt,top=12pt,bottom=11pt,before skip=3pt,after skip=18pt]
    \poppill{#2}{popink}{popcream}\par\vspace{9pt}
    {\raggedright\hyphenpenalty=10000\exhyphenpenalty=10000\popdisplay\fontsize{22}{24}\selectfont\color{popcream}\MakeUppercase{#1}\par}
  \end{tcolorbox}
}
% Section numérotée : pastille orange avec le numéro + titre Archivo
\newcounter{popsec}
\newcommand{\coursec}[1]{%
  \stepcounter{popsec}\setcounter{popsub}{0}%
  \par\vspace{16pt}\noindent
  \begin{minipage}{\linewidth}
  \tikz[baseline=-0.7ex]\node[draw=popink,line width=1.6pt,fill=poporange,rounded corners=4pt,minimum size=22pt,inner sep=0]{\popdisplay\fontsize{12}{12}\selectfont\color{popcream}\thepopsec};%
  \hspace{8pt}{\popdisplay\fontsize{15}{17}\selectfont\color{popink}\MakeUppercase{#1}}\par
  \vspace{4pt}\noindent{\color{poporange}\rule{36pt}{3.6pt}}
  \end{minipage}\par\nopagebreak\vspace{8pt}
}
\newcounter{popsub}
\newcommand{\courssub}[1]{%
  \stepcounter{popsub}%
  \par\vspace{9pt}\noindent{\popxbold\fontsize{11.5}{13}\selectfont\color{popink}{\color{poporange}\thepopsec.\thepopsub}\hspace{6pt}#1}\par\nopagebreak\vspace{4pt}
}

% ============================================================
% Boîtes pédagogiques — même recette : fond blanc, bordure épaisse colorée,
% ombre dure encre, badge plein. Titre optionnel [#1].
% ============================================================
\tcbset{popbase/.style={breakable,enhanced,colback=white,boxrule=2.3pt,arc=9pt,
  shadow={3pt}{-3pt}{0pt}{popink},left=14pt,right=14pt,top=11pt,bottom=11pt,before skip=11pt,after skip=15pt,
  fontupper=\popsemi\fontsize{10.6}{14.5}\selectfont\color{popink},
  fontlower=\popsemi\fontsize{10.6}{14.5}\selectfont\color{popink}}}
% Coche dessinée (le glyphe ✓ n'existe pas dans les polices du document)
\newcommand{\popcheck}[1]{\tikz[baseline=-0.5ex]\draw[#1,line width=1.6pt,line cap=round,line join=round](-0.13,0.0)--(-0.04,-0.09)--(0.13,0.1);}
\newcommand{\popboxhead}[3]{\poppill{#1}{#2}{white}\ifx\relax#3\relax\else\hspace{6pt}{\popxbold\fontsize{10.6}{12}\selectfont\color{popink}#3}\fi\par\vspace{7pt}}

\newtcolorbox{aretenir}[1][]{popbase,colframe=popgreen,before upper={\popboxhead{À retenir}{popgreen}{#1}}}
\newtcolorbox{exemplebox}[1][]{popbase,colframe=poporange,before upper={\popboxhead{Exemple}{poporange}{#1}}}
\newtcolorbox{definition}[1][]{popbase,colframe=popblue,before upper={\popboxhead{Définition}{popblue}{#1}}}
\newtcolorbox{propriete}[1][]{popbase,colframe=poppurple,before upper={\popboxhead{Propriété}{poppurple}{#1}}}
\newtcolorbox{methode}[1][]{popbase,colframe=popgold,before upper={\popboxhead{Méthode}{popgold}{#1}}}
\newtcolorbox{attention}[1][]{popbase,colframe=poppink,before upper={\popboxhead{Attention}{poppink}{#1}}}
\newtcolorbox{experience}[1][]{popbase,colframe=popblue,colback=popblueL,before upper={\popboxhead{Expérience}{popblue}{#1}}}
% « Le savais-tu ? » : carte violette pleine (contexte, culture, Cameroun)
\newtcolorbox{savaistu}[1][]{popbase,colback=poppurple,colframe=popink,
  fontupper=\popsemi\fontsize{10.4}{14.2}\selectfont\color{white},
  before upper={\poppill{Le savais-tu\,?}{white}{poppurple}\ifx\relax#1\relax\else\hspace{6pt}{\popxbold\color{white}#1}\fi\par\vspace{7pt}}}
% Exercice résolu : énoncé en haut, \tcblower puis la solution sur fond crème
\newcounter{popexo}\newcounter{popexr}
\newtcolorbox{exoresolu}[1][]{popbase,colframe=poporange,colbacklower=popcream,
  segmentation style={popink,line width=1.2pt,dashed},
  before upper={\stepcounter{popexr}\popboxhead{Exercice résolu \thepopexr}{poporange}{#1}},
  before lower={\poppill{Solution}{popink}{popcream}\par\vspace{6pt}}}
% Exercice (non résolu en place) — la correction est dans la partie Corrigés
\newtcolorbox{exercice}[1][]{popbase,colframe=popink,
  before upper={\stepcounter{popexo}\popboxhead{Exercice \thepopexo}{popink}{#1}}}
% Corrigé
\newtcolorbox{corrige}[1][]{popbase,colframe=popgreen,colback=popgreenL,
  before upper={\popboxhead{Corrigé}{popgreen}{#1}}}
% Objectifs / prérequis / bilan
\newtcolorbox{objectifs}{popbase,colframe=popink,colback=poporangeL,
  before upper={\popboxhead{Objectifs de la leçon}{popink}{}}}
\newtcolorbox{prerequis}{popbase,colframe=popink,colback=popblueL,
  before upper={\popboxhead{Prérequis}{popblue}{}}}
\newtcolorbox{bilan}{popbase,colframe=popink,colback=white,boxrule=2.8pt,
  before upper={\popboxhead{Fiche bilan}{poporange}{L'essentiel à connaître pour l'examen}}}
% Figure encadrée avec légende
\newtcolorbox{popfigure}[1][]{enhanced,breakable=false,colback=white,colframe=popline,boxrule=1.4pt,arc=8pt,
  left=10pt,right=10pt,top=10pt,bottom=8pt,before skip=10pt,after skip=12pt,halign=center,
  lower separated=false,halign lower=center,
  fontlower=\popsemi\fontsize{9}{11.5}\selectfont\color{popmuted},
  #1}
\newcommand{\legende}[1]{\par\vspace{6pt}{\popsemi\fontsize{9}{11.5}\selectfont\color{popmuted}#1}}

% ============================================================
% Signature OnBuch+
% ============================================================
\newcommand{\poplogoinline}[1]{{\poplogo\fontsize{#1}{#1}\selectfont\color{popink}OnBuch\color{poporange}+}}
\newcommand{\signatureonbuch}{%
  \begin{tcolorbox}[enhanced,colback=popink,colframe=popink,boxrule=2.2pt,arc=10pt,
      drop shadow=poporange,left=14pt,right=14pt,top=10pt,bottom=10pt,before skip=0pt,after skip=0pt]
    \tikz[baseline=-0.7ex]\node[circle,fill=popgreen,draw=popcream,line width=1.3pt,minimum size=22pt,inner sep=0]{\popcheck{white}};%
    \hspace{9pt}\begin{minipage}[c]{0.62\linewidth}
      {\popxbold\fontsize{12}{13}\selectfont\color{popcream}Signé\ \textbullet\ {\poplogo OnBuch\color{poporange}+}}\par\vspace{2pt}
      {\raggedright\popsemi\fontsize{8}{10.5}\selectfont\color{popline}\MakeUppercase{Équipe pédagogique OnBuch+}\\\MakeUppercase{Conforme au programme officiel MINESEC}\par}
    \end{minipage}\hfill
    {\poplogo\fontsize{22}{22}\selectfont\color{popcream}OnBuch\color{poporange}+}
  \end{tcolorbox}%
}

% ============================================================
% Page de garde
% #1 titre · #2 matière · #3 niveau · #4 module · #5 n° leçon · #6 durée ·
% #7 description / objectifs (texte ou itemize)
% ============================================================
\newcommand{\pagegarde}[7]{%
  \thispagestyle{empty}
  \newgeometry{a4paper,margin=1.7cm,top=1.6cm,bottom=1.4cm}%
  \begin{tikzpicture}[remember picture,overlay]
    \fill[poporange!18] ([xshift=-3.2cm,yshift=-3.6cm]current page.north east) circle (4.6cm);
    \fill[poppurple!14] ([xshift=2.4cm,yshift=7.6cm]current page.south west) circle (3.8cm);
  \end{tikzpicture}%
  {\poplogo\fontsize{30}{30}\selectfont\color{popink}OnBuch\color{poporange}+}\hfill
  \raisebox{0.6ex}{\poppill{Cours signé \textbullet\ Premium}{popink}{popcream}}\par
  \vspace{1pt}\popsquiggle{poppurple}\par
  \vspace{8pt}
  \begin{tcolorbox}[enhanced,colback=poporange,colframe=popink,boxrule=2.8pt,arc=13pt,
      drop shadow=popink,left=18pt,right=18pt,top=12pt,bottom=13pt,after skip=10pt]
    \poppill{#2 \textbullet\ #3}{popink}{popcream}\hspace{5pt}\poppill{#4}{white}{popink}\par\vspace{8pt}
    {\popdisplay\fontsize{14}{14}\selectfont\color{popink}LEÇON #5}\par\vspace{4pt}
    {\raggedright\hyphenpenalty=10000\exhyphenpenalty=10000\popdisplay\fontsize{25}{27}\selectfont\color{popcream}\MakeUppercase{#1}\par}
    \vspace{8pt}
    {\popxbold\fontsize{9.5}{9.5}\selectfont\color{popink}\MakeUppercase{Durée conseillée : #6}}
  \end{tcolorbox}
  \begin{tcolorbox}[enhanced,colback=white,colframe=popink,boxrule=2.2pt,arc=10pt,
      drop shadow=popink,left=16pt,right=16pt,top=10pt,bottom=10pt,after skip=8pt]
    \poppill{Dans cette leçon}{poppurple}{white}\par\vspace{5pt}
    \popsemi\fontsize{9.3}{12.6}\selectfont\color{popink}#7
  \end{tcolorbox}
  \noindent\begin{minipage}[t]{0.487\linewidth}
    \begin{tcolorbox}[enhanced,colback=popgreen,colframe=popink,boxrule=2.2pt,arc=10pt,
        drop shadow=popink,left=11pt,right=11pt,top=10pt,bottom=10pt,height=3.1cm,valign=center]
      \tcbox[colback=white,colframe=popink,boxrule=1.3pt,arc=5pt,boxsep=0pt,nobeforeafter,
        left=3pt,right=3pt,top=3pt,bottom=3pt,valign=center]{\qrcode[height=1.9cm]{https://play.google.com/store/apps/details?id=com.luvvix.onbuch&hl=fr}}%
      \hspace{8pt}\begin{minipage}[c]{0.5\linewidth}
        {\popdisplay\fontsize{11}{12.5}\selectfont\color{white}\MakeUppercase{Télécharge OnBuch}}\par\vspace{4pt}
        {\raggedright\popsemi\fontsize{8.4}{11}\selectfont\color{white}Gratuit \textbullet\ Google Play\\Tuteur IA, annales, concours, résultats\par}
      \end{minipage}
    \end{tcolorbox}
  \end{minipage}\hfill
  \begin{minipage}[t]{0.487\linewidth}
    \begin{tcolorbox}[enhanced,colback=poppurple,colframe=popink,boxrule=2.2pt,arc=10pt,
        drop shadow=popink,left=11pt,right=11pt,top=10pt,bottom=10pt,height=3.1cm,valign=center]
      \tikz[baseline=-0.7ex]\node[circle,fill=white,draw=popink,line width=1.3pt,minimum size=1.6cm,inner sep=0]{\popdisplay\fontsize{24}{24}\selectfont\color{poppurple}+};%
      \hspace{8pt}\begin{minipage}[c]{0.6\linewidth}
        {\popdisplay\fontsize{11}{12.5}\selectfont\color{white}\MakeUppercase{Passe à OnBuch+}}\par\vspace{4pt}
        {\raggedright\popsemi\fontsize{8.4}{11}\selectfont\color{white}Tous les cours signés, Tuteur IA illimité, fiches PDF — onglet OnBuch+ de l'app\par}
      \end{minipage}
    \end{tcolorbox}
  \end{minipage}
  \vspace{8pt}
  \popcontenu
  \vfill
  \signatureonbuch
  \restoregeometry
  \newpage
}

% Rangée « Ce cours contient » (page de garde)
\newcommand{\poptile}[3]{% #1 couleur #2 gros texte #3 légende
  \begin{tcolorbox}[enhanced,colback=#1,colframe=popink,boxrule=2pt,arc=9pt,shadow={2.5pt}{-2.5pt}{0pt}{popink},
      left=6pt,right=6pt,top=9pt,bottom=9pt,halign=center,valign=center,height=1.75cm,width=\linewidth,nobeforeafter]
    {\popdisplay\fontsize{12.5}{13}\selectfont\color{white}\MakeUppercase{#2}}\par\vspace{3pt}
    {\centering\popsemi\fontsize{7.8}{9.5}\selectfont\color{white}#3\par}
  \end{tcolorbox}}
\newcommand{\popcontenu}{%
  \par\noindent{\popxboldls\fontsize{7.5}{7.5}\selectfont\color{popmuted}CE COURS SIGNÉ CONTIENT}\par\vspace{7pt}
  \noindent\begin{minipage}[t]{0.235\linewidth}\poptile{popblue}{Cours}{complet, illustré, pas à pas}\end{minipage}\hfill
  \begin{minipage}[t]{0.235\linewidth}\poptile{poporange}{Méthodes}{et exercices résolus}\end{minipage}\hfill
  \begin{minipage}[t]{0.235\linewidth}\poptile{poppink}{Exercices}{type examen + corrigés}\end{minipage}\hfill
  \begin{minipage}[t]{0.235\linewidth}\poptile{popgreen}{Fiche bilan}{l'essentiel à retenir}\end{minipage}\par}

% Sommaire
\newtcolorbox{sommaire}{enhanced,colback=white,colframe=popink,boxrule=2.2pt,arc=10pt,
  drop shadow=popink,left=18pt,right=18pt,top=13pt,bottom=13pt,before skip=6pt,after skip=20pt,
  before upper={\poppill{Sommaire}{poppurple}{white}\par\vspace{9pt}\popsemi\fontsize{10.6}{14.5}\selectfont\color{popink}}}

% Signature de fin de cours — poussée tout en bas de la dernière page
\newcommand{\findecours}{%
  \par\vfill\nopagebreak
  \begin{center}\popsquiggle{poporange}\end{center}\vspace{4pt}
  \signatureonbuch
}
% Glyphes absents de Fira Math : redéfinis à partir de glyphes disponibles
\AtBeginDocument{%
  \def\setminus{\mathbin{\backslash}}\let\smallsetminus\setminus
  \def\vdots{\mathord{\vbox{\baselineskip3.2pt\lineskiplimit0pt\kern4pt\hbox{.}\hbox{.}\hbox{.}}}}%
  \def\ddots{\mathinner{\mkern1mu\raise6.5pt\hbox{.}\mkern2mu\raise3.6pt\hbox{.}\mkern2mu\raise0.7pt\hbox{.}\mkern1mu}}%
  \def\triangle{\mathord{\scalebox{0.9}{$\symup{\Delta}$}}}%
  \def\nabla{\mathord{\raisebox{\depth}{\scalebox{1}[-1]{$\symup{\Delta}$}}}}%
  \def\checkmark{\popcheck{popdark}}%
  \def\star{\mathbin{\ast}}\def\bigstar{\mathord{\ast}}%
  \def\diamond{\mathbin{\scalebox{0.6}{$\Diamond$}}}%
  \def\bigcup{\mathop{\mathchoice{\raisebox{-0.3em}{\scalebox{1.7}{$\cup$}}}{\scalebox{1.25}{$\cup$}}{\cup}{\cup}}\displaylimits}%
  \def\bigcap{\mathop{\mathchoice{\raisebox{-0.3em}{\scalebox{1.7}{$\cap$}}}{\scalebox{1.25}{$\cap$}}{\cap}{\cap}}\displaylimits}%
  \def\lesseqqgtr{\mathrel{\lessgtr}}\def\bigoplus{\mathop{\oplus}\displaylimits}%
}
\providecommand{\ssi}{si et seulement si} % abréviation fréquente

%% FIGFIX-BEGIN v4 — correctifs globaux des figures (docs/onbuch-generation/tools/figfix)
\usepackage{adjustbox}\usepackage{etoolbox}
% toute figure TikZ/pgfplots plus large que la ligne est réduite à la largeur disponible
\BeforeBeginEnvironment{tikzpicture}{\begin{adjustbox}{max width=\linewidth}}
\AfterEndEnvironment{tikzpicture}{\end{adjustbox}}
% légendes pgfplots lisibles par-dessus les courbes
\pgfplotsset{every axis/.append style={legend style={fill=white,fill opacity=0.92,text opacity=1,draw=black!15,inner sep=3pt}}}
% étiquettes d'arêtes (syntaxe edge["..."]) avec fond blanc
\tikzset{every edge quotes/.append style={fill=white,inner sep=1.2pt}}
%% FIGFIX-END
\begin{document}

\pagegarde{Réactions d'oxydoréduction spontanées et nombres d'oxydation}{Chimie}{Tle TI}{Module 3 · Oxydoréduction}{09}{7 h}{Cette leçon introduit les réactions d'oxydoréduction à travers l'action des acides sur les métaux et celle des ions métalliques sur les métaux, puis généralise la notion grâce aux nombres d'oxydation. L'élève apprend à écrire et équilibrer les équations-bilans, en solution aqueuse comme par voie sèche, et à interpréter des phénomènes quotidiens comme la formation de la rouille.
\vspace{4pt}
\begin{multicols}{2}\small
\begin{itemize}
\item À la fin de cette leçon, je sais réaliser expérimentalement l'action des acides chlorhydrique et sulfurique dilués sur les métaux (Zn, Fe, Cu, Al, Ag, Au) et identifier les produits formés.
\item À la fin de cette leçon, je sais définir oxydation, réduction, oxydant, réducteur et couple redox du point de vue électronique.
\item À la fin de cette leçon, je sais écrire les demi-équations électroniques et en déduire l'équation-bilan d'une réaction d'oxydoréduction en solution aqueuse.
\item À la fin de cette leçon, je sais réaliser et interpréter l'action du zinc ou du fer sur les ions Cu2+, Zn2+ et Fe2+.
\item À la fin de cette leçon, je sais écrire les équations-bilans des réactions Fe2+/MnO4- en milieu acide, éthanol/Cr2O7 2- et I2/S2O3 2-.
\item À la fin de cette leçon, je sais calculer le nombre d'oxydation d'un élément dans une espèce chimique.
\end{itemize}\end{multicols}}

\begin{sommaire}
\begin{multicols}{2}
\begin{enumerate}
\item Action d'un acide dilué sur un métal
\item Oxydation, réduction et couples redox : le point de vue électronique
\item Action d'un ion métallique sur un métal
\item Autres réactions d'oxydoréduction en solution aqueuse
\item Généralisation : les nombres d'oxydation
\item Réactions d'oxydoréduction par voie sèche et synthèse
\item Travaux pratiques
\item Méthodes et exercices résolus
\item Exercices
\item Corrigés des exercices
\item Fiche bilan
\end{enumerate}
\end{multicols}
\end{sommaire}

\begin{objectifs}
\begin{itemize}
\item À la fin de cette leçon, je sais réaliser expérimentalement l'action des acides chlorhydrique et sulfurique dilués sur les métaux (Zn, Fe, Cu, Al, Ag, Au) et identifier les produits formés.
\item À la fin de cette leçon, je sais définir oxydation, réduction, oxydant, réducteur et couple redox du point de vue électronique.
\item À la fin de cette leçon, je sais écrire les demi-équations électroniques et en déduire l'équation-bilan d'une réaction d'oxydoréduction en solution aqueuse.
\item À la fin de cette leçon, je sais réaliser et interpréter l'action du zinc ou du fer sur les ions Cu2+, Zn2+ et Fe2+.
\item À la fin de cette leçon, je sais écrire les équations-bilans des réactions Fe2+/MnO4- en milieu acide, éthanol/Cr2O7 2- et I2/S2O3 2-.
\item À la fin de cette leçon, je sais calculer le nombre d'oxydation d'un élément dans une espèce chimique.
\item À la fin de cette leçon, je sais équilibrer une équation d'oxydoréduction en utilisant les nombres d'oxydation.
\item À la fin de cette leçon, je sais identifier une réaction d'oxydoréduction par voie sèche et expliquer la formation de la rouille.
\end{itemize}
\end{objectifs}

\begin{prerequis}
\begin{itemize}
\item Structure de l'atome, électrons de valence et ions monoatomiques (classe de Première)
\item Écriture et équilibrage d'une équation-bilan chimique (conservation des éléments et des charges)
\item Ions en solution aqueuse : tests d'identification des ions Cu2+, Fe2+, Fe3+ (précipités d'hydroxydes)
\item Couples acide/base et réactions acido-basiques (notions de début d'année de Terminale)
\item Formules des espèces courantes : HCl, H2SO4, MnO4-, Cr2O7 2-, S2O3 2-, éthanol C2H5OH
\end{itemize}
\end{prerequis}

\newpage

\coursec{Action d'un acide dilué sur un métal}

À Douala, les carrosseries des taxis rouillent rapidement sous l'effet de l'air humide et salin, tandis que les bijoux en or des mamans gardent leur éclat pendant des décennies. Pourquoi certains métaux se dégradent-ils et pas d'autres ? Pourquoi le plongeur de la brasserie utilise-t-il du dichromate de potassium pour détecter l'alcool dans l'haleine d'un chauffeur ? Et comment le fer d'une toiture en tôle peut-il être protégé par de simples plaques de zinc ? Toutes ces questions trouvent leur réponse dans les \cle{réactions d'oxydoréduction}, que cette leçon nous apprend à identifier, écrire et prévoir.

Notre voyage commence par une expérience simple, que tu as peut-être déjà vue au laboratoire : que se passe-t-il lorsqu'on verse un acide dilué sur un métal ? La réponse à cette question apparemment banale contient en germe toute la théorie de l'oxydoréduction.

\courssub{Expérience : action de l'acide chlorhydrique dilué sur le zinc}

\begin{experience}[Attaque du zinc par l'acide chlorhydrique dilué]
\textbf{Dispositif et protocole.} On introduit quelques grammes de \cle{grenaille de zinc} (petits grains de zinc métallique) dans un tube à essais. On verse ensuite de l'acide chlorhydrique dilué (\ce{H3O+} + \ce{Cl-}) en excès. On bouche le tube et on recueille le gaz dégagé par déplacement d'eau ou directement dans un second tube renversé.

\textbf{Observations.}
\begin{itemize}
\item Il se produit une \textbf{effervescence} : des bulles de gaz se forment à la surface du zinc et se dégagent de la solution ;
\item la grenaille de zinc \textbf{disparaît progressivement} : le métal est \emph{attaqué} ;
\item le tube s'échauffe légèrement : la réaction est \textbf{exothermique} ;
\item en approchant une flamme (bûchette en feu) de l'ouverture du tube contenant le gaz recueilli, on entend une \textbf{légère détonation} (un « aboiement » caractéristique) ;
\item si l'on ajoute quelques gouttes de soude (solution d'hydroxyde de sodium) à la solution obtenue en fin de réaction, il apparaît un \textbf{précipité blanc gélatineux}.
\end{itemize}

\textbf{Interprétation.} La détonation à la flamme est le test caractéristique du \cle{dihydrogène} \ce{H2} : le gaz dégagé est donc \ce{H2}. Le précipité blanc avec la soude est de l'hydroxyde de zinc \ce{Zn(OH)2} : il prouve la présence des ions \ce{Zn^2+} en solution. Le zinc métal s'est donc transformé en ions zinc \ce{Zn^2+}, tandis que les ions hydrogène de l'acide se sont transformés en dihydrogène gazeux.
\end{experience}

\begin{popfigure}
\begin{center}
\begin{tikzpicture}[scale=0.95, every node/.style={font=\small}]
% Tube à essais gauche : réaction
\draw[thick, popblue] (-2.2,0) -- (-2.2,3.2);
\draw[thick, popblue] (-0.8,0) -- (-0.8,3.2);
\draw[thick, popblue] (-2.2,0) arc (180:360:0.7);
% Liquide
\fill[popblueL] (-2.15,0.15) -- (-2.15,1.9) -- (-0.85,1.9) -- (-0.85,0.15) arc (180:360:0.65) -- cycle;
% Grenaille de zinc
\foreach \x/\y in {-1.9/0.35,-1.5/0.3,-1.1/0.4,-1.7/0.55,-1.3/0.6}
  \fill[popmuted] (\x,\y) circle (0.09);
% Bulles
\foreach \x/\y in {-1.8/1.0,-1.4/1.3,-1.1/1.1,-1.6/1.6,-1.2/1.7}
  \draw[popblue, fill=white] (\x,\y) circle (0.05);
\node[popdark] at (-1.5,-0.7) {grenaille de zinc};
\node[popdark] at (-1.5,2.35) {\ce{HCl} dilué};
% Flèche vers tube de recueillement
\draw[-{Stealth}, thick, poporange] (-0.6,2.6) -- (0.9,2.6);
% Tube renversé de recueillement
\draw[thick, popgreen] (1.1,3.4) -- (1.1,1.2);
\draw[thick, popgreen] (2.3,3.4) -- (2.3,1.2);
\draw[thick, popgreen] (1.1,3.4) arc (180:0:0.6);
\fill[popgreenL] (1.15,3.35) arc (180:0:0.55) -- (2.25,2.0) -- (1.15,2.0) -- cycle;
\node[popdark] at (1.7,1.55) {gaz \ce{H2}};
% Flamme
\draw[thick, popgold] (3.1,1.0) -- (3.1,1.9);
\fill[poporange] (3.1,2.15) .. controls (2.9,2.5) and (3.3,2.5) .. (3.1,2.85) .. controls (3.25,2.55) and (2.95,2.55) .. (3.1,2.15);
\node[popdark] at (3.1,0.7) {flamme};
% Détonation (remplace starburst par un nœud standard)
\node[draw=popink, fill=poppinkL, rounded corners=3pt, inner sep=3pt, font=\small\bfseries] at (2.4,3.9) {détonation !};
\end{tikzpicture}
\end{center}
\legende{Dispositif expérimental : attaque de la grenaille de zinc par l'acide chlorhydrique dilué, recueillement du dihydrogène et test à la flamme (détonation caractéristique).}
\end{popfigure}

\begin{definition}[Le test du dihydrogène]
Le \cle{dihydrogène} \ce{H2} est un gaz incolore et inodore, très inflammable. Son test d'identification est la \textbf{détonation} : au contact d'une flamme, le mélange dihydrogène--dioxygène de l'air réagit brutalement en produisant une petite explosion sonore appelée « aboiement » :
\[ \ce{2H2(g) + O2(g) -> 2H2O(g)} \]
Tout dégagement gazeux qui produit une détonation à la flamme met donc en évidence du dihydrogène.
\end{definition}

\courssub{Équation-bilan de la réaction et rôle des ions spectateurs}

Que s'est-il réellement passé au niveau microscopique ? Deux transformations ont eu lieu simultanément :

\begin{itemize}
\item chaque atome de zinc a \textbf{perdu deux électrons} pour devenir un ion zinc : $\ce{Zn -> Zn^2+ + 2e-}$ ;
\item chaque ion hydrogène \ce{H+} (présent en solution sous forme \ce{H3O+}) a \textbf{capté un électron} : $\ce{2H+ + 2e- -> H2}$.
\end{itemize}

Les électrons perdus par le zinc sont exactement ceux gagnés par les ions hydrogène : il y a eu \textbf{transfert d'électrons} du métal vers l'acide. En additionnant les deux demi-équations, les électrons s'éliminent et l'on obtient l'équation-bilan :

\[ \ce{Zn(s) + 2H+(aq) -> Zn^2+(aq) + H2(g)} \]

Et les ions chlorure \ce{Cl-} ? Ils étaient présents avant la réaction, ils sont présents après, inchangés : ils n'ont participé à \textbf{aucun} transfert d'électrons. On dit que ce sont des \cle{ions spectateurs}. C'est pourquoi ils n'apparaissent pas dans l'équation-bilan ionique.

\begin{attention}[Erreur fréquente : écrire \ce{HCl} dans l'équation-bilan]
Beaucoup d'élèves écrivent $\ce{Zn + 2HCl -> ZnCl2 + H2}$. Cette écriture n'est pas fausse en apparence, mais elle est \textbf{trompeuse} : en solution aqueuse, l'acide chlorhydrique est \emph{totalement dissocié} en \ce{H+} (solvaté en \ce{H3O+}) et \ce{Cl-}, et le chlorure de zinc est dissocié en \ce{Zn^2+} et \ce{Cl-}. Le seul réactif acide qui agit réellement est l'ion \ce{H+} ; l'ion \ce{Cl-} est spectateur. Au Baccalauréat TI, on exige l'écriture \textbf{ionique} : $\ce{Zn + 2H+ -> Zn^2+ + H2}$.
\end{attention}

\begin{methode}[Identifier les produits d'une attaque acide]
\begin{enumerate}
\item \textbf{Identifier le gaz} : recueillir le gaz dégagé et approcher une flamme. Détonation $\Rightarrow$ \ce{H2}.
\item \textbf{Identifier le cation métallique} en solution par un test de précipitation : avec la soude, \ce{Zn^2+} donne un précipité blanc de \ce{Zn(OH)2}, \ce{Fe^2+} un précipité vert sale de \ce{Fe(OH)2}, \ce{Al^3+} un précipité blanc de \ce{Al(OH)3} (soluble dans un excès de soude).
\item \textbf{Écrire les deux demi-équations} (perte d'électrons par le métal, gain par \ce{H+}), puis les combiner en éliminant les électrons.
\end{enumerate}
\end{methode}

\courssub{Généralisation au fer et à l'aluminium}

La même expérience, réalisée avec de la limaille de fer ou de la poudre d'aluminium, donne des résultats analogues : effervescence de dihydrogène et formation du cation métallique correspondant.

\textbf{Avec le fer} : le fer donne l'ion fer(II) \ce{Fe^2+} (mis en évidence par le précipité vert de \ce{Fe(OH)2} avec la soude) :
\[ \ce{Fe(s) + 2H+(aq) -> Fe^2+(aq) + H2(g)} \]
Demi-équations : $\ce{Fe -> Fe^2+ + 2e-}$ et $\ce{2H+ + 2e- -> H2}$.

\textbf{Avec l'aluminium} : l'aluminium donne l'ion \ce{Al^3+}. Chaque atome d'aluminium perd \emph{trois} électrons : $\ce{Al -> Al^3+ + 3e-}$. Pour équilibrer les échanges électroniques avec $\ce{2H+ + 2e- -> H2}$, il faut 2 atomes d'aluminium (qui perdent $6e^-$) pour 6 ions \ce{H+} (qui gagnent $6e^-$) :
\[ \ce{2Al(s) + 6H+(aq) -> 2Al^3+(aq) + 3H2(g)} \]

\begin{attention}[Bien équilibrer les électrons échangés]
Le nombre d'électrons \textbf{perdus} par le métal doit être exactement égal au nombre d'électrons \textbf{gagnés} par les ions \ce{H+}. Avec l'aluminium, oublier ce point conduit à l'équation fausse $\ce{Al + 2H+ -> Al^3+ + H2}$, qui ne conserve ni la charge ni la matière. Vérifie toujours la conservation de la charge électrique de part et d'autre de la flèche.
\end{attention}

\courssub{Action de l'acide sulfurique dilué : le rôle central de l'ion \ce{H+}}

Reprenons l'expérience en remplaçant l'acide chlorhydrique par de l'\textbf{acide sulfurique dilué} (\ce{2H+} + \ce{SO4^2-}). Les observations sont \emph{identiques} : effervescence de \ce{H2}, disparition du zinc, précipité blanc de \ce{Zn(OH)2} avec la soude. L'équation-bilan ionique est la même :
\[ \ce{Zn(s) + 2H+(aq) -> Zn^2+(aq) + H2(g)} \]
Cette fois, ce sont les ions sulfate \ce{SO4^2-} qui sont spectateurs. La conclusion est capitale : \textbf{quel que soit l'acide dilué non oxydant utilisé, le réactif commun est l'ion \ce{H+}}. C'est lui qui capte les électrons du métal. L'anion de l'acide (\ce{Cl-}, \ce{SO4^2-}) ne joue aucun rôle dans le transfert électronique.

\begin{aretenir}[À retenir]
Les acides chlorhydrique et sulfurique \textbf{dilués} attaquent certains métaux en dégageant du \cle{dihydrogène} et en formant le \cle{cation métallique} correspondant. L'équation-bilan ionique fait toujours intervenir le même réactif acide : l'ion \ce{H+}. Le métal \textbf{perd des électrons}, l'ion \ce{H+} les \textbf{capte} : c'est une réaction d'\cle{oxydoréduction}.
\end{aretenir}

\courssub{Le cas du cuivre, de l'argent et de l'or : tous les métaux ne réagissent pas}

Renouvelons l'expérience avec un fil de cuivre, une pièce d'argent ou un bijou en or : \textbf{aucune effervescence, aucun dégagement gazeux, le métal reste intact}, même après plusieurs heures. Les acides chlorhydrique et sulfurique dilués sont \emph{sans action} sur le cuivre, l'argent et l'or.

Comment l'expliquer ? La réaction $\ce{M + n H+ -> M^{n+} + \frac{n}{2} H2}$ n'est possible que si le métal \textbf{accepte de céder ses électrons} aux ions \ce{H+}, c'est-à-dire s'il est \emph{plus réducteur} que le dihydrogène. Le zinc, le fer et l'aluminium sont des métaux \textbf{plus réducteurs que \ce{H2}} : ils donnent spontanément leurs électrons à \ce{H+}. Le cuivre, l'argent et l'or sont \textbf{moins réducteurs que \ce{H2}} : ils « accrochent » trop leurs électrons pour les céder aux ions hydrogène. Voilà pourquoi les bijoux en or traversent les décennies sans s'altérer, tandis que la tôle de fer rouille !

\begin{center}
\begin{tabularx}{\linewidth}{|l|c|X|X|}
\hline
\rowcolor{popblueL} \textbf{Métal} & \textbf{Attaqué par \ce{HCl} dilué ?} & \textbf{Produits} & \textbf{Équation-bilan ionique} \\
\hline
Zinc \ce{Zn} & Oui (vive) & \ce{Zn^2+} + \ce{H2} & $\ce{Zn + 2H+ -> Zn^2+ + H2}$ \\
\hline
Fer \ce{Fe} & Oui & \ce{Fe^2+} + \ce{H2} & $\ce{Fe + 2H+ -> Fe^2+ + H2}$ \\
\hline
Aluminium \ce{Al} & Oui (lente au départ) & \ce{Al^3+} + \ce{H2} & $\ce{2Al + 6H+ -> 2Al^3+ + 3H2}$ \\
\hline
Cuivre \ce{Cu} & Non & --- & pas de réaction \\
\hline
Argent \ce{Ag} & Non & --- & pas de réaction \\
\hline
Or \ce{Au} & Non & --- & pas de réaction \\
\hline
\end{tabularx}
\end{center}

\begin{savaistu}[Pourquoi l'aluminium semble-t-il inerte à l'air libre ?]
L'aluminium est pourtant un métal \emph{très} réducteur, bien plus que le fer ! Alors pourquoi les casseroles et les tôles d'aluminium, si répandues dans nos cuisines et nos quartiers, ne se dégradent-elles pas ? Parce qu'au contact de l'air, l'aluminium se recouvre instantanément d'une \textbf{couche très fine (quelques nanomètres) d'oxyde d'aluminium \ce{Al2O3}}, appelée \cle{alumine}. Cette couche, dure, adhérente et \emph{imperméable}, isole le métal de l'extérieur : on dit que l'aluminium est \textbf{passivé}. C'est d'ailleurs pourquoi l'attaque par l'acide dilué est lente au départ : l'acide doit d'abord dissoudre la couche d'alumine avant d'atteindre le métal. Le fer, lui, forme une rouille poreuse et friable qui laisse passer l'air et l'eau : la corrosion se poursuit en profondeur.
\end{savaistu}

\courssub{Interprétation électronique : naissance de l'oxydoréduction}

Résumons le point de vue électronique sur l'attaque du zinc :
\begin{itemize}
\item le zinc \textbf{perd} des électrons : $\ce{Zn -> Zn^2+ + 2e-}$ ; on dit que le zinc est \textbf{oxydé} et qu'il joue le rôle de \cle{réducteur} ;
\item les ions \ce{H+} \textbf{gagnent} des électrons : $\ce{2H+ + 2e- -> H2}$ ; on dit que \ce{H+} est \textbf{réduit} et qu'il joue le rôle d'\cle{oxydant}.
\end{itemize}
Une réaction au cours de laquelle il y a \textbf{transfert d'électrons} entre un réducteur (qui les cède) et un oxydant (qui les capte) est une \cle{réaction d'oxydoréduction}. Nous définirons ces termes rigoureusement dans la section suivante ; retiens dès maintenant l'idée intuitive : \emph{le métal donne, l'acide prend}.

\begin{exoresolu}[Volume de dihydrogène dégagé par l'attaque du zinc]
On attaque \SI{6,54}{g} de grenaille de zinc par un excès d'acide chlorhydrique dilué. On donne $M(\ce{Zn}) = \SI{65,4}{g.mol^{-1}}$ et le volume molaire dans les CNTP $V_m = \SI{22,4}{L.mol^{-1}}$. Déterminer la quantité et le volume de dihydrogène dégagés dans les CNTP.
\tcblower
\textbf{Solution détaillée.}

Étape 1 : quantité de matière de zinc consommée (l'acide est en excès, le zinc est le réactif limitant) :
\[ n(\ce{Zn}) = \frac{m}{M} = \frac{6,54}{65,4} = \SI{0,100}{mol}. \]

Étape 2 : exploitation de l'équation-bilan $\ce{Zn + 2H+ -> Zn^2+ + H2}$. D'après les coefficients stœchiométriques, $1$ mol de zinc produit $1$ mol de dihydrogène :
\[ n(\ce{H2}) = n(\ce{Zn}) = \SI{0,100}{mol}. \]

Étape 3 : volume de gaz dégagé dans les CNTP :
\[ V(\ce{H2}) = n(\ce{H2}) \times V_m = 0,100 \times 22,4 = \SI{2,24}{L}. \]

\textbf{Conclusion} : l'attaque de \SI{6,54}{g} de zinc dégage \SI{2,24}{L} de dihydrogène dans les CNTP. Remarque le réflexe à avoir : passer par les \emph{quantités de matière} et les coefficients de l'équation-bilan, jamais directement des grammes aux litres.
\end{exoresolu}

\begin{exercice}[Attaque du fer par l'acide sulfurique dilué]
On fait réagir \SI{2,8}{g} de limaille de fer avec un excès d'acide sulfurique dilué. On donne $M(\ce{Fe}) = \SI{56}{g.mol^{-1}}$ et $V_m = \SI{22,4}{L.mol^{-1}}$ dans les CNTP.
\begin{enumerate}
\item Écrire les deux demi-équations électroniques puis l'équation-bilan ionique de la réaction.
\item Calculer le volume de dihydrogène dégagé dans les CNTP.
\item Quel test permet d'identifier le gaz recueilli ? Quel réactif permet de mettre en évidence les ions formés en solution, et qu'observe-t-on ?
\end{enumerate}
\end{exercice}

\begin{corrige}[Exercice 1]
\textbf{1.} Oxydation du fer : $\ce{Fe -> Fe^2+ + 2e-}$ ; réduction des ions hydrogène : $\ce{2H+ + 2e- -> H2}$. Équation-bilan : $\ce{Fe + 2H+ -> Fe^2+ + H2}$.

\textbf{2.} $n(\ce{Fe}) = \dfrac{2,8}{56} = \SI{0,050}{mol}$, donc $n(\ce{H2}) = \SI{0,050}{mol}$ et $V(\ce{H2}) = 0,050 \times 22,4 = \SI{1,12}{L}$.

\textbf{3.} Le gaz produit une détonation (aboiement) au contact d'une flamme : c'est \ce{H2}. L'ajout de soude à la solution donne un précipité vert sale d'hydroxyde de fer(II) \ce{Fe(OH)2}, caractéristique des ions \ce{Fe^2+}.
\end{corrige}

En résumé, cette première étude expérimentale nous a révélé trois idées fondamentales : certains métaux cèdent spontanément des électrons aux ions \ce{H+} des acides dilués, d'autres non ; ce transfert d'électrons est le cœur d'une réaction d'oxydoréduction ; et il existe une \emph{hiérarchie} entre les métaux, des plus réducteurs (zinc, fer, aluminium) aux moins réducteurs (cuivre, argent, or). C'est cette hiérarchie que nous allons exploiter dans la suite de la leçon, en faisant réagir cette fois un métal avec les ions d'un \emph{autre} métal.

\coursec{Oxydation, réduction et couples redox : le point de vue électronique}

Dans la section précédente, tu as observé expérimentalement qu'un clou de zinc plongé dans l'acide chlorhydrique dilué se ronge lentement pendant qu'un dégagement gazeux de dihydrogène se produit. L'équation-bilan de cette transformation s'écrit :

\[ \ce{Zn(s) + 2H+(aq) -> Zn^{2+}(aq) + H2(g)} \]

Cette équation est correctement équilibrée, mais elle cache l'essentiel : \emph{que s'est-il réellement passé au niveau microscopique ?} Pourquoi le zinc disparaît-il ? D'où vient le dihydrogène ? La réponse tient en un mot : les \cle{électrons}. Dans cette section, nous allons ouvrir le « capot » de cette réaction et découvrir le mécanisme fondamental de toutes les réactions d'oxydoréduction : le \cle{transfert d'électrons}. C'est la clé qui te permettra de comprendre et d'écrire toutes les réactions de ce module, du dosage au permanganate jusqu'à la formation de la rouille.

\courssub{Analyse microscopique de la réaction entre le zinc et les ions \ce{H+}}

Regardons de très près ce qui se passe à la surface du zinc.

L'atome de zinc \ce{Zn} possède des électrons périphériques qu'il peut céder assez facilement. Lorsqu'il rencontre des ions hydrogène \ce{H+} en solution, chaque atome de zinc \textbf{cède deux électrons} et se transforme en ion zinc \ce{Zn^{2+}}, qui passe en solution :

\[ \ce{Zn -> Zn^{2+} + 2e-} \]

Ces deux électrons ne restent pas libres dans l'eau (un électron libre ne peut pas exister durablement en solution aqueuse) : ils sont immédiatement \textbf{captés par deux ions \ce{H+}}, qui se transforment en une molécule de dihydrogène gazeux :

\[ \ce{2H+ + 2e- -> H2} \]

\begin{popfigure}
\begin{center}
\begin{tikzpicture}[scale=1.0, every node/.style={font=\small}]
% Atome de zinc
\node[draw=popblue, thick, fill=popblue!35, circle, minimum size=1.5cm] (Zn) at (0,0) {\textbf{Zn}};
\node[below=0.85cm of Zn, text=popblue] {atome de zinc (donneur)};
% Ions H+
\node[draw=poporange, thick, fill=poporange!35, circle, minimum size=0.85cm] (H1) at (5.2,0.9) {\textbf{\ce{H+}}};
\node[draw=poporange, thick, fill=poporange!35, circle, minimum size=0.85cm] (H2) at (5.2,-0.9) {\textbf{\ce{H+}}};
\node[below=0.6cm of H2, text=poporange] {ions hydrogène (accepteurs)};
% Flèches d'électrons
\draw[-{Stealth[length=3mm]}, very thick, poppink] (Zn.30) .. controls (2.4,1.5) and (3.6,1.6) .. (H1.150)
  node[midway, above, text=poppink] {$e^-$};
\draw[-{Stealth[length=3mm]}, very thick, poppink] (Zn.-30) .. controls (2.4,-1.5) and (3.6,-1.6) .. (H2.210)
  node[midway, below, text=poppink] {$e^-$};
% Flèche de réaction
\draw[-{Stealth[length=3mm]}, thick, popdark] (6.1,0) -- (7.1,0);
% Produits
\node[draw=popblue, thick, fill=popblue!15, circle, minimum size=1.1cm] (Zn2) at (8.1,0.7) {\textbf{\ce{Zn^{2+}}}};
\node[draw=poporange, thick, fill=poporange!15, circle, minimum size=0.65cm] (H3) at (7.85,-0.8) {\textbf{H}};
\node[draw=poporange, thick, fill=poporange!15, circle, minimum size=0.65cm] (H4) at (8.5,-0.8) {\textbf{H}};
\node[below=0.5cm of H4, text=popdark] {\ce{Zn^{2+}} en solution et \ce{H2} gazeux};
\end{tikzpicture}
\end{center}
\legende{Transfert de deux électrons d'un atome de zinc vers deux ions hydrogène : le zinc est oxydé en ion \ce{Zn^{2+}}, les ions \ce{H+} sont réduits en dihydrogène \ce{H2}.}
\end{popfigure}

Le bilan net est donc un \textbf{transfert direct d'électrons} du zinc vers les ions hydrogène. Aucun électron n'apparaît dans l'équation-bilan globale, car les deux électrons cédés par le zinc sont exactement les deux électrons captés par les ions \ce{H+} : tout se passe comme si les électrons « transitaient » d'un partenaire à l'autre.

\begin{aretenir}[L'essentiel du mécanisme]
Une réaction d'oxydoréduction est une réaction au cours de laquelle il y a \cle{transfert d'électrons} entre les réactifs. Les électrons cédés par l'un sont intégralement captés par l'autre : il n'y a ni création ni disparition d'électrons.
\end{aretenir}

\courssub{Oxydation et réduction : les définitions électroniques}

L'expérience précédente met en évidence deux phénomènes simultanés et indissociables :

\begin{itemize}
\item le zinc \textbf{perd} deux électrons : on dit qu'il subit une \cle{oxydation} ;
\item les ions \ce{H+} \textbf{gagnent} chacun un électron : on dit qu'ils subissent une \cle{réduction}.
\end{itemize}

\begin{definition}[Oxydation et réduction (point de vue électronique)]
\begin{itemize}
\item L'\cle{oxydation} est une \textbf{perte d'électron(s)} par une espèce chimique.
\item La \cle{réduction} est un \textbf{gain d'électron(s)} par une espèce chimique.
\end{itemize}
Oxydation et réduction se produisent toujours \textbf{simultanément} : on ne peut pas gagner des électrons si personne ne les cède. L'ensemble constitue une réaction d'\cle{oxydoréduction} (ou réaction « redox »).
\end{definition}

\begin{attention}[Un piège de vocabulaire classique]
Ne te fie pas au sens courant du mot « oxydation » : une oxydation n'implique \textbf{pas nécessairement} l'élément oxygène ! Le zinc est oxydé par les ions \ce{H+} en l'absence totale de dioxygène. Retiens le moyen mnémotechnique : \textbf{oxydation = perte d'électrons} (les deux mots commencent par une voyelle), \textbf{réduction = gain d'électrons}.
\end{attention}

\courssub{Oxydant et réducteur}

Chaque partenaire du transfert joue un rôle bien précis :

\begin{definition}[Oxydant et réducteur]
\begin{itemize}
\item Un \cle{oxydant} est une espèce chimique capable de \textbf{capter (accepter) un ou plusieurs électrons}. Au cours de la réaction, l'oxydant est \textbf{réduit}.
\item Un \cle{réducteur} est une espèce chimique capable de \textbf{céder (donner) un ou plusieurs électrons}. Au cours de la réaction, le réducteur est \textbf{oxydé}.
\end{itemize}
\end{definition}

Dans notre exemple : l'ion \ce{H+} capte des électrons, c'est l'\textbf{oxydant} (il est réduit en \ce{H2}) ; le zinc \ce{Zn} cède des électrons, c'est le \textbf{réducteur} (il est oxydé en \ce{Zn^{2+}}).

\begin{attention}[Le paradoxe à bien maîtriser]
C'est la source d'erreur numéro un au baccalauréat : \textbf{l'oxydant est réduit} et \textbf{le réducteur est oxydé}. Le nom décrit l'action que l'espèce \emph{provoque chez son partenaire}, pas ce qu'elle subit. Le réducteur « réduit » l'autre... et donc s'oxyde lui-même. Prends le temps de vérifier cette logique sur chaque exercice.
\end{attention}

\courssub{Les demi-équations électroniques}

Chacun des deux phénomènes (oxydation et réduction) se traduit par une écriture formelle appelée \cle{demi-équation électronique}, dans laquelle les électrons apparaissent explicitement.

Pour le couple du zinc, l'oxydation du réducteur s'écrit :

\[ \ce{Zn = Zn^{2+} + 2e-} \]

Pour le couple de l'hydrogène, la réduction de l'oxydant s'écrit :

\[ \ce{2H+ + 2e- = H2} \]

On utilise le signe « = » (et non une flèche) car une demi-équation est une écriture formelle qui peut être lue dans les deux sens : de gauche à droite c'est une réduction, de droite à gauche c'est une oxydation.

\begin{propriete}[Double conservation dans une demi-équation]
Toute demi-équation électronique doit vérifier \textbf{simultanément} :
\begin{enumerate}
\item la \textbf{conservation des éléments} : même nombre d'atomes de chaque élément des deux côtés ;
\item la \textbf{conservation de la charge électrique} : la somme algébrique des charges (en comptant chaque électron comme une charge $-1$) est identique des deux côtés.
\end{enumerate}
\end{propriete}

Vérifions sur $\ce{Zn = Zn^{2+} + 2e-}$ : un atome de zinc de chaque côté (éléments conservés) ; charge à droite : $(+2) + 2\times(-1) = 0$, charge à gauche : $0$ (charges conservées). La demi-équation est correcte.

\begin{attention}[Erreur fréquente]
Écrire \ce{Zn -> Zn^{2+}} sans les électrons est \textbf{faux} : la charge n'est pas conservée ($0 \neq +2$). Les électrons sont obligatoires dans une demi-équation. En revanche, ils devront disparaître de l'équation-bilan finale, nous y reviendrons.
\end{attention}

\courssub{La notion de couple oxydant/réducteur}

Une même espèce « famille » peut exister sous deux formes : une forme oxydée (qui a perdu des électrons, ou peut en capter) et une forme réduite (qui a gagné des électrons, ou peut en céder). Ces deux formes sont liées par une demi-équation et constituent un \cle{couple redox}.

\begin{definition}[Couple oxydant/réducteur]
Un \cle{couple redox}, noté \textbf{Ox/Red}, associe un oxydant Ox et un réducteur Red reliés par la demi-équation électronique :
\[ \ce{Ox + $n$e- = Red} \]
où $n$ est le nombre d'électrons échangés. L'oxydant et le réducteur d'un même couple sont dits \emph{conjugués}.
\end{definition}

Voici les couples que tu as déjà rencontrés ou que tu rencontreras très souvent en Terminale TI :

\begin{center}
\begin{tabularx}{\linewidth}{|l|X|l|}
\hline
\rowcolor{popblueL} \textbf{Couple Ox/Red} & \textbf{Demi-équation électronique} & \textbf{Remarque} \\
\hline
\ce{Zn^{2+}}/\ce{Zn} & \ce{Zn^{2+} + 2e- = Zn} & métal/ion métallique \\
\hline
\ce{H+}/\ce{H2} & \ce{2H+ + 2e- = H2} & action des acides sur les métaux \\
\hline
\ce{Cu^{2+}}/\ce{Cu} & \ce{Cu^{2+} + 2e- = Cu} & dépôt de cuivre rouge \\
\hline
\ce{Fe^{3+}}/\ce{Fe^{2+}} & \ce{Fe^{3+} + e- = Fe^{2+}} & un seul électron échangé \\
\hline
\end{tabularx}
\end{center}

\begin{exemplebox}[Le couple \ce{Fe^{3+}}/\ce{Fe^{2+}}]
L'ion fer(III) \ce{Fe^{3+}} (couleur rouille en solution) peut capter \textbf{un seul} électron pour devenir ion fer(II) \ce{Fe^{2+}} (couleur vert pâle) :
\[ \ce{Fe^{3+} + e- = Fe^{2+}} \]
Vérification des charges : à gauche $(+3) + (-1) = +2$ ; à droite $+2$. Ce couple est extrêmement utilisé : il intervient dans de nombreux dosages (par exemple le dosage des ions \ce{Fe^{2+}} par le permanganate, que tu étudieras dans la section 4) et dans les phénomènes de corrosion du fer.
\end{exemplebox}

\begin{savaistu}[D'où vient le mot « réduction » ?]
Le terme « réduction » est né dans les hauts fourneaux de la métallurgie, bien avant qu'on découvre l'électron ! « Réduire » un minerai, c'était le \emph{réduire} (diminuer) en métal : par exemple, transformer l'oxyde de fer \ce{Fe2O3} du minerai en fer métallique \ce{Fe}. On sait aujourd'hui que cette transformation est bien un gain d'électrons ($\ce{Fe^{3+}}$ dans l'oxyde devient \ce{Fe} métal), ce qui a donné son sens moderne au mot. À l'inverse, le mot « oxydation » vient de l'oxygène, premier oxydant connu des chimistes du XVIII\textsuperscript{e} siècle avec Lavoisier.
\end{savaistu}

\courssub{Une réaction d'oxydoréduction met en jeu deux couples}

Reprenons la réaction du zinc avec les ions \ce{H+}. Elle fait intervenir \textbf{deux couples} : \ce{Zn^{2+}}/\ce{Zn} et \ce{H+}/\ce{H2}. Le réducteur de l'un (\ce{Zn}) a cédé ses électrons à l'oxydant de l'autre (\ce{H+}).

\begin{propriete}[Règle fondamentale]
Toute réaction d'oxydoréduction met en jeu \textbf{deux couples redox} Ox$_1$/Red$_1$ et Ox$_2$/Red$_2$ : c'est l'\textbf{oxydant de l'un des couples} qui réagit avec le \textbf{réducteur de l'autre couple}. L'équation-bilan s'écrit :
\[ \ce{Ox1 + Red2 -> Red1 + Ox2} \]
\end{propriete}

\begin{popfigure}
\begin{center}
\begin{tikzpicture}[scale=1.0, every node/.style={font=\small}]
% Couple 1
\node[draw=popblue, thick, rounded corners, fill=popblue!15, minimum width=3.4cm, minimum height=0.9cm] (c1) at (0,2) {\ce{Ox1}/\ce{Red1} \; ($\ce{Ox1} + n_1 e^- = \ce{Red1}$)};
% Couple 2
\node[draw=poporange, thick, rounded corners, fill=poporange!15, minimum width=3.4cm, minimum height=0.9cm] (c2) at (0,-2) {\ce{Ox2}/\ce{Red2} \; ($\ce{Ox2} + n_2 e^- = \ce{Red2}$)};
% Flèches de réaction
\draw[-{Stealth[length=3mm]}, very thick, popgreen] (-1.2,1.5) -- (-1.2,-1.5)
  node[midway, left, text=popgreen, align=center] {\ce{Ox1} capte\\ des électrons\\ (il est réduit)};
\draw[-{Stealth[length=3mm]}, very thick, poppink] (1.2,-1.5) -- (1.2,1.5)
  node[midway, right, text=poppink, align=center] {\ce{Red2} cède\\ des électrons\\ (il est oxydé)};
% Équation résultante
\node[text=popdark, font=\small\bfseries] at (0,-3.6) {Réaction : \ce{Ox1} + \ce{Red2} -> \ce{Red1} + \ce{Ox2}};
\end{tikzpicture}
\end{center}
\legende{Schéma général d'une réaction d'oxydoréduction : l'oxydant du couple 1 capte les électrons cédés par le réducteur du couple 2.}
\end{popfigure}

Dans le cas du zinc et de l'acide : \ce{Ox1} = \ce{H+}, \ce{Red1} = \ce{H2}, \ce{Ox2} = \ce{Zn^{2+}}, \ce{Red2} = \ce{Zn}. La réaction est donc $\ce{H+} + \ce{Zn} \to \ce{H2} + \ce{Zn^{2+}}$ (à équilibrer).

\courssub{Méthode : construire l'équation-bilan à partir des demi-équations}

Voici la méthode systématique qui te servira dans tout le module, y compris pour les réactions complexes en milieu acide.

\begin{methode}[Écrire l'équation-bilan d'une réaction d'oxydoréduction]
\begin{enumerate}
\item \textbf{Identifier les deux couples} mis en jeu et repérer quel réactif est l'oxydant, quel réactif est le réducteur.
\item \textbf{Écrire les deux demi-équations} électroniques, en vérifiant la conservation des éléments et des charges.
\item \textbf{Multiplier} chaque demi-équation par un coefficient tel que le nombre d'électrons \textbf{cédés} soit égal au nombre d'électrons \textbf{captés} (plus petit multiple commun).
\item \textbf{Additionner membre à membre} les deux demi-équations : les électrons s'éliminent.
\item \textbf{Vérifier} l'équation-bilan obtenue : conservation des éléments et des charges, et absence totale d'électrons.
\end{enumerate}
\end{methode}

\begin{attention}[Les électrons n'apparaissent JAMAIS dans l'équation-bilan]
C'est une règle absolue : dans l'équation-bilan finale d'une réaction d'oxydoréduction, il ne doit rester \textbf{aucun électron}. Si tu obtiens des « $e^-$ » dans ton équation finale, c'est que tu n'as pas correctement multiplié les demi-équations. Les électrons ne font que transiter : le nombre d'électrons cédés par le réducteur est \emph{exactement égal} au nombre d'électrons captés par l'oxydant.
\end{attention}

\begin{exoresolu}[Action de l'acide chlorhydrique sur l'aluminium]
L'aluminium \ce{Al} (couple \ce{Al^{3+}}/\ce{Al}) est attaqué par l'acide chlorhydrique dilué (couple \ce{H+}/\ce{H2}). Écris les demi-équations électroniques, identifie l'oxydant et le réducteur, puis établis l'équation-bilan de la réaction.
\tcblower
\textbf{Étape 1 — Couples et rôles.} L'aluminium métal \ce{Al} est un réducteur (il va s'oxyder en \ce{Al^{3+}}) ; l'ion \ce{H+} est l'oxydant (il va se réduire en \ce{H2}).

\textbf{Étape 2 — Demi-équations.}
\begin{align*}
\text{Oxydation : } & \ce{Al = Al^{3+} + 3e-} \\
\text{Réduction : } & \ce{2H+ + 2e- = H2}
\end{align*}

\textbf{Étape 3 — Multiplication.} L'aluminium cède 3 électrons, chaque molécule de \ce{H2} formée en capte 2. Le plus petit multiple commun est 6 : on multiplie la première demi-équation par 2 et la seconde par 3.
\begin{align*}
\ce{Al = Al^{3+} + 3e-} \quad &(\times 2) \\
\ce{2H+ + 2e- = H2} \quad &(\times 3)
\end{align*}

\textbf{Étape 4 — Addition.} Les $6\,e^-$ cédés compensent exactement les $6\,e^-$ captés :
\[ \ce{2Al + 6H+ -> 2Al^{3+} + 3H2} \]

\textbf{Étape 5 — Vérification.} Éléments : 2 Al et 6 H de chaque côté. Charges : à gauche $6 \times (+1) = +6$ ; à droite $2 \times (+3) = +6$. Aucun électron n'apparaît : l'équation est correcte. Cette réaction, lente au début à cause de la couche d'alumine protectrice, dégage du dihydrogène et chauffe sensiblement (transformation exothermique).
\end{exoresolu}

\begin{exercice}[À toi de jouer]
Le fer métal \ce{Fe} réagit avec les ions hydrogène \ce{H+} pour donner des ions fer(II) \ce{Fe^{2+}} et du dihydrogène (couples \ce{Fe^{2+}}/\ce{Fe} et \ce{H+}/\ce{H2}). Écris les deux demi-équations électroniques, précise qui est oxydé et qui est réduit, puis établis l'équation-bilan.
\end{exercice}

\begin{corrige}[Exercice : fer et ions \ce{H+}]
Oxydation du réducteur (le fer) : $\ce{Fe = Fe^{2+} + 2e-}$. Réduction de l'oxydant (l'ion \ce{H+}) : $\ce{2H+ + 2e- = H2}$. Le nombre d'électrons échangés est déjà le même (2) : il suffit d'additionner membre à membre :
\[ \ce{Fe + 2H+ -> Fe^{2+} + H2} \]
Le fer est oxydé (il perd 2 électrons), les ions \ce{H+} sont réduits (ils gagnent des électrons). Vérification des charges : $+2$ de chaque côté.
\end{corrige}

\begin{aretenir}[Ce qu'il faut retenir de cette section]
\begin{itemize}
\item Une réaction d'\cle{oxydoréduction} est un \textbf{transfert d'électrons} entre réactifs.
\item \textbf{Oxydation} = perte d'électrons ; \textbf{réduction} = gain d'électrons.
\item L'\textbf{oxydant} capte des électrons (il est réduit) ; le \textbf{réducteur} cède des électrons (il est oxydé).
\item Un \textbf{couple redox} Ox/Red s'écrit $\ce{Ox + $n$e- = Red}$ ; toute réaction redox met en jeu deux couples : l'oxydant de l'un réagit avec le réducteur de l'autre.
\item L'équation-bilan s'obtient en combinant les demi-équations de façon à \textbf{éliminer les électrons} ; éléments et charges doivent être conservés à chaque étape.
\end{itemize}
\end{aretenir}

Tu maîtrises maintenant le langage électronique de l'oxydoréduction. Dans la section suivante, nous appliquerons ces outils à une nouvelle famille d'expériences : l'action d'un ion métallique sur un métal, comme celle du zinc plongé dans une solution de sulfate de cuivre — une expérience spectaculaire qui te permettra de classer les métaux selon leur pouvoir réducteur.

\coursec{Action d'un ion métallique sur un métal}

Dans la section précédente, nous avons appris à reconnaître une oxydation et une réduction grâce aux transferts d'électrons, et nous avons défini les couples oxydant/réducteur. Nous avons aussi vu qu'un acide dilué attaque certains métaux : c'est l'ion \ce{H3O+} qui capte les électrons du métal. Mais l'ion \ce{H3O+} n'est pas le seul oxydant présent en solution aqueuse. Que se passe-t-il si l'on plonge un métal dans une solution contenant les ions \emph{d'un autre métal} ? Un bijoutier qui trempe un clou en fer dans une solution de cuivre, ou un plongeur qui fixe un bloc de zinc sur la coque d'une pirogue à Kribi, utilisent sans le savoir la même chimie. C'est ce que nous allons découvrir maintenant.

\courssub{Expérience : le zinc dans une solution de sulfate de cuivre(II)}

\begin{experience}[Le zinc « dévore » les ions cuivre]
\textbf{Dispositif et protocole.} On plonge une lame de zinc bien décapée (gris brillant) dans un bécher contenant une solution de sulfate de cuivre(II) \ce{CuSO4}, de couleur bleue due aux ions \ce{Cu^2+}.

\textbf{Observations.} Au bout de quelques minutes :
\begin{itemize}
\item la lame de zinc se recouvre d'un \cle{dépôt rouge-orangé} ;
\item la solution bleue se \cle{décolore} progressivement ;
\item la lame de zinc s'use : elle perd de la masse.
\end{itemize}

\textbf{Tests d'identification.} Le dépôt rouge est du \cle{cuivre métallique} \ce{Cu} (couleur caractéristique, conducteur). La solution décolorée, additionnée de soude, donne un précipité blanc gélatineux \textbf{soluble en excès de soude} : c'est le test caractéristique des ions \ce{Zn^2+}.

\textbf{Interprétation.} Les ions \ce{Cu^2+} ont disparu de la solution (décoloration) pour se transformer en cuivre métal déposé ; le zinc s'est transformé en ions \ce{Zn^2+} passés en solution.
\end{experience}

\begin{popfigure}
\begin{tikzpicture}[scale=0.9]
% Bécher avant
\draw[thick] (0,0) -- (0,3) (2.6,0) -- (2.6,3) (0,0) -- (2.6,0);
\fill[popblue!35] (0.05,0.05) rectangle (2.55,2.2);
\draw[popblue!60] (0.05,2.2) -- (2.55,2.2);
% lame de zinc
\fill[popmuted!60] (1.05,0) rectangle (1.55,3.3);
\node[rotate=90] at (1.3,3.7) {\small lame de zinc};
\node at (1.3,-0.6) {\small solution bleue de \ce{CuSO4}};
\node at (1.3,-1.2) {\small\textbf{AVANT}};
% flèche temps
\draw[-{Stealth},very thick,poporange] (3.2,1.5) -- (4.4,1.5);
\node at (3.8,1.9) {\small quelques};
\node at (3.8,1.1) {\small minutes};
% Bécher après
\begin{scope}[xshift=5cm]
\draw[thick] (0,0) -- (0,3) (2.6,0) -- (2.6,3) (0,0) -- (2.6,0);
\fill[popblue!8] (0.05,0.05) rectangle (2.55,2.2);
\draw[popblue!40] (0.05,2.2) -- (2.55,2.2);
% lame amincie + dépôt cuivre
\fill[popmuted!50] (1.12,0.0) rectangle (1.48,3.3);
\fill[poporange!85] (1.02,0.0) rectangle (1.12,1.9);
\fill[poporange!85] (1.48,0.0) rectangle (1.58,1.9);
\draw[-{Stealth},poporange] (2.9,1.2) -- (1.62,1.2);
\node[poporange] at (3.6,1.2) {\small dépôt de};
\node[poporange] at (3.6,0.8) {\small cuivre};
\node at (1.3,-0.6) {\small solution décolorée : ions \ce{Zn^2+}};
\node at (1.3,-1.2) {\small\textbf{APRÈS}};
\end{scope}
\end{tikzpicture}
\legende{Lame de zinc plongée dans une solution de sulfate de cuivre(II) : dépôt rouge de cuivre et décoloration de la solution.}
\end{popfigure}

\courssub{Interprétation électronique : les couples \ce{Zn^2+}/\ce{Zn} et \ce{Cu^2+}/\ce{Cu}}

L'équation-bilan de la réaction observée s'écrit :

\[ \ce{Zn(s) + Cu^2+(aq) -> Zn^2+(aq) + Cu(s)} \]

Décortiquons ce transfert d'électrons à l'aide des deux couples mis en jeu :

\begin{itemize}
\item le zinc \cle{perd} deux électrons : il est \textbf{oxydé}, c'est le \textbf{réducteur} :
\[ \ce{Zn -> Zn^2+ + 2e-} \qquad \text{couple } \ce{Zn^2+}/\ce{Zn} \]
\item l'ion cuivre(II) \cle{gagne} deux électrons : il est \textbf{réduit}, c'est l'\textbf{oxydant} :
\[ \ce{Cu^2+ + 2e- -> Cu} \qquad \text{couple } \ce{Cu^2+}/\ce{Cu} \]
\end{itemize}

Les électrons ne circulent pas librement dans la solution : ils passent \emph{directement} du zinc aux ions \ce{Cu^2+} au contact de la lame. La réaction se produit dès la mise en contact, à température ambiante, sans chauffage ni apport d'énergie : c'est une réaction \cle{spontanée}.

\begin{definition}[Réaction d'oxydoréduction spontanée]
Une réaction d'\cle{oxydoréduction} est dite \cle{spontanée} si elle se produit dès la mise en contact des réactifs, sans apport d'énergie extérieur (ni chauffage, ni courant électrique). Toutes les réactions étudiées dans cette section sont spontanées.
\end{definition}

\begin{attention}[Du cuivre, pas de la rouille !]
Le dépôt \textbf{rouge-orangé} observé sur le zinc est du \textbf{cuivre métallique}, et non de la rouille. La rouille est un oxyde de fer hydraté, de couleur \textbf{brun-rouille}, qui ne se forme que sur le fer en présence de dioxygène et d'eau. Confondre les deux est une erreur classique au baccalauréat : regarde la couleur et demande-toi quel métal est présent dans la solution.
\end{attention}

\courssub{Généralisation : d'autres couples métalliques}

\textbf{Le fer dans une solution de cuivre(II).} Plongeons maintenant un clou en fer dans la même solution bleue de \ce{CuSO4}. On observe le même phénomène : dépôt rouge de cuivre sur le clou, décoloration de la solution, qui prend peu à peu une teinte verdâtre caractéristique des ions \ce{Fe^2+}. L'équation-bilan est :

\[ \ce{Fe(s) + Cu^2+(aq) -> Fe^2+(aq) + Cu(s)} \]

Le fer (couple \ce{Fe^2+}/\ce{Fe}) est oxydé, l'ion \ce{Cu^2+} est réduit. Le fer, comme le zinc, réduit donc les ions cuivre(II).

\textbf{Le zinc dans une solution de fer(II).} Une lame de zinc plongée dans une solution de sulfate de fer(II) \ce{FeSO4} se recouvre d'un dépôt \textbf{gris métallique} de fer, tandis que la teinte verdâtre de la solution s'atténue :

\[ \ce{Zn(s) + Fe^2+(aq) -> Zn^2+(aq) + Fe(s)} \]

Le zinc réduit aussi les ions fer(II).

\textbf{Le fer dans une solution de zinc(II) : rien ne se passe.} En revanche, un clou en fer plongé dans une solution de \ce{ZnSO4} reste \emph{rigoureusement inchangé}, même après plusieurs jours : aucun dépôt, aucune coloration nouvelle. La réaction

\[ \ce{Fe(s) + Zn^2+(aq) ->[ne se produit pas] } \]

n'a \textbf{pas lieu}. La réaction entre un métal et un ion métallique ne se produit donc que \textbf{dans un seul sens} : le zinc réduit \ce{Fe^2+}, mais le fer ne réduit pas \ce{Zn^2+}.

\courssub{Prévision d'une réaction : le classement des pouvoirs réducteurs}

Ces expériences permettent de classer qualitativement les métaux selon leur \cle{pouvoir réducteur}, c'est-à-dire leur aptitude à céder des électrons :

\begin{center}
\begin{tabularx}{\linewidth}{|l|X|}
\hline
\rowcolor{popblueL} \textbf{Expérience} & \textbf{Conclusion} \\
\hline
\ce{Zn} réduit \ce{Fe^2+} & le zinc est plus réducteur que le fer \\
\hline
\ce{Fe} réduit \ce{Cu^2+} & le fer est plus réducteur que le cuivre \\
\hline
\ce{Fe} ne réduit pas \ce{Zn^2+} & confirmation : Zn $>$ Fe \\
\hline
\end{tabularx}
\end{center}

D'où le classement par pouvoir réducteur décroissant :

\[ \ce{Zn} \;>\; \ce{Fe} \;>\; \ce{Cu} \;>\; \ce{Ag} \;>\; \ce{Au} \]

\begin{methode}[Prévoir si une réaction a lieu entre un métal M et un ion M'$^{2+}$]
\begin{enumerate}
\item Identifier les deux couples mis en jeu : \ce{M^2+}/\ce{M} et \ce{\text{M'}^2+}/\ce{\text{M'}}.
\item Comparer les \cle{pouvoirs réducteurs} des deux métaux à l'aide du classement expérimental.
\item Conclure : la réaction \textbf{a lieu} si le métal M introduit est \textbf{plus réducteur} que le métal M' de l'ion ; le métal le plus réducteur cède ses électrons à l'ion de l'autre couple. Sinon, il ne se passe \textbf{rien}.
\item Écrire les deux demi-équations électroniques, puis l'équation-bilan en veillant à ce que les électrons échangés se compensent.
\end{enumerate}
\end{methode}

\begin{exemplebox}[Prévision rapide]
Peut-on conserver une solution de sulfate de cuivre(II) dans un récipient en fer ? Non : le fer étant plus réducteur que le cuivre, la réaction \ce{Fe + Cu^2+ -> Fe^2+ + Cu} attaquerait le récipient. En revanche, on peut sans risque conserver une solution de \ce{ZnSO4} dans un récipient en fer, car le fer est moins réducteur que le zinc.
\end{exemplebox}

\begin{exoresolu}[Le clou qui récupère le cuivre]
\textbf{Énoncé.} On plonge un clou en fer de masse $m = \SI{5,6}{g}$ dans \SI{200}{mL} d'une solution de sulfate de cuivre(II) de concentration $C = \SI{0,50}{mol.L^{-1}}$. On donne $M(\ce{Fe}) = \SI{56}{g.mol^{-1}}$ et $M(\ce{Cu}) = \SI{63,5}{g.mol^{-1}}$. Déterminer le réactif limitant, la masse de cuivre déposée et la masse de fer restante.
\tcblower
\textbf{Solution.} Équation-bilan : $\ce{Fe + Cu^2+ -> Fe^2+ + Cu}$.

Quantités initiales :
\[ n(\ce{Fe}) = \frac{5,6}{56} = \SI{0,10}{mol} \qquad n(\ce{Cu^2+}) = C \times V = 0,50 \times 0,200 = \SI{0,10}{mol} \]

Tableau d'avancement (en mol) :

\begin{center}
\begin{tabular}{lcccc}
\hline
 & \ce{Fe} & \ce{Cu^2+} & \ce{Fe^2+} & \ce{Cu} \\
\hline
État initial & 0,10 & 0,10 & 0 & 0 \\
État final & $0,10 - x$ & $0,10 - x$ & $x$ & $x$ \\
\hline
\end{tabular}
\end{center}

Les deux réactifs sont en proportions stœchiométriques : $x_{max} = \SI{0,10}{mol}$, tout est consommé.

Masse de cuivre déposée :
\[ m(\ce{Cu}) = 0,10 \times 63,5 = \SI{6,35}{g} \]

Masse de fer restante : $m(\ce{Fe}) = 0$ : le clou est entièrement consommé (s'il est resté assez longtemps dans la solution, il s'est dissous en totalité en libérant \SI{6,35}{g} de cuivre rouge).
\end{exoresolu}

\begin{exemplebox}[L'hydrométallurgie du cuivre]
La réaction $\ce{Fe + Cu^2+ -> Fe^2+ + Cu}$ est exploitée industriellement en \cle{hydrométallurgie} : on fait percoler des solutions minières pauvres contenant des ions \ce{Cu^2+} sur de la ferraille (vieux rails, tôles usagées). Le cuivre se dépose sur le fer et l'on récupère ainsi un métal précieux à partir de solutions très diluées, à faible coût. C'est le même principe que ton clou dans le bécher, mais à l'échelle de milliers de tonnes !
\end{exemplebox}

\courssub{Application : la protection du fer par le zinc et la formation de la rouille}

\textbf{La rouille : une oxydoréduction de tous les jours.} Un clou, une tôle, un portail en fer exposés à l'air humide se recouvrent peu à peu d'une couche brun-rouille friable : la \cle{rouille}, un oxyde de fer hydraté de formule approchée \ce{Fe2O3, xH2O}. Le fer est \textbf{oxydé} par le \textbf{dioxygène de l'air} en présence d'\textbf{eau} (les deux sont indispensables : un clou dans de l'huile bouillie ou dans l'air sec ne rouille pas). Schématiquement, en présence d'eau :

\[ \ce{4Fe + 3O2 + 6H2O -> 4Fe(OH)3 -> 2Fe2O3, 3H2O} \]

Le fer perd des électrons (oxydation), le dioxygène les capte (réduction). La rouille, contrairement à la couche d'alumine sur l'aluminium, est \emph{poreuse} : elle ne protège pas le métal, et la corrosion progresse jusqu'à la destruction complète de l'objet.

\textbf{Le zinc, gardien du fer : l'anode sacrificielle.} Le zinc étant \emph{plus réducteur} que le fer, il s'oxyde \emph{à la place} du fer dès que les deux métaux sont en contact en milieu humide : le zinc « se sacrifie » et le fer reste intact. C'est le principe de l'\cle{anode sacrificielle} : on fixe des blocs de zinc sur les coques en acier des bateaux (dans les ports de Douala et de Limbe, par exemple) ou sur les pipelines enterrés ; le zinc se corrode lentement et on le remplace périodiquement, tandis que la coque en fer est préservée.

\begin{popfigure}
\begin{tikzpicture}[scale=0.85]
% Clou non protégé
\draw[thick,popmuted] (0.3,0) -- (0.3,2.6);
\draw[thick,popmuted] (0.3,2.6) arc (180:0:0.45);
\draw[thick,popmuted] (1.2,2.6) -- (1.2,0.4);
\draw[thick,popmuted] (1.2,0.4) -- (0.75,-0.3) -- (0.3,0);
\fill[poporange!70] (0.15,0.1) rectangle (1.35,1.1);
\node at (0.75,1.6) {\small\ce{Fe}};
\node[poporange!90!black] at (0.75,0.55) {\small rouille};
\node at (0.75,-0.9) {\small\textbf{Clou non protégé}};
\node at (0.75,-1.4) {\small \ce{Fe} oxydé par \ce{O2} + \ce{H2O}};
% Clou protégé
\begin{scope}[xshift=6.5cm]
\draw[thick,popmuted] (0.3,0) -- (0.3,2.6);
\draw[thick,popmuted] (0.3,2.6) arc (180:0:0.45);
\draw[thick,popmuted] (1.2,2.6) -- (1.2,0.4);
\draw[thick,popmuted] (1.2,0.4) -- (0.75,-0.3) -- (0.3,0);
\fill[popblue!50] (1.2,0.6) rectangle (1.9,1.6);
\node at (1.55,1.1) {\small\ce{Zn}};
\node at (0.75,1.6) {\small\ce{Fe}};
\draw[-{Stealth},thick,poppurple] (1.2,1.05) -- (0.4,1.05);
\node[poppurple] at (0.8,1.35) {\small $e^-$};
\node at (0.75,-0.9) {\small\textbf{Clou protégé}};
\node at (0.75,-1.4) {\small le zinc s'oxyde à la place du fer};
\end{scope}
\end{tikzpicture}
\legende{Corrosion d'un clou en fer et protection par un bloc de zinc (anode sacrificielle) : le zinc, plus réducteur, cède ses électrons et s'oxyde à la place du fer.}
\end{popfigure}

\begin{savaistu}[La galvanisation sauve les toits du Cameroun]
Les tôles ondulées qui couvrent tant de maisons à Yaoundé, Douala ou Bafoussam sont des tôles d'acier \textbf{galvanisées} : elles sont recouvertes d'une fine couche de zinc. Même si la couche est rayée, le zinc continue de protéger le fer par le mécanisme d'anode sacrificielle, car il est plus réducteur : c'est lui qui s'oxyde. C'est pourquoi une tôle galvanisée résiste des années aux pluies tropicales, alors qu'une tôle nue rouillerait en quelques mois. Le mot « galvanisation » honore le physicien italien Luigi Galvani, pionnier de l'électricité animale au XVIII\textsuperscript{e} siècle.
\end{savaistu}

\begin{aretenir}[L'essentiel de la section]
\begin{itemize}
\item Un métal plongé dans une solution contenant les ions d'un \textbf{autre} métal peut les réduire : \ce{Zn + Cu^2+ -> Zn^2+ + Cu}, \ce{Fe + Cu^2+ -> Fe^2+ + Cu}, \ce{Zn + Fe^2+ -> Zn^2+ + Fe}.
\item La réaction est \textbf{spontanée} et n'a lieu que \textbf{dans un sens} : le métal le \textbf{plus réducteur} cède ses électrons aux ions de l'autre couple.
\item Classement qualitatif des pouvoirs réducteurs : $\ce{Zn} > \ce{Fe} > \ce{Cu} > \ce{Ag} > \ce{Au}$.
\item La \textbf{rouille} résulte de l'oxydation du fer par le dioxygène en présence d'eau ; le zinc protège le fer par \textbf{anode sacrificielle} (galvanisation, coques de bateaux).
\end{itemize}
\end{aretenir}

\coursec{Autres réactions d'oxydoréduction en solution aqueuse}

Dans la section précédente, nous avons vu qu'un métal comme le zinc peut céder des électrons aux ions $\ce{Cu^2+}$ : c'est la réaction d'oxydoréduction entre un ion métallique et un métal. Mais l'oxydoréduction ne se limite pas aux métaux et à leurs ions ! En solution aqueuse, de nombreuses espèces dissoutes — anions oxygénés comme $\ce{MnO4-}$ ou $\ce{Cr2O7^2-}$, molécules organiques comme l'éthanol, molécules simples comme le diiode — échangent des électrons. Ces réactions, souvent spectaculaires car accompagnées de changements de couleur, sont à la base des \cle{dosages redox} (titrages) que tu réaliseras en travaux pratiques et qui tombent régulièrement au Baccalauréat TI. Étudions trois réactions classiques, incontournables.

\courssub{Réaction entre les ions fer(II) et les ions permanganate en milieu acide}

\textbf{Observation expérimentale.} Le permanganate de potassium, bien connu des infirmiers comme antiseptique, forme en solution une couleur \emph{violette} intense, due à l'ion \cle{permanganate} $\ce{MnO4-}$. Les ions fer(II) $\ce{Fe^2+}$, eux, donnent une solution vert très pâle, presque incolore. Lorsqu'on verse goutte à goutte une solution acidifiée de permanganate de potassium dans une solution de sulfate de fer(II), chaque goutte violette se \emph{décolore instantanément} au contact de la solution de fer(II) : la solution devient progressivement jaune-orangé, couleur des ions $\ce{Fe^3+}$ qui se forment.

\begin{experience}[Décoloration du permanganate par les ions fer(II)]
\textbf{Dispositif et protocole :} dans un erlenmeyer, introduire \SI{20,0}{mL} d'une solution de sulfate de fer(II) acidifiée avec quelques gouttes d'acide sulfurique dilué. Verser progressivement, à la burette, une solution de permanganate de potassium (\ce{K+} + \ce{MnO4-}) de couleur violette.\par
\textbf{Observations :} tant qu'il reste des ions $\ce{Fe^2+}$, chaque goutte de permanganate se décolore immédiatement. La solution vire lentement au jaune-orangé (ions $\ce{Fe^3+}$). Quand tout le fer(II) est consommé, la première goutte en excès donne à la solution une teinte rose persistante : c'est l'\emph{équivalence} du titrage.\par
\textbf{Interprétation :} l'ion $\ce{MnO4-}$ (violet) est transformé en ion $\ce{Mn^2+}$ (quasi incolore) : il est \emph{réduit}. L'ion $\ce{Fe^2+}$ est transformé en $\ce{Fe^3+}$ : il est \emph{oxydé}. C'est bien une réaction d'oxydoréduction.
\end{experience}

\textbf{Écriture des demi-équations.} Le couple du fer est simple à traiter : l'ion fer(II) cède un électron :
\[ \ce{Fe^2+ = Fe^3+ + e-} \]
Le couple $\ce{MnO4-}/\ce{Mn^2+}$ est plus délicat, car l'ion permanganate contient des atomes d'oxygène qu'il faut évacuer sous forme d'eau, en milieu acide. Sa demi-équation est :
\[ \ce{MnO4- + 8H+ + 5e- = Mn^2+ + 4H2O} \]

\begin{methode}[Équilibrer une demi-équation en milieu acide]
Pour équilibrer une demi-équation électronique en milieu acide, procède toujours dans le même ordre :
\begin{enumerate}
\item \textbf{Équilibrer l'élément caractéristique} (autre que O et H) : ici un atome de Mn de chaque côté.
\item \textbf{Équilibrer l'oxygène avec des molécules d'eau} $\ce{H2O}$ : $\ce{MnO4-}$ possède 4 atomes d'oxygène, on écrit donc $4\,\ce{H2O}$ à droite.
\item \textbf{Équilibrer l'hydrogène avec des ions} $\ce{H+}$ : $4\,\ce{H2O}$ apportent 8 atomes d'hydrogène, on écrit donc $8\,\ce{H+}$ à gauche.
\item \textbf{Équilibrer les charges avec des électrons} $\ce{e-}$ : à gauche, la charge totale est $-1 + 8 = +7$ ; à droite, elle vaut $+2$. Il faut donc ajouter $5$ électrons à gauche pour que $+7 - 5 = +2$.
\item \textbf{Vérifier} la conservation des éléments et des charges.
\end{enumerate}
\end{methode}

\textbf{Équation-bilan.} Pour que les électrons n'apparaissent pas dans le bilan, il faut que le nombre d'électrons cédés par $\ce{Fe^2+}$ soit égal au nombre d'électrons captés par $\ce{MnO4-}$. On multiplie donc la demi-équation du fer par 5 :
\begin{align*}
\ce{MnO4- + 8H+ + 5e-} &\ce{= Mn^2+ + 4H2O} \\
\ce{5 Fe^2+} &\ce{= 5 Fe^3+ + 5 e-} \\ \intertext{\hrulefill}
\ce{MnO4- + 5Fe^2+ + 8H+} &\ce{-> Mn^2+ + 5Fe^3+ + 4H2O}
\end{align*}

\begin{exemplebox}[Vérification de la conservation des charges]
Une équation-bilan correcte doit conserver \emph{à la fois} les éléments et la charge électrique. Vérifions pour :
\[ \ce{MnO4- + 5Fe^2+ + 8H+ -> Mn^2+ + 5Fe^3+ + 4H2O} \]
\begin{itemize}
\item \textbf{Charge à gauche :} $(-1) + 5\times(+2) + 8\times(+1) = -1 + 10 + 8 = +17$.
\item \textbf{Charge à droite :} $(+2) + 5\times(+3) + 0 = +17$.
\end{itemize}
La charge $+17$ est conservée : l'équation est correctement équilibrée. Les éléments le sont aussi : 1 Mn, 5 Fe, 8 H et 4 O de chaque côté.
\end{exemplebox}

\begin{attention}[Erreur fréquente : oublier les $\ce{H+}$ et les $\ce{H2O}$]
L'erreur la plus classique au Bac est d'écrire $\ce{MnO4- + 5e- -> Mn^2+}$ en oubliant les $8\,\ce{H+}$ et les $4\,\ce{H2O}$. Cette écriture est \textbf{fausse} : ni les atomes d'oxygène ni les charges ne sont conservés. Retiens le réflexe : dès qu'une espèce contient de l'oxygène ($\ce{MnO4-}$, $\ce{Cr2O7^2-}$, $\ce{NO3-}$\ldots), la demi-équation en milieu acide fait intervenir $\ce{H+}$ et $\ce{H2O}$. Autre piège : le milieu doit être \emph{acide} ; sans acide sulfurique ajouté, la réaction ne se produit pas correctement (formation de $\ce{MnO2}$ brun en milieu neutre/basique).
\end{attention}

\courssub{Réaction entre l'éthanol et les ions dichromate : le principe de l'alcootest}

\textbf{Observation expérimentale.} L'ion \cle{dichromate} $\ce{Cr2O7^2-}$ donne aux solutions une couleur \emph{orange}. En milieu acide, en présence d'éthanol $\ce{CH3CH2OH}$, la solution vire au \emph{vert} : c'est la couleur des ions chrome(III) $\ce{Cr^3+}$ qui se forment. Ce virage spectaculaire de l'orange au vert est l'une des réactions les plus célèbres de la chimie appliquée.

\begin{popfigure}
\begin{center}
\begin{tikzpicture}[scale=0.9, every node/.style={font=\small}]
% Tube 1 : dichromate orange
\draw[thick, popdark] (0,0) -- (0,3) (1.4,0) -- (1.4,3);
\draw[thick, popdark] (0,0) arc (180:360:0.7);
\fill[poporange!70] (0.02,0.15) arc (180:360:0.68) -- (1.38,2.0) -- (0.02,2.0) -- cycle;
\draw[popdark] (0,2.0) -- (1.4,2.0);
\node[below, popdark] at (0.7,-0.5) {$\ce{Cr2O7^2-}$ (orange)};
\node[above, popdark] at (0.7,3.1) {Tube 1};
% Flèche avec éthanol
\draw[-{Stealth[length=3mm]}, very thick, popblue] (1.9,1.6) -- (3.3,1.6);
\node[above, popblue] at (2.6,1.7) {$+$ éthanol};
\node[below, popblue] at (2.6,1.4) {milieu acide};
% Tube 2 : vert
\draw[thick, popdark] (3.8,0) -- (3.8,3) (5.2,0) -- (5.2,3);
\draw[thick, popdark] (3.8,0) arc (180:360:0.7);
\fill[popgreen!60] (3.82,0.15) arc (180:360:0.68) -- (5.18,2.0) -- (3.82,2.0) -- cycle;
\draw[popdark] (3.8,2.0) -- (5.2,2.0);
\node[below, popdark] at (4.5,-0.5) {$\ce{Cr^3+}$ (vert)};
\node[above, popdark] at (4.5,3.1) {Tube 2};
\end{tikzpicture}
\end{center}
\legende{Virage colorimétrique du dichromate : la solution orange vire au vert lorsque l'éthanol est oxydé en acide éthanoïque en milieu acide.}
\end{popfigure}

\textbf{Les couples mis en jeu.} L'ion dichromate est l'oxydant du couple $\ce{Cr2O7^2-}/\ce{Cr^3+}$. Sa demi-équation s'établit avec la méthode vue plus haut (7 atomes d'oxygène $\Rightarrow$ $7\,\ce{H2O}$, donc $14\,\ce{H+}$, puis 6 électrons) :
\[ \ce{Cr2O7^2- + 14H+ + 6e- = 2Cr^3+ + 7H2O} \]
L'éthanol est le réducteur du couple $\ce{CH3COOH}/\ce{C2H5OH}$ : il est oxydé en \cle{acide éthanoïque} (l'acide du vinaigre !) :
\[ \ce{CH3CH2OH + H2O = CH3COOH + 4H+ + 4e-} \]
Pour combiner les deux demi-équations, on cherche le plus petit commun multiple de 6 et 4, soit 12 : on multiplie la première par 2 et la seconde par 3 :
\[ \ce{2Cr2O7^2- + 3CH3CH2OH + 16H+ -> 4Cr^3+ + 3CH3COOH + 11H2O} \]
Vérifions les charges : à gauche, $2\times(-2) + 16 = +12$ ; à droite, $4\times(+3) = +12$. L'équation est correcte.

\begin{savaistu}[L'éthylotest : la chimie au service de la sécurité routière]
Les premiers alcootests chimiques (éthylotests), utilisés par les forces de l'ordre — y compris au Cameroun lors des contrôles routiers sur l'axe Yaoundé--Douala — reposaient exactement sur cette réaction. Le conducteur soufflait dans un ballon ; l'air expiré traversait un tube contenant du dichromate de potassium acidifié, de couleur orange. Si l'air contenait des vapeurs d'éthanol, le réactif virait au vert, et la longueur de la zone verdie indiquait la quantité d'alcool. Aujourd'hui, les éthylotests électroniques ont pris le relais, mais le principe historique reste un magnifique exemple d'application directe de l'oxydoréduction. À noter : l'oxydation de l'éthanol en acide éthanoïque est aussi ce qui transforme le vin de palme en vinaigre lorsqu'on le laisse trop longtemps à l'air !
\end{savaistu}

\courssub{Réaction entre le diiode et les ions thiosulfate}

\textbf{Observation expérimentale.} Le diiode $\ce{I2}$ dissous dans l'eau (en présence d'iodure de potassium pour le solubiliser sous forme de $\ce{I3-}$) donne une solution \emph{brune} (ou jaune-orangé selon la concentration). L'ion \cle{thiosulfate} $\ce{S2O3^2-}$ est incolore. Quand on ajoute une solution de thiosulfate de sodium à une solution de diiode, la coloration brune \emph{s'atténue progressivement} puis disparaît totalement : le diiode a été entièrement consommé.

\textbf{Les couples mis en jeu.} Le diiode est l'oxydant du couple $\ce{I2}/\ce{I-}$ ; il est réduit en ions iodure incolores :
\[ \ce{I2 + 2e- = 2I-} \]
L'ion thiosulfate est le réducteur du couple $\ce{S4O6^2-}/\ce{S2O3^2-}$ ; il est oxydé en ion \cle{tétrathionate} :
\[ \ce{2S2O3^2- = S4O6^2- + 2e-} \]
Les deux demi-équations échangent le même nombre d'électrons (2) : on les additionne directement :
\[ \ce{I2 + 2S2O3^2- -> 2I- + S4O6^2-} \]
Cette réaction est \emph{rapide}, \emph{totale} et \emph{unique} : elle possède toutes les qualités requises pour servir de \emph{réaction support de titrage}.

\textbf{Le rôle de l'empois d'amidon.} Comment repérer précisément la fin de la réaction ? La décoloration du diiode est visible, mais peu nette quand la solution devient très pâle. On ajoute alors quelques gouttes d'\cle{empois d'amidon} : en présence de diiode, il forme un complexe d'un \emph{bleu nuit} intense, très facile à observer. À l'équivalence, dès que la dernière trace de $\ce{I2}$ disparaît, la solution passe brutalement du bleu nuit à l'incolore. L'empois d'amidon joue donc le rôle d'\emph{indicateur de fin de réaction}.

\begin{popfigure}
\begin{center}
\begin{tikzpicture}[scale=0.95, every node/.style={font=\small}]
% Support et burette
\draw[very thick, popdark] (-2.6,0) -- (-2.6,5.6);
\draw[very thick, popdark] (-3.2,0) -- (-2.0,0);
\draw[thick, popdark] (-2.6,4.6) -- (-1.6,4.6);
% Burette
\draw[thick, popdark] (-1.45,2.9) -- (-1.45,5.4);
\draw[thick, popdark] (-1.05,2.9) -- (-1.05,5.4);
\fill[popblueL] (-1.43,3.6) rectangle (-1.07,5.35);
\draw[popdark] (-1.43,3.6) -- (-1.07,3.6);
% graduations
\foreach \y in {3.9,4.3,4.7,5.1} \draw[popdark] (-1.45,\y) -- (-1.35,\y);
% robinet et pointe
\draw[thick, popdark] (-1.45,2.9) -- (-1.45,2.5) -- (-1.25,2.3);
\draw[thick, popdark] (-1.05,2.9) -- (-1.05,2.5) -- (-1.25,2.3);
\draw[thick, popdark] (-1.7,2.62) -- (-0.8,2.62);
\draw[thick, popdark] (-1.25,2.3) -- (-1.25,1.9);
% goutte
\fill[popblue] (-1.25,1.55) circle (0.07);
\node[right, popdark, align=left] at (-0.5,4.4) {burette : solution de\\ thiosulfate de sodium};
% Erlenmeyer
\draw[thick, popdark] (-2.35,0) -- (-2.35,0.9) -- (-1.6,1.9);
\draw[thick, popdark] (-0.15,0) -- (-0.15,0.9) -- (-0.9,1.9);
\draw[thick, popdark] (-1.6,1.9) -- (-1.6,2.2) (-0.9,1.9) -- (-0.9,2.2);
\draw[thick, popdark] (-1.6,2.2) arc (120:60:0.8);
% liquide bleu nuit (calculé pour être à l'intérieur des parois)
% Paroi gauche : x = -2.35 + y*(0.75/1.9). À y=0.75 -> x = -2.05
% Paroi droite : x = -0.15 - y*(0.75/1.9). À y=0.75 -> x = -0.45
\fill[popblue!85!black] (-2.05,0.05) -- (-2.05,0.75) -- (-0.45,0.75) -- (-0.45,0.05) -- cycle;
\draw[popdark] (-2.05,0.75) -- (-0.45,0.75);
\node[right, popdark, align=left] at (0.3,0.6) {erlenmeyer : diiode $+$\\ empois d'amidon (bleu nuit)};
\end{tikzpicture}
\end{center}
\legende{Montage de titrage iodométrique : la burette contient la solution de thiosulfate ; l'erlenmeyer contient le diiode et quelques gouttes d'empois d'amidon. À l'équivalence, la solution passe du bleu nuit à l'incolore.}
\end{popfigure}

\begin{exoresolu}[Titrage du diiode par le thiosulfate]
\textbf{Énoncé.} On dose une solution de diiode par une solution de thiosulfate de sodium de concentration $C = \SI{0,100}{mol.L^{-1}}$. L'équivalence est repérée par la décoloration de l'empois d'amidon pour un volume versé $V_E = \SI{12,4}{mL}$. Calculer la quantité de matière de diiode présente dans l'erlenmeyer.
\tcblower
\textbf{Solution détaillée.} L'équation du titrage est :
\[ \ce{I2 + 2S2O3^2- -> 2I- + S4O6^2-} \]
\textbf{Étape 1 : quantité de thiosulfate versée à l'équivalence.}
\[ n(\ce{S2O3^2-}) = C \times V_E = 0,100 \times 12,4 \times 10^{-3} = \SI{1,24e-3}{mol} \]
\textbf{Étape 2 : exploitation de la stœchiométrie.} D'après les coefficients de l'équation, 1 mole de diiode réagit avec 2 moles d'ions thiosulfate. À l'équivalence, les réactifs ont été mélangés dans les proportions stœchiométriques :
\[ n(\ce{I2}) = \frac{n(\ce{S2O3^2-})}{2} = \frac{1,24 \times 10^{-3}}{2} = \SI{6,2e-4}{mol} \]
\textbf{Conclusion :} l'erlenmeyer contenait $\SI{6,2e-4}{mol}$ de diiode. La relation $n(\ce{I2}) = \dfrac{n(\ce{S2O3^2-})}{2}$ est une relation stœchiométrique à connaître parfaitement pour le Bac : elle revient dans presque tous les exercices de titrage iodométrique.
\end{exoresolu}

\courssub{Pourquoi ces réactions sont-elles si importantes ?}

Ces trois réactions ne sont pas de simples curiosités de laboratoire : elles illustrent les conditions qu'une réaction doit remplir pour servir de \emph{support de titrage}, et elles seront au cœur de tes séances de travaux pratiques.

\begin{aretenir}[L'essentiel de la section]
\begin{itemize}
\item En milieu acide, l'ion permanganate $\ce{MnO4-}$ (violet) oxyde les ions $\ce{Fe^2+}$ en $\ce{Fe^3+}$ et est réduit en $\ce{Mn^2+}$ (quasi incolore) :
\[ \ce{MnO4- + 5Fe^2+ + 8H+ -> Mn^2+ + 5Fe^3+ + 4H2O} \]
\item En milieu acide, l'ion dichromate $\ce{Cr2O7^2-}$ (orange) oxyde l'éthanol en acide éthanoïque et est réduit en $\ce{Cr^3+}$ (vert) : c'est le principe de l'alcootest chimique.
\[ \ce{2Cr2O7^2- + 3CH3CH2OH + 16H+ -> 4Cr^3+ + 3CH3COOH + 11H2O} \]
\item Le diiode $\ce{I2}$ (brun) oxyde les ions thiosulfate et est réduit en ions iodure $\ce{I-}$ (incolores) :
\[ \ce{I2 + 2S2O3^2- -> 2I- + S4O6^2-} \]
L'empois d'amidon (bleu nuit en présence de $\ce{I2}$) sert d'indicateur de fin de réaction.
\item Pour équilibrer une demi-équation en milieu acide : équilibrer O avec $\ce{H2O}$, puis H avec $\ce{H+}$, puis les charges avec $\ce{e-}$.
\item À l'équivalence d'un titrage, les réactifs sont dans les proportions stœchiométriques : par exemple $n(\ce{I2}) = n(\ce{S2O3^2-})/2$.
\end{itemize}
\end{aretenir}

\begin{exercice}[À toi de jouer]
On acidifie un mélange contenant des ions dichromate et on y introduit de l'éthanol en excès.\par
1. Écrire les deux demi-équations électroniques des couples $\ce{Cr2O7^2-}/\ce{Cr^3+}$ et $\ce{CH3COOH}/\ce{C2H5OH}$.\par
2. En déduire l'équation-bilan de la réaction et vérifier la conservation des charges.\par
3. On titre ensuite une solution de diiode par une solution de thiosulfate de concentration $\SI{5,0e-2}{mol.L^{-1}}$. L'équivalence est obtenue pour $V_E = \SI{16,0}{mL}$. Calculer la quantité de diiode dosée.
\end{exercice}

\begin{corrige}[Exercice]
1. Couple du dichromate : $\ce{Cr2O7^2- + 14H+ + 6e- = 2Cr^3+ + 7H2O}$. Couple de l'éthanol : $\ce{CH3CH2OH + H2O = CH3COOH + 4H+ + 4e-}$.\par
2. On multiplie la première par 2 et la seconde par 3 (ppcm de 6 et 4 égal à 12 électrons) :
\[ \ce{2Cr2O7^2- + 3CH3CH2OH + 16H+ -> 4Cr^3+ + 3CH3COOH + 11H2O} \]
Charges : à gauche $2\times(-2) + 16 = +12$ ; à droite $4\times(+3) = +12$. La charge est conservée.\par
3. Quantité de thiosulfate : $n(\ce{S2O3^2-}) = C \times V_E = 5,0\times10^{-2} \times 16,0\times10^{-3} = \SI{8,0e-4}{mol}$. D'après la stœchiométrie :
\[ n(\ce{I2}) = \frac{n(\ce{S2O3^2-})}{2} = \SI{4,0e-4}{mol} \]
\end{corrige}

Tu maîtrises désormais les trois réactions d'oxydoréduction fondamentales en solution aqueuse. Mais comment reconnaître une oxydoréduction lorsqu'aucun transfert d'électrons n'est évident, par exemple dans les réactions par voie sèche ? C'est toute la puissance de la notion de \emph{nombre d'oxydation}, que nous allons découvrir dans la prochaine section.

\coursec{Généralisation : les nombres d'oxydation}

Dans les sections précédentes, nous avons étudié des réactions d'oxydoréduction en solution aqueuse : action des ions permanganate sur les ions fer(II), du dichromate sur l'éthanol, du diiode sur le thiosulfate. Dans tous ces cas, nous avons pu écrire des \cle{demi-équations électroniques} mettant en jeu un transfert d'électrons bien visible. Mais que dire de la combustion du méthane dans le dioxygène, ou de la formation de la rouille à l'air libre ? Aucun ion, aucun transfert d'électrons évident n'apparaît. Pourtant, ces réactions \emph{sont} des oxydoréductions. Pour les reconnaître et les équilibrer, les chimistes ont inventé un outil remarquable : le \cle{nombre d'oxydation}. C'est l'objet de cette section, qui te donnera une définition \emph{générale} de l'oxydoréduction, valable en solution aqueuse comme par voie sèche.

\courssub{Les limites de la définition électronique}

Considérons la combustion du méthane, gaz utilisé dans les cuisinières à gaz de nos villes :

\[ \ce{CH4 + 2 O2 -> CO2 + 2 H2O} \]

Cette réaction dégage beaucoup d'énergie, elle est totale et rapide une fois amorcée. Mais où sont les électrons transférés ? Le méthane \ce{CH4} et le dioxygène \ce{O2} sont des molécules \emph{neutres}, liées par des liaisons \emph{covalentes} : aucun ion n'intervient, aucune demi-équation électronique ne s'écrit naturellement. Pourtant, le carbone a bien été « oxydé » (il s'est combiné à l'oxygène) et l'oxygène a été « réduit ».

De même, la formation de la rouille sur les portails métalliques de nos concessions pendant la saison des pluies :

\[ \ce{4 Fe + 3 O2 -> 2 Fe2O3} \]

met en jeu des espèces neutres. La définition électronique (perte ou gain réel d'électrons) atteint ici ses limites. Il faut un point de vue plus large, fondé sur une \emph{charge fictive} attribuée à chaque atome : le nombre d'oxydation.

\courssub{Définition du nombre d'oxydation}

L'idée est simple : dans une liaison covalente, les électrons sont \emph{partagés} entre deux atomes, mais pas toujours équitablement. L'atome le plus \cle{électronégatif} attire davantage le doublet de liaison. Faisons alors une fiction comptable : attribuons \emph{tous} les électrons de chaque liaison à l'atome le plus électronégatif, comme si toutes les liaisons étaient devenues ioniques. La charge que porterait alors chaque atome est son nombre d'oxydation.

\begin{definition}[Nombre d'oxydation]
Le \cle{nombre d'oxydation} (noté \cle{n.o.}) d'un atome dans une espèce chimique est la \textbf{charge électrique fictive} qu'il porterait si les électrons de chaque liaison covalente étaient entièrement attribués à l'atome le plus électronégatif des deux. Il s'exprime en \textbf{chiffres romains} précédés du signe : $+\mathrm{I}$, $-\mathrm{II}$, $+\mathrm{VI}$, etc.
\end{definition}

\begin{attention}[Ne pas confondre n.o. et charge réelle]
Le nombre d'oxydation est une \textbf{charge fictive}, pas une charge réelle. On l'écrit en \textbf{chiffres romains avec le signe devant} : $+\mathrm{II}$, $-\mathrm{II}$, $+\mathrm{VII}$. La charge réelle d'un ion s'écrit en chiffres arabes avec le signe après : \ce{Fe^2+}, \ce{MnO4-}. Écrire « n.o. du manganèse = 7+ » est une erreur classique au baccalauréat : la bonne écriture est $+\mathrm{VII}$.
\end{attention}

\courssub{Règles de calcul du nombre d'oxydation}

Les règles suivantes, appliquées dans l'ordre, permettent de déterminer le n.o. de n'importe quel élément dans n'importe quelle espèce.

\begin{aretenir}[Règles de détermination du nombre d'oxydation]
\begin{enumerate}
\item Le n.o. d'un élément dans un \textbf{corps simple} est nul : n.o.$(\mathrm{Fe}) = 0$, n.o.$(\mathrm{O}$ dans $\ce{O2}) = 0$, n.o.$(\mathrm{H}$ dans $\ce{H2}) = 0$.
\item Le n.o. d'un \textbf{ion monoatomique} est égal à sa charge : n.o.$(\ce{Fe^2+}) = +\mathrm{II}$, n.o.$(\ce{Cl-}) = -\mathrm{I}$.
\item Le n.o. de l'\textbf{oxygène} vaut $-\mathrm{II}$ dans presque tous les composés. \emph{Exception} : dans les peroxydes (comme l'eau oxygénée \ce{H2O2}), il vaut $-\mathrm{I}$.
\item Le n.o. de l'\textbf{hydrogène} vaut $+\mathrm{I}$ dans presque tous les composés. \emph{Exception} : dans les hydrures métalliques (comme \ce{NaH}), il vaut $-\mathrm{I}$.
\item La \textbf{somme des n.o.} de tous les atomes est \textbf{nulle} dans une molécule neutre, et \textbf{égale à la charge} de l'ion dans un ion polyatomique.
\end{enumerate}
\end{aretenir}

\begin{exemplebox}[Calculs de nombres d'oxydation]
\textbf{a) Soufre dans l'acide sulfurique \ce{H2SO4}.} Molécule neutre : $2 \times (+\mathrm{I}) + x + 4 \times (-\mathrm{II}) = 0$, soit $2 + x - 8 = 0$, d'où $x = +6$ : n.o.$(\mathrm{S}) = +\mathrm{VI}$.

\textbf{b) Manganèse dans l'ion permanganate \ce{MnO4-}.} Ion de charge $-1$ : $x + 4 \times (-2) = -1$, soit $x = +7$ : n.o.$(\mathrm{Mn}) = +\mathrm{VII}$.

\textbf{c) Chrome dans l'ion dichromate \ce{Cr2O7^2-}.} $2x + 7 \times (-2) = -2$, soit $2x = 12$, d'où $x = +6$ : n.o.$(\mathrm{Cr}) = +\mathrm{VI}$.

\textbf{d) Carbone dans l'éthanol \ce{C2H5OH}, de formule développée \chemfig{CH_3-CH_2-OH}.} Molécule neutre : $2x + 6 \times (+1) + (-2) = 0$, soit $2x = -4$, d'où $x = -2$ en \emph{moyenne} : n.o. moyen$(\mathrm{C}) = -\mathrm{II}$. (Atome par atome : le carbone du groupe \ce{CH3} est à $-\mathrm{III}$, celui lié à l'oxygène à $-\mathrm{I}$.)
\end{exemplebox}

\begin{savaistu}[Un même élément, plusieurs degrés d'oxydation]
Un même élément peut exister à plusieurs nombres d'oxydation. Le fer, par exemple, existe aux degrés $+\mathrm{II}$ (ion \ce{Fe^2+}, vert pâle en solution) et $+\mathrm{III}$ (ion \ce{Fe^3+}, rouille orangée) : c'est ce passage de l'un à l'autre qui explique la couleur de la rouille sur les tôles et les carrosseries. Le manganèse est encore plus riche : $+\mathrm{II}$ (\ce{Mn^2+}, presque incolore), $+\mathrm{IV}$ (\ce{MnO2}, brun, dans les piles), $+\mathrm{VII}$ (\ce{MnO4-}, violet intense). C'est cette variété qui fait la richesse de la chimie des métaux de transition étudiée à la SONARA ou dans les laboratoires d'analyse.
\end{savaistu}

\courssub{Oxydoréduction : les nouvelles définitions}

Le nombre d'oxydation permet de reformuler toutes les définitions de l'oxydoréduction de manière \emph{générale}.

\begin{definition}[Oxydation, réduction, oxydant, réducteur (point de vue des n.o.)]
\begin{itemize}
\item Une \cle{oxydation} est une \textbf{augmentation} du nombre d'oxydation d'un élément.
\item Une \cle{réduction} est une \textbf{diminution} du nombre d'oxydation d'un élément.
\item Un \cle{oxydant} est une espèce dont un élément voit son n.o. \textbf{diminuer} (il se réduit).
\item Un \cle{réducteur} est une espèce dont un élément voit son n.o. \textbf{augmenter} (il s'oxyde).
\end{itemize}
Une réaction est une \cle{oxydoréduction} si et seulement si \textbf{au moins un élément change de nombre d'oxydation} au cours de la réaction.
\end{definition}

Vérifions la cohérence avec la définition électronique : perdre des électrons (négatifs) fait \emph{augmenter} la charge fictive, donc le n.o. — les deux points de vue coïncident. Reprenons le couple \ce{MnO4-}/\ce{Mn^2+} : le manganèse passe de $+\mathrm{VII}$ à $+\mathrm{II}$, son n.o. diminue de 5 unités, ce qui correspond exactement à la capture de 5 électrons dans la demi-équation $\ce{MnO4- + 8 H+ + 5 e- -> Mn^2+ + 4 H2O}$. La variation de n.o. \emph{compte les électrons échangés}.

\begin{center}
\begin{tabularx}{\linewidth}{|l|X|X|X|}
\hline
\rowcolor{popblueL} \textbf{Couple} & \textbf{n.o. initial} & \textbf{n.o. final} & \textbf{Variation (électrons)} \\
\hline
\ce{MnO4-}/\ce{Mn^2+} & Mn : $+\mathrm{VII}$ & Mn : $+\mathrm{II}$ & diminution de 5 $\rightarrow$ gain de $5\,e^-$ (réduction) \\
\hline
\ce{Fe^3+}/\ce{Fe^2+} & Fe : $+\mathrm{III}$ & Fe : $+\mathrm{II}$ & diminution de 1 $\rightarrow$ gain de $1\,e^-$ (réduction) \\
\hline
\ce{Fe^2+} oxydé en \ce{Fe^3+} & Fe : $+\mathrm{II}$ & Fe : $+\mathrm{III}$ & augmentation de 1 $\rightarrow$ perte de $1\,e^-$ (oxydation) \\
\hline
\ce{Cr2O7^2-}/\ce{Cr^3+} & Cr : $+\mathrm{VI}$ & Cr : $+\mathrm{III}$ & diminution de 3 par Cr (réduction) \\
\hline
\end{tabularx}
\end{center}

\begin{popfigure}
\begin{center}
\begin{tikzpicture}[scale=1]
% Axe gradué
\draw[->, thick, popdark] (-4.6,0) -- (4.9,0) node[below left=2pt and 0pt, font=\small] {nombre d'oxydation};
\foreach \x/\lab in {-4/{$-\mathrm{IV}$},-3/{$-\mathrm{III}$},-2/{$-\mathrm{II}$},-1/{$-\mathrm{I}$},0/{$0$},1/{$+\mathrm{I}$},2/{$+\mathrm{II}$},3/{$+\mathrm{III}$},4/{$+\mathrm{IV}$}}{
  \draw[popdark] (\x,0.08) -- (\x,-0.08);
  \node[below, font=\footnotesize, popdark] at (\x,-0.08) {\lab};
}
% Flèche oxydation (vers la droite)
\draw[->, very thick, poporange] (-2,0.55) -- (2,0.55);
\node[above, poporange, font=\small\bfseries] at (0,0.55) {Oxydation : le n.o. augmente};
\node[below, poporange, font=\footnotesize] at (0,0.42) {perte d'électrons (réducteur)};
% Flèche réduction (vers la gauche)
\draw[->, very thick, popblue] (3,-0.75) -- (-1,-0.75);
\node[below, popblue, font=\small\bfseries] at (1,-0.78) {Réduction : le n.o. diminue};
\node[below, popblue, font=\footnotesize] at (1,-1.25) {gain d'électrons (oxydant)};
% Exemples ponctuels
\fill[popgreen] (2,0) circle (2.2pt); \node[above=8pt, font=\footnotesize, popgreen] at (2,0.1) {\ce{Fe^2+} : $+\mathrm{II}$};
\fill[popgreen] (3,0) circle (2.2pt); \node[above=8pt, font=\footnotesize, popgreen] at (3,0.1) {\ce{Fe^3+} : $+\mathrm{III}$};
\end{tikzpicture}
\end{center}
\legende{Échelle des nombres d'oxydation : toute augmentation du n.o. (vers la droite) est une oxydation ; toute diminution (vers la gauche) est une réduction.}
\end{popfigure}

\courssub{Méthode d'équilibrage par les nombres d'oxydation}

La grande force des nombres d'oxydation est de permettre d'\textbf{équilibrer} une équation d'oxydoréduction sans écrire les demi-équations électroniques. Le principe physique est la conservation de la charge : \emph{le nombre d'électrons cédés par le réducteur est égal au nombre d'électrons captés par l'oxydant}, autrement dit la variation totale de n.o. de l'élément oxydé compense exactement celle de l'élément réduit.

\begin{methode}[Équilibrer une équation d'oxydoréduction par les n.o.]
\begin{enumerate}
\item \textbf{Calculer} les n.o. de tous les éléments dans les réactifs et les produits.
\item \textbf{Repérer} les éléments dont le n.o. varie, et écrire les variations : $\Delta$ n.o. de l'élément oxydé (augmentation) et de l'élément réduit (diminution).
\item \textbf{Ajuster} les coefficients des espèces oxydante et réductrice pour que l'augmentation totale égale la diminution totale (électrons cédés = électrons captés).
\item \textbf{Équilibrer} ensuite les autres éléments (O, H, charges) comme dans une équation ordinaire.
\item \textbf{Vérifier} la conservation des atomes et des charges.
\end{enumerate}
\end{methode}

\begin{exoresolu}[Combustion du méthane : une oxydoréduction sans ions]
On considère la combustion complète du méthane : $\ce{CH4 + O2 -> CO2 + H2O}$. Montrer qu'il s'agit d'une oxydoréduction et équilibrer l'équation par la méthode des nombres d'oxydation.
\tcblower
\textbf{Étape 1 — Calcul des n.o.}
\begin{itemize}
\item Dans \ce{CH4} : $x + 4 \times (+1) = 0 \Rightarrow$ n.o.$(\mathrm{C}) = -\mathrm{IV}$ ; n.o.$(\mathrm{H}) = +\mathrm{I}$.
\item Dans \ce{O2} (corps simple) : n.o.$(\mathrm{O}) = 0$.
\item Dans \ce{CO2} : $x + 2 \times (-2) = 0 \Rightarrow$ n.o.$(\mathrm{C}) = +\mathrm{IV}$.
\item Dans \ce{H2O} : n.o.$(\mathrm{O}) = -\mathrm{II}$, n.o.$(\mathrm{H}) = +\mathrm{I}$.
\end{itemize}
\textbf{Étape 2 — Variations.} Le carbone passe de $-\mathrm{IV}$ à $+\mathrm{IV}$ : augmentation de $8$ (oxydation, \ce{CH4} est le réducteur). Chaque oxygène passe de $0$ à $-\mathrm{II}$ : diminution de $2$ par atome (réduction, \ce{O2} est l'oxydant). Le n.o. de H ne change pas.

\textbf{Étape 3 — Compensation.} Il faut que la diminution totale compense l'augmentation de $8$ : il faut donc $4$ atomes d'oxygène réduits, soit $2\,\ce{O2}$. D'où :
\[ \ce{CH4 + 2 O2 -> CO2 + 2 H2O} \]
\textbf{Étape 4 — Vérification.} C : $1 = 1$ ; H : $4 = 4$ ; O : $4 = 4$. L'équation est équilibrée. Cette réaction, totale et fortement exothermique, est bien une oxydoréduction, alors qu'aucune demi-équation électronique n'était apparente.
\end{exoresolu}

\begin{exoresolu}[Équilibrage de la réaction du dichromate sur l'éthanol]
En milieu acide, l'ion dichromate \ce{Cr2O7^2-} oxyde l'éthanol \ce{CH3CH2OH} en acide éthanoïque \ce{CH3COOH} et se réduit en ion \ce{Cr^3+} (c'est le principe de l'éthylotest). Équilibrer l'équation par les n.o.
\tcblower
\textbf{Étape 1 — n.o.} Cr : $+\mathrm{VI}$ dans \ce{Cr2O7^2-}, $+\mathrm{III}$ dans \ce{Cr^3+}. Carbone fonctionnel de l'éthanol (celui lié à O) : $-\mathrm{I}$ ; dans \ce{CH3COOH}, ce même carbone : $+\mathrm{III}$.

\textbf{Étape 2 — Variations.} Chaque Cr diminue de $3$ ; comme \ce{Cr2O7^2-} contient $2$ Cr, la diminution totale est de $6$ par ion dichromate. Le carbone fonctionnel augmente de $4$ par molécule d'éthanol.

\textbf{Étape 3 — Compensation.} Le plus petit multiple commun de $6$ et $4$ est $12$ : il faut $2\,\ce{Cr2O7^2-}$ (diminution $2 \times 6 = 12$) et $3\,\ce{CH3CH2OH}$ (augmentation $3 \times 4 = 12$).

\textbf{Étape 4 — Équilibrage de O et H} avec \ce{H+} et \ce{H2O} :
\[ \ce{2 Cr2O7^2- + 3 CH3CH2OH + 16 H+ -> 4 Cr^3+ + 3 CH3COOH + 11 H2O} \]
\textbf{Vérification des charges :} $2 \times (-2) + 16 = +12$ à gauche ; $4 \times (+3) = +12$ à droite. L'équation est correcte. On observe le virage de l'orange (dichromate) au vert (ion chrome III).
\end{exoresolu}

\begin{attention}[Pièges fréquents avec les nombres d'oxydation]
\begin{itemize}
\item \textbf{Compter tous les atomes} de l'élément : dans \ce{Cr2O7^2-}, la variation de n.o. concerne \emph{deux} atomes de chrome ; oublier ce facteur 2 conduit à des coefficients faux.
\item \textbf{Ne pas oublier les exceptions} : n.o. de l'oxygène $= -\mathrm{I}$ dans les peroxydes (\ce{H2O2}), n.o. de l'hydrogène $= -\mathrm{I}$ dans les hydrures (\ce{NaH}).
\item \textbf{Un n.o. peut être fractionnaire} en moyenne : dans \ce{Fe3O4} (magnétite de la rouille), n.o. moyen du fer $= +\frac{8}{3}$, car le solide contient en réalité un fer $+\mathrm{II}$ et deux fers $+\mathrm{III}$.
\item \textbf{Toutes les réactions ne sont pas des oxydoréductions} : dans $\ce{HCl + NaOH -> NaCl + H2O}$, aucun n.o. ne varie ; c'est une réaction acido-basique, pas redox.
\end{itemize}
\end{attention}

\begin{exercice}[Pour t'entraîner]
\begin{enumerate}
\item Détermine le nombre d'oxydation du soufre dans \ce{SO2}, \ce{SO3}, \ce{H2S} et \ce{S2O3^2-} (n.o. moyen).
\item Détermine le n.o. de l'iode dans \ce{I2} et dans \ce{I-}, puis montre que la réaction $\ce{I2 + 2 S2O3^2- -> 2 I- + S4O6^2-}$ est bien une oxydoréduction.
\item Équilibre par la méthode des n.o. la formation de la rouille : $\ce{Fe + O2 -> Fe2O3}$.
\end{enumerate}
\end{exercice}

\begin{corrige}[Exercice]
\textbf{1.} \ce{SO2} : $x + 2(-2) = 0 \Rightarrow +\mathrm{IV}$. \ce{SO3} : $+\mathrm{VI}$. \ce{H2S} : $x + 2 = 0 \Rightarrow -\mathrm{II}$. \ce{S2O3^2-} : $2x - 6 = -2 \Rightarrow x = +2$, soit $+\mathrm{II}$ en moyenne.

\textbf{2.} Dans \ce{I2} (corps simple) : n.o.$(\mathrm{I}) = 0$ ; dans \ce{I-} : $-\mathrm{I}$. L'iode diminue de $1$ par atome (réduction). Dans \ce{S2O3^2-}, le soufre vaut $+\mathrm{II}$ en moyenne ; dans \ce{S4O6^2-} : $4x - 12 = -2 \Rightarrow x = +\frac{5}{2}$, soit $+\mathrm{II{,}5}$ : le soufre augmente (oxydation). Deux éléments changent de n.o. : c'est bien une oxydoréduction.

\textbf{3.} Fe : $0 \rightarrow +\mathrm{III}$ (augmentation de $3$ par atome) ; O : $0 \rightarrow -\mathrm{II}$ (diminution de $2$ par atome). Pour compenser : $4$ Fe (augmentation totale $12$) et $6$ O, soit $3\,\ce{O2}$ (diminution totale $12$) : $\ce{4 Fe + 3 O2 -> 2 Fe2O3}$.
\end{corrige}

\begin{aretenir}[L'essentiel de la section]
\begin{itemize}
\item Le \cle{nombre d'oxydation} est la charge fictive d'un atome si toutes les liaisons étaient ioniques ; il s'écrit en chiffres romains signés ($+\mathrm{VI}$, $-\mathrm{II}$).
\item Règles clés : corps simple $\rightarrow 0$ ; ion monoatomique $\rightarrow$ sa charge ; O $\rightarrow -\mathrm{II}$ (sauf peroxydes) ; H $\rightarrow +\mathrm{I}$ (sauf hydrures) ; somme des n.o. $= 0$ (molécule) ou $=$ charge de l'ion.
\item Oxydation $=$ \textbf{augmentation} du n.o. ; réduction $=$ \textbf{diminution} du n.o. Une réaction est une oxydoréduction si au moins un élément change de n.o.
\item La variation de n.o. compte les électrons échangés : elle permet d'équilibrer les équations en égalisant électrons cédés et captés, y compris pour les réactions sans ions (combustions, voie sèche).
\end{itemize}
\end{aretenir}

Cet outil en main, nous pouvons maintenant quitter la solution aqueuse et aborder les réactions d'oxydoréduction \emph{par voie sèche} — combustions, métallurgie, corrosion — que nous synthétiserons dans la dernière section de cette leçon.

\coursec{Réactions d'oxydoréduction par voie sèche et synthèse}

Dans la section précédente, tu as découvert les \cle{nombres d'oxydation}, un outil puissant qui permet de repérer une oxydoréduction même quand le transfert d'électrons n'est pas visible sous forme d'ions en solution. Jusqu'ici, toutes nos réactions se déroulaient \emph{en solution aqueuse} : ions \ce{H+} attaquant le zinc, ions \ce{Cu^2+} réagissant avec le fer, ions \ce{MnO4-} oxydant \ce{Fe^2+}. Mais le feu qui rougit la lame d'un forgeron de Foumban, la combustion du bois dans la cuisine, la soudure des rails de chemin de fer : aucune de ces transformations ne met en jeu des ions en solution. Sont-ce pour autant des oxydoréductions ? C'est toute la question de cette dernière section — et la réponse, tu vas le voir, est un magnifique \emph{oui}.

\courssub{Notion de réaction d'oxydoréduction par voie sèche}

\textbf{Intuition d'abord.} Quand la paille de fer brûle dans le dioxygène avec une pluie d'étincelles, il n'y a ni eau, ni ions, ni demi-équation électronique évidente. Pourtant, le fer passe du métal brillant à un oxyde noir, et le dioxygène disparaît. Quelque chose de profond s'est passé entre les atomes. Les nombres d'oxydation vont nous le révéler.

\begin{definition}[Réaction d'oxydoréduction par voie sèche]
Une réaction d'\cle{oxydoréduction par voie sèche} est une réaction d'oxydoréduction qui se déroule \textbf{hors de toute solution aqueuse}, entre des réactifs \textbf{non dissous} : solides, liquides purs ou gaz. Elle nécessite souvent une \textbf{température élevée} (apport initial d'énergie thermique) ou une amorce (flamme, étincelle).
\end{definition}

Le critère d'identification ne peut plus être l'échange d'ions en solution : c'est la \textbf{variation des nombres d'oxydation} qui devient l'outil décisif.

\begin{methode}[Identifier une oxydoréduction par voie sèche]
\begin{enumerate}
\item Écrire l'équation-bilan équilibrée de la réaction.
\item Attribuer le nombre d'oxydation de \textbf{chaque élément} dans les réactifs et dans les produits (rappel : corps simple $\Rightarrow$ n.o. $= 0$ ; oxygène combiné $\Rightarrow$ $-\mathrm{II}$ sauf peroxydes ; hydrogène combiné $\Rightarrow$ $+\mathrm{I}$ sauf hydrures).
\item Comparer les n.o. avant et après : si au moins un élément voit son n.o. \textbf{augmenter} (oxydation) et un autre le voit \textbf{diminuer} (réduction), la réaction est une oxydoréduction.
\item Identifier l'oxydant (l'espèce dont l'élément voit son n.o. diminuer) et le réducteur (l'espèce dont l'élément voit son n.o. augmenter).
\end{enumerate}
\end{methode}

\courssub{Exemple 1 : la combustion du fer dans le dioxygène}

\begin{experience}[Combustion de la paille de fer dans le dioxygène]
\textbf{Dispositif et protocole.} On fixe un petit morceau de paille de fer (laine d'acier) à l'extrémité d'un fil de fer, on l'enflamme, puis on le plonge dans un flacon contenant du dioxygène pur, avec un fond de sable ou d'eau pour protéger le verre.

\textbf{Observations.} La combustion, lente dans l'air, devient \textbf{vive et éclatante} dans le dioxygène : gerbe d'étincelles, lumière blanche-jaune intense. Il se forme des particules \textbf{noires} qui retombent au fond du flacon : c'est l'oxyde magnétique de fer \ce{Fe3O4} (magnétite).

\textbf{Interprétation.} Le fer a réagi avec le dioxygène selon :
\[ \ce{3Fe + 2O2 -> Fe3O4} \]
\end{experience}

\begin{popfigure}
\begin{center}
\begin{tikzpicture}[scale=0.95, every node/.style={font=\small}]
% Flacon
\draw[thick, popdark] (-2.2,0) -- (-2.2,4.2) (-2.2,4.2) -- (-0.9,4.2) (-0.9,4.2) -- (-0.9,5.0);
\draw[thick, popdark] (2.2,0) -- (2.2,4.2) (2.2,4.2) -- (0.9,4.2) (0.9,4.2) -- (0.9,5.0);
\draw[thick, popdark] (-2.2,0) -- (2.2,0);
% Fond de sable
\fill[popgoldL] (-2.15,0) rectangle (2.15,0.45);
\draw[popgold!70!black] (-2.15,0.45) .. controls (-1,0.55) and (1,0.35) .. (2.15,0.45);
\node[popgold!60!black, font=\footnotesize] at (0,0.22) {sable (protection)};
% Dioxygène
\node[popblue, font=\footnotesize] at (1.5,3.6) {\ce{O2}};
\node[popblue, font=\footnotesize] at (-1.6,2.6) {\ce{O2}};
\node[popblue, font=\footnotesize] at (0.3,3.1) {\ce{O2}};
% Fil de fer
\draw[thick, popmuted] (0,5.2) -- (0,2.6);
% Paille de fer en combustion
\fill[poppink!60] (0,2.35) circle (0.28);
\draw[poppink, thick] (0,2.35) circle (0.28);
% Étincelles
\foreach \a in {20,60,100,140,200,250,300,340}{
  \draw[poporange, thick, -{Stealth[length=1.5mm]}] (0,2.35) -- ++(\a:0.75);
}
\foreach \a in {40,120,220,320}{
  \draw[popgold, thick] (0,2.35) -- ++(\a:1.05);
}
\node[popink, font=\footnotesize, align=center] at (0,1.35) {étincelles};
% Particules noires Fe3O4
\foreach \p in {(-1.2,0.55),(0.1,0.6),(1.1,0.52),(-0.5,0.62),(0.7,0.58)}{
  \fill[popdark] \p circle (0.07);
}
\node[popdark, font=\footnotesize] at (0,-0.45) {particules noires de \ce{Fe3O4}};
% Étiquette flacon
\node[popmuted, font=\footnotesize] at (0,4.6) {flacon de dioxygène};
\end{tikzpicture}
\end{center}
\legende{Combustion vive de la paille de fer dans le dioxygène : gerbe d'étincelles et formation de magnétite \ce{Fe3O4} noire.}
\end{popfigure}

\textbf{Analyse par les nombres d'oxydation.} Attribuons les n.o. :
\begin{itemize}
\item Dans \ce{Fe} (corps simple) : n.o.(Fe) $= 0$. Dans \ce{O2} (corps simple) : n.o.(O) $= 0$.
\item Dans \ce{Fe3O4} : l'oxygène est au n.o. $-\mathrm{II}$. La neutralité impose $3\,x + 4\times(-2) = 0$, soit $x = +\dfrac{8}{3}$. Le fer passe donc au n.o. moyen $+\dfrac{8}{3}$ (en réalité, \ce{Fe3O4} contient deux ions \ce{Fe^3+} et un ion \ce{Fe^2+} par motif, d'où cette valeur fractionnaire).
\end{itemize}

Le n.o. du fer \textbf{augmente} ($0 \to +\frac{8}{3}$) : le fer est \textbf{oxydé}, c'est le réducteur. Le n.o. de l'oxygène \textbf{diminue} ($0 \to -\mathrm{II}$) : le dioxygène est \textbf{réduit}, c'est l'oxydant. C'est bien une oxydoréduction — sans le moindre ion en solution !

\begin{attention}[Toute combustion dans le dioxygène est une oxydoréduction]
Ne cherche pas des ions pour décider si une réaction est une oxydoréduction. \textbf{Toute combustion dans le dioxygène} (métal, carbone, hydrocarbure, alcool) est une oxydoréduction : le dioxygène (n.o. $0$) est toujours réduit en oxygène combiné (n.o. $-\mathrm{II}$), et l'élément combustible est toujours oxydé. Le critère, c'est la \textbf{variation des n.o.}, pas la présence d'ions.
\end{attention}

\courssub{Exemple 2 : combustion du carbone et du méthane}

\textbf{Combustion du carbone} (charbon de bois dans la cuisine) :
\[ \ce{C + O2 -> CO2} \]
n.o. du carbone : $0 \to +\mathrm{IV}$ (oxydation) ; n.o. de l'oxygène : $0 \to -\mathrm{II}$ (réduction). Le carbone est le réducteur, \ce{O2} l'oxydant.

\textbf{Combustion complète du méthane} (gaz de cuisine, biogaz) :
\[ \ce{CH4 + 2O2 -> CO2 + 2H2O} \]
Vérifions le carbone : dans \ce{CH4}, $x + 4\times(+1) = 0$ donc n.o.(C) $= -\mathrm{IV}$ ; dans \ce{CO2}, n.o.(C) $= +\mathrm{IV}$. Le carbone est oxydé de $-\mathrm{IV}$ à $+\mathrm{IV}$ : chaque atome de carbone « perd » formellement 8 électrons, récupérés par les 4 atomes d'oxygène qui passent chacun de $0$ à $-\mathrm{II}$. Le compte est bon : $8 = 4 \times 2$.

\begin{attention}[Piège fréquent : le n.o. négatif du carbone]
Beaucoup d'élèves croient qu'un nombre d'oxydation ne peut pas être négatif. Faux ! Dans \ce{CH4}, le carbone, plus électronégatif que l'hydrogène, porte le n.o. $-\mathrm{IV}$. Une oxydation correspond à une \textbf{augmentation algébrique} du n.o. : passer de $-\mathrm{IV}$ à $+\mathrm{IV}$ est bien une augmentation.
\end{attention}

\courssub{Exemple 3 : réduction de l'oxyde de cuivre(II)}

\textbf{Par le dihydrogène.} En chauffant de l'oxyde de cuivre(II) noir dans un courant de dihydrogène, on observe l'apparition de cuivre métallique rouge et de buée (eau) sur les parois froides du tube :
\[ \ce{CuO + H2 ->[\Delta] Cu + H2O} \]
n.o. du cuivre : $+\mathrm{II} \to 0$ (réduction, \ce{CuO} est l'oxydant) ; n.o. de l'hydrogène : $0 \to +\mathrm{I}$ (oxydation, \ce{H2} est le réducteur).

\textbf{Par le carbone.} En chauffant fortement un mélange de \ce{CuO} et de poudre de charbon :
\[ \ce{2CuO + C ->[\Delta] 2Cu + CO2} \]
Le cuivre passe de $+\mathrm{II}$ à $0$ (réduction), le carbone de $0$ à $+\mathrm{IV}$ (oxydation). C'est le principe ancestral de la métallurgie du cuivre : le charbon « arrache » l'oxygène de l'oxyde métallique.

\courssub{Exemple 4 : réaction d'un métal avec le dichlore}

Le sodium brûle dans le dichlore avec une flamme jaune vive ; le magnésium s'y enflamme également :
\[ \ce{2Na + Cl2 -> 2NaCl} \qquad \ce{Mg + Cl2 -> MgCl2} \]
n.o. du sodium : $0 \to +\mathrm{I}$ (oxydation) ; n.o. du chlore : $0 \to -\mathrm{I}$ (réduction). Le dichlore joue ici le rôle d'oxydant, exactement comme le dioxygène. Remarque au passage que le transfert d'électrons est cette fois très concret : \ce{NaCl} est un solide \textbf{ionique}, constitué d'ions \ce{Na+} et \ce{Cl-} — le sodium a littéralement cédé un électron au chlore. La voie sèche et la voie ionique se rejoignent !

\begin{exemplebox}[L'aluminothermie : souder les rails par la chimie]
La réaction entre l'aluminium en poudre et l'oxyde de fer(III), amorcée par un ruban de magnésium, est \textbf{extrêmement exothermique} (température atteinte : environ \SI{2500}{\celsius}) :
\[ \ce{2Al + Fe2O3 -> 2Fe + Al2O3} \]
n.o. de l'aluminium : $0 \to +\mathrm{III}$ (oxydation, Al est le réducteur) ; n.o. du fer : $+\mathrm{III} \to 0$ (réduction, \ce{Fe2O3} est l'oxydant). La chaleur dégagée est telle que le fer produit est \textbf{liquide} : on le coule entre deux rails pour les souder. Cette réaction classique au Baccalauréat illustre parfaitement une oxydoréduction par voie sèche, sans eau ni ions.
\end{exemplebox}

\begin{savaistu}[Le haut fourneau : l'oxydoréduction à l'échelle industrielle]
La production mondiale de fer repose sur une oxydoréduction par voie sèche géante : dans le \textbf{haut fourneau}, le minerai \ce{Fe2O3} est réduit vers \SI{1500}{\celsius} par le monoxyde de carbone, lui-même formé par combustion du coke :
\[ \ce{Fe2O3 + 3CO -> 2Fe + 3CO2} \]
Le fer passe de $+\mathrm{III}$ à $0$ (réduction), le carbone de CO de $+\mathrm{II}$ à $+\mathrm{IV}$ (oxydation). Chaque année, plus d'un milliard de tonnes de fonte sont ainsi produites dans le monde — la chimie redox est littéralement le squelette de notre civilisation industrielle.
\end{savaistu}

\courssub{Exercice résolu et tableau de synthèse}

\begin{exoresolu}[Identifier et équilibrer une oxydoréduction par voie sèche]
\textbf{Énoncé.} Le tungstène, métal des filaments de lampes, est préparé par réduction de l'oxyde \ce{WO3} par le dihydrogène à chaud. (1) Écrire l'équation-bilan. (2) Montrer, à l'aide des nombres d'oxydation, qu'il s'agit d'une oxydoréduction, et identifier oxydant et réducteur.
\tcblower
\textbf{Solution.} (1) Les produits sont le tungstène métallique et l'eau :
\[ \ce{WO3 + 3H2 ->[\Delta] W + 3H2O} \]
(2) Attribuons les n.o. : dans \ce{WO3}, $x + 3\times(-2) = 0$ donc n.o.(W) $= +\mathrm{VI}$ ; dans W métal, n.o.(W) $= 0$. Dans \ce{H2}, n.o.(H) $= 0$ ; dans \ce{H2O}, n.o.(H) $= +\mathrm{I}$.
Le tungstène voit son n.o. \textbf{diminuer} de $+\mathrm{VI}$ à $0$ : il est réduit, donc \ce{WO3} est l'\textbf{oxydant}. L'hydrogène voit son n.o. \textbf{augmenter} de $0$ à $+\mathrm{I}$ : il est oxydé, donc \ce{H2} est le \textbf{réducteur}. Les variations de n.o. étant non nulles, la réaction est bien une oxydoréduction par voie sèche. Bilan électronique : W gagne formellement 6 électrons, les 6 atomes d'hydrogène en perdent chacun 1 : $6 = 6$, l'équation est cohérente.
\end{exoresolu}

\begin{center}
\begin{tabularx}{\linewidth}{|l|X|X|l|}
\hline
\rowcolor{popblueL} \textbf{Réaction} & \textbf{Élément oxydé (n.o.)} & \textbf{Élément réduit (n.o.)} & \textbf{Domaine} \\
\hline
\ce{Zn + 2H+ -> Zn^2+ + H2} & Zn : $0 \to +\mathrm{II}$ & H : $+\mathrm{I} \to 0$ & solution aqueuse \\
\hline
\ce{Fe + Cu^2+ -> Fe^2+ + Cu} & Fe : $0 \to +\mathrm{II}$ & Cu : $+\mathrm{II} \to 0$ & solution aqueuse \\
\hline
\ce{MnO4- + 5Fe^2+ + 8H+ -> Mn^2+ + 5Fe^3+ + 4H2O} & Fe : $+\mathrm{II} \to +\mathrm{III}$ & Mn : $+\mathrm{VII} \to +\mathrm{II}$ & solution aqueuse \\
\hline
\ce{3Fe + 2O2 -> Fe3O4} & Fe : $0 \to +\frac{8}{3}$ & O : $0 \to -\mathrm{II}$ & voie sèche \\
\hline
\ce{CH4 + 2O2 -> CO2 + 2H2O} & C : $-\mathrm{IV} \to +\mathrm{IV}$ & O : $0 \to -\mathrm{II}$ & voie sèche \\
\hline
\ce{2CuO + C -> 2Cu + CO2} & C : $0 \to +\mathrm{IV}$ & Cu : $+\mathrm{II} \to 0$ & voie sèche \\
\hline
\ce{2Na + Cl2 -> 2NaCl} & Na : $0 \to +\mathrm{I}$ & Cl : $0 \to -\mathrm{I}$ & voie sèche \\
\hline
\ce{2Al + Fe2O3 -> 2Fe + Al2O3} & Al : $0 \to +\mathrm{III}$ & Fe : $+\mathrm{III} \to 0$ & voie sèche \\
\hline
\end{tabularx}
\end{center}

\courssub{Bilan de la leçon et ouverture}

\begin{aretenir}[Les deux visages d'une même famille de réactions]
L'oxydoréduction peut être décrite par \textbf{deux points de vue complémentaires} :
\begin{itemize}
\item le point de vue \textbf{électronique} : l'oxydation est une \textbf{perte} d'électrons, la réduction un \textbf{gain} d'électrons ; il est parfaitement adapté aux solutions aqueuses, où les demi-équations électroniques s'écrivent naturellement ;
\item le point de vue des \textbf{nombres d'oxydation} : l'oxydation est une \textbf{augmentation} du n.o., la réduction une \textbf{diminution} du n.o. ; il est universel et s'applique aussi bien aux solutions qu'aux réactions par voie sèche.
\end{itemize}
Dans les deux cas, oxydation et réduction sont \textbf{toujours simultanées}, et l'oxydant d'un couple capte ce que le réducteur de l'autre couple cède.
\end{aretenir}

Toutes les réactions étudiées dans cette leçon sont \textbf{spontanées} : le réducteur le plus fort cède ses électrons à l'oxydant le plus fort, sans intervention extérieure. Mais cette spontanéité soulève deux questions passionnantes. Peut-on \emph{exploiter} ce transfert d'électrons spontané pour produire du courant électrique ? C'est le principe des \textbf{piles}, que nous étudierons bientôt. Et peut-on \emph{forcer} une réaction redox à se produire dans le sens non spontané, en imposant un courant ? C'est le domaine de l'\textbf{électrolyse}, au cœur de la production industrielle de l'aluminium, du chlore ou du zinc. Tu tiens maintenant tous les outils — demi-équations, couples, nombres d'oxydation — pour aborder ces deux chapitres avec confiance.

\newpage

\coursec{Travaux pratiques}

\courssub{Objectifs du TP}

À l'issue de cette séance de travaux pratiques, tu seras capable de :
\begin{itemize}
  \item réaliser expérimentalement l'action des acides chlorhydrique et sulfurique \cle{dilués} sur différents métaux (zinc, fer, cuivre, aluminium) ;
  \item identifier le \cle{gaz dégagé} (dihydrogène) par le test à la flamme et les \cle{ions métalliques} formés (\ce{Zn^2+}, \ce{Fe^2+}, \ce{Al^3+}) par précipitation avec la soude ;
  \item réaliser l'action du zinc et du fer sur les ions \ce{Cu^2+}, \ce{Zn^2+} et \ce{Fe^2+} et observer dépôts métalliques et décolorations ;
  \item écrire les \cle{demi-équations électroniques} et les \cle{équations-bilans} des réactions d'oxydoréduction observées ;
  \item conclure sur le sens spontané des réactions entre couples redox.
\end{itemize}

\courssub{Matériel et produits}

\begin{center}
\rowcolors{2}{popblueL}{white}
\begin{tabularx}{\linewidth}{|l|X|}
\hline
\rowcolor{popblue!40} \textbf{Matériel (par groupe de 4)} & \textbf{Produits} \\
\hline
8 tubes à essai + portoir & Grenaille de zinc, limaille de fer, tournures de cuivre, petits morceaux d'aluminium \\
\hline
1 pipette jaugée de \SI{5}{mL} + poire & Acide chlorhydrique dilué (\SI{1}{mol.L^{-1}}), acide sulfurique dilué (\SI{0,5}{mol.L^{-1}}) \\
\hline
1 spatule, 1 pince & Solution de soude (\SI{1}{mol.L^{-1}}) \\
\hline
1 lame de zinc, 1 clou en fer décapé & Solutions de \ce{CuSO4} (\SI{0,5}{mol.L^{-1}}), \ce{ZnSO4} (\SI{0,5}{mol.L^{-1}}), \ce{FeSO4} fraîchement préparée (\SI{0,5}{mol.L^{-1}}) \\
\hline
Bec Bunsen ou bougie, allumettes, bouchon & Papier pH, eau distillée, papier de verre (pour décaper) \\
\hline
\end{tabularx}
\end{center}

\begin{attention}[Sécurité : les gestes qui sauvent]
\begin{itemize}
  \item \textbf{Acides dilués (pictogramme « corrosif » : deux gouttes attaquant une main et une surface)} : port obligatoire de la blouse, des lunettes et des gants. En cas de contact avec la peau, rincer abondamment à l'eau froide pendant plusieurs minutes et prévenir l'enseignant.
  \item \textbf{Soude (même pictogramme corrosif)} : manipuler avec précaution, ne jamais pipetter à la bouche.
  \item \textbf{Test du dihydrogène} : le mélange \ce{H2}/air est \textbf{explosif}. Boucher le tube dès le dégagement gazeux, éloigner le tube du visage et des camarades, puis approcher la flamme. Ne jamais chauffer un tube bouché.
  \item \textbf{Orientation des tubes} : l'ouverture d'un tube à essai ne doit jamais pointer vers toi ni vers un voisin.
  \item \textbf{Versement de l'acide} : toujours verser l'acide \emph{dans} le tube contenant déjà le métal, jamais l'inverse (règle : « l'acide dans l'eau » pour toute dilution).
  \item Les solutions de sulfate de cuivre sont \textbf{toxiques pour l'environnement} : récupérer les résidus dans le bidon prévu, ne rien jeter à l'évier.
  \item \textbf{Manipulation de l'aluminium} : la couche d'alumine superficielle retarde la réaction ; ne pas forcer le chauffage pour « accélérer », le dégagement d'\ce{H2} peut devenir soudain et violent.
\end{itemize}
\end{attention}

\courssub{Schéma du dispositif}

\begin{popfigure}
\begin{center}
\begin{tikzpicture}[scale=0.9, >=Stealth]
% Tube à essai
\draw[thick, popblue] (0,0) -- (0,3.2);
\draw[thick, popblue] (1.2,0) -- (1.2,3.2);
\draw[thick, popblue] (0,0) arc (180:360:0.6);
% Liquide
\fill[popblue!25] (0.05,0.35) -- (0.05,1.6) -- (1.15,1.6) -- (1.15,0.35) arc (180:360:0.55);
\draw[popblue!60] (0.05,1.6) -- (1.15,1.6);
% Métal au fond
\fill[popmuted!60] (0.25,0.25) circle (0.09);
\fill[popmuted!60] (0.5,0.18) circle (0.09);
\fill[popmuted!60] (0.75,0.28) circle (0.09);
\fill[popmuted!60] (0.95,0.2) circle (0.09);
% Bulles
\foreach \x/\y/\r in {0.35/1.0/0.05, 0.6/1.3/0.06, 0.8/0.9/0.045, 0.5/2.0/0.07, 0.75/2.4/0.06, 0.4/2.7/0.05}
  \draw[popblue!70] (\x,\y) circle (\r);
% Bouchon
\fill[poporange!70] (0.02,3.2) rectangle (1.18,3.45);
% Flamme
\draw[thick, popgold] (2.6,2.6) -- (2.6,3.6);
\fill[poporange] (2.6,3.9) ellipse (0.12 and 0.28);
\fill[popgold!80] (2.6,3.85) ellipse (0.06 and 0.16);
% Labels
\node[align=left, font=\small] at (4.6,2.2) {Bouchon puis flamme :};
\node[align=left, font=\small, poporange] at (4.6,1.7) {détonation $\Rightarrow$ \ce{H2}};
\node[align=left, font=\small] at (0.6,-0.7) {Métal + acide dilué};
\end{tikzpicture}
\end{center}
\legende{Dispositif d'attaque d'un métal par un acide dilué et identification du dihydrogène par le test à la flamme.}
\end{popfigure}

\courssub{Protocole}

\textbf{Partie A : action des acides dilués sur les métaux}
\begin{enumerate}
  \item Numéroter 8 tubes à essai. Introduire dans chacun une pointe de spatule de métal : grenaille de zinc (tubes 1 et 2), limaille de fer (tubes 3 et 4), tournures de cuivre (tubes 5 et 6), aluminium (tubes 7 et 8).
  \item Ajouter \SI{5}{mL} d'acide chlorhydrique dilué dans les tubes 1, 3, 5, 7 et \SI{5}{mL} d'acide sulfurique dilué dans les tubes 2, 4, 6, 8.
  \item Observer pendant 2 minutes : dégagement gazeux ou non, vitesse de la réaction, échauffement éventuel du tube.
  \item \textbf{Test du gaz} : pour les tubes où un dégagement est visible, boucher rapidement le tube, attendre 30 secondes, puis approcher une flamme de l'ouverture. Noter le résultat (détonation caractéristique : le gaz est le dihydrogène \ce{H2}).
  \item \textbf{Test des ions formés} : prélever quelques gouttes de chaque solution ayant réagi et ajouter de la soude goutte à goutte. Noter la couleur du précipité : blanc pour \ce{Zn(OH)2} et \ce{Al(OH)3} (ce dernier se redissout dans un excès de soude), vert pour \ce{Fe(OH)2} (qui noircit à l'air par oxydation en \ce{Fe(OH)3}).
\end{enumerate}

\textbf{Partie B : action d'un métal sur un ion métallique}
\begin{enumerate}
  \item Verser \SI{5}{mL} de solution de \ce{CuSO4} (bleue) dans un tube et y plonger la lame de zinc. Observer pendant 5 minutes.
  \item Verser \SI{5}{mL} de solution de \ce{CuSO4} dans un second tube et y plonger le clou en fer décapé. Observer.
  \item Verser \SI{5}{mL} de solution de \ce{ZnSO4} dans un troisième tube et y plonger un clou en fer décapé. Observer : se passe-t-il quelque chose ?
  \item Verser \SI{5}{mL} de solution de \ce{FeSO4} (verdâtre) dans un quatrième tube et y ajouter de la grenaille de zinc. Observer la décoloration.
  \item Consigner toutes les observations dans un tableau et rédiger le compte rendu.
\end{enumerate}

\begin{savaistu}[Et si un produit manque au laboratoire ?]
Pas de sulfate de fer(II) ? On peut le préparer juste avant la séance en faisant réagir de la limaille de fer avec de l'acide sulfurique dilué, puis en filtrant. Pas de lame de zinc ? La grenaille de zinc convient très bien, et les boîtes de conserve découpées fournissent du fer. Le bleu de la solution de cuivre peut aussi être obtenu en dissolvant de la « pierre bleue » (sulfate de cuivre vendu en quincaillerie comme fongicide). Un bec Bunsen absent peut être remplacé par une simple bougie pour le test du dihydrogène.
\end{savaistu}

\courssub{Observations et résultats attendus}

\begin{center}
\rowcolors{2}{popgreenL}{white}
\begin{tabularx}{\linewidth}{|l|X|X|X|}
\hline
\rowcolor{popgreen!40} \textbf{Métal} & \textbf{Avec HCl dilué} & \textbf{Avec \ce{H2SO4} dilué} & \textbf{Test soude} \\
\hline
Zinc & Dégagement gazeux vif ; détonation à la flamme & Dégagement vif ; détonation & Précipité blanc (\ce{Zn(OH)2}) \\
\hline
Fer & Dégagement lent ; détonation & Dégagement lent ; détonation & Précipité vert (\ce{Fe(OH)2}) noircissant à l'air \\
\hline
Aluminium & Dégagement après un temps de latence (couche d'alumine) & Dégagement modéré après latence & Précipité blanc soluble dans l'excès de soude (\ce{Al(OH)3} amphotère) \\
\hline
Cuivre & Aucune réaction & Aucune réaction & --- \\
\hline
\end{tabularx}
\end{center}

\begin{center}
\rowcolors{2}{poporangeL}{white}
\begin{tabularx}{\linewidth}{|X|X|}
\hline
\rowcolor{poporange!40} \textbf{Expérience} & \textbf{Observation attendue} \\
\hline
Zn dans \ce{CuSO4} & Dépôt rougeâtre de cuivre sur le zinc ; la solution bleue se décolore progressivement \\
\hline
Fe (clou) dans \ce{CuSO4} & Dépôt rougeâtre de cuivre sur le clou ; décoloration de la solution \\
\hline
Fe (clou) dans \ce{ZnSO4} & Aucune transformation visible \\
\hline
Zn dans \ce{FeSO4} & Dépôt grisâtre de fer ; la teinte verdâtre de la solution disparaît \\
\hline
\end{tabularx}
\end{center}

\courssub{Exploitation}

\begin{methode}[Rappel : rédaction des demi-équations et de l'équation-bilan]
\begin{enumerate}
  \item Identifier les deux couples redox en présence (ex. \ce{Zn^2+}/\ce{Zn} et \ce{Cu^2+}/\ce{Cu}).
  \item Écrire la demi-équation d'\textbf{oxydation} (perte d'électrons) pour le réducteur.
  \item Écrire la demi-équation de \textbf{réduction} (gain d'électrons) pour l'oxydant.
  \item Multiplier les demi-équations par les coefficients nécessaires pour égaler le nombre d'électrons.
  \item Additionner les deux demi-équations et simplifier les électrons (et les espèces communes) pour obtenir l'\textbf{équation-bilan}.
\end{enumerate}
\end{methode}

\begin{exoresolu}[Questions guidées et réponses]
\textbf{Q1.} Quel gaz se dégage lors de l'action de l'acide chlorhydrique sur le zinc ? Écrire l'équation-bilan de la réaction.

\textbf{Q2.} Pourquoi le cuivre ne réagit-il pas avec les acides dilués ?

\textbf{Q3.} Écrire les demi-équations électroniques et l'équation-bilan de la réaction entre le zinc et les ions \ce{Cu^2+}.

\textbf{Q4.} Le clou de fer déposé dans la solution de \ce{ZnSO4} ne subit aucune transformation. Interpréter.

\textbf{Q5.} Classer les trois métaux Zn, Fe et Cu du plus réducteur au moins réducteur.

\textbf{Q6.} Écrire les demi-équations et l'équation-bilan pour la réaction de l'aluminium avec l'acide chlorhydrique dilué.
\tcblower
\textbf{R1.} Le gaz qui détone à la flamme est le \cle{dihydrogène} \ce{H2}. Les ions \ce{H3O+} de l'acide sont réduits (ils gagnent des électrons) et le zinc est oxydé (il perd des électrons). Demi-équations :
\begin{align*}
\text{Oxydation : } & \ce{Zn(s) -> Zn^2+(aq) + 2e-} \\
\text{Réduction : } & \ce{2H3O+(aq) + 2e- -> H2(g) + 2H2O(l)}
\end{align*}
Équation-bilan :
\[ \ce{Zn(s) + 2H3O+(aq) -> Zn^2+(aq) + H2(g) + 2H2O(l)} \]

\textbf{R2.} Le cuivre est un métal \emph{peu réducteur} (couple \ce{Cu^2+}/\ce{Cu} au potentiel standard élevé) : il ne peut pas céder d'électrons aux ions \ce{H3O+} des acides dilués non oxydants. Il n'est donc pas attaqué (contrairement à l'acide nitrique, oxydant, qui l'attaque via le couple \ce{NO3-}/\ce{NO} ou \ce{NO2}).

\textbf{R3.} Couples en présence : \ce{Zn^2+}/\ce{Zn} et \ce{Cu^2+}/\ce{Cu}. Le zinc est le réducteur, l'ion cuivre(II) est l'oxydant.
\begin{align*}
\text{Oxydation : } & \ce{Zn(s) -> Zn^2+(aq) + 2e-} \\
\text{Réduction : } & \ce{Cu^2+(aq) + 2e- -> Cu(s)}
\end{align*}
Équation-bilan (les électrons ne figurent pas) :
\[ \ce{Zn(s) + Cu^2+(aq) -> Zn^2+(aq) + Cu(s)} \]
Le dépôt rougeâtre est le cuivre métallique formé ; la décoloration traduit la disparition des ions \ce{Cu^2+} bleus.

\textbf{R4.} Aucune réaction ne se produit : le fer est un réducteur \emph{moins fort} que le zinc (couple \ce{Fe^2+}/\ce{Fe} moins réducteur que \ce{Zn^2+}/\ce{Zn}), il ne peut donc pas réduire les ions \ce{Zn^2+}. La réaction spontanée aurait lieu dans l'autre sens : le zinc réduit \ce{Fe^2+} (tube 4 de la partie B).

\textbf{R5.} D'après les expériences : le zinc réduit \ce{Fe^2+} et \ce{Cu^2+}, le fer réduit \ce{Cu^2+} mais pas \ce{Zn^2+}. D'où le classement du plus réducteur au moins réducteur (pouvoir réducteur décroissant) :
\[ \ce{Zn} > \ce{Fe} > \ce{Cu} \]

\textbf{R6.} Pour l'aluminium, le couple est \ce{Al^3+}/\ce{Al} (3 électrons). L'oxydant est toujours \ce{H3O+}/\ce{H2} (2 électrons). Il faut égaler à 6 électrons (PPCM de 2 et 3).
\begin{align*}
\text{Oxydation (x3) : } & \ce{2Al(s) -> 2Al^3+(aq) + 6e-} \\
\text{Réduction (x3) : } & \ce{6H3O+(aq) + 6e- -> 3H2(g) + 6H2O(l)}
\end{align*}
Équation-bilan :
\[ \ce{2Al(s) + 6H3O+(aq) -> 2Al^3+(aq) + 3H2(g) + 6H2O(l)} \]
Souvent simplifiée en milieu acide par : $\ce{2Al + 6H+ -> 2Al^3+ + 3H2}$.
\end{exoresolu}

\courssub{Synthèse des demi-équations électroniques de référence}

\begin{aretenir}[Demi-équations des couples redox rencontrés]
\begin{center}
\begin{tabularx}{\linewidth}{|l|l|l|}
\hline
\rowcolor{poppurple!40} \textbf{Couple redox} & \textbf{Demi-équation de réduction (gain e-)} & \textbf{Demi-équation d'oxydation (perte e-)} \\
\hline
\ce{H3O+}/\ce{H2} & \ce{2H3O+ + 2e- -> H2 + 2H2O} & \ce{H2 + 2H2O -> 2H3O+ + 2e-} \\
\hline
\ce{Cu^2+}/\ce{Cu} & \ce{Cu^2+ + 2e- -> Cu} & \ce{Cu -> Cu^2+ + 2e-} \\
\hline
\ce{Fe^2+}/\ce{Fe} & \ce{Fe^2+ + 2e- -> Fe} & \ce{Fe -> Fe^2+ + 2e-} \\
\hline
\ce{Zn^2+}/\ce{Zn} & \ce{Zn^2+ + 2e- -> Zn} & \ce{Zn -> Zn^2+ + 2e-} \\
\hline
\ce{Al^3+}/\ce{Al} & \ce{Al^3+ + 3e- -> Al} & \ce{Al -> Al^3+ + 3e-} \\
\hline
\end{tabularx}
\end{center}
\textbf{Règle du sens spontané :} Le réducteur le plus fort (en haut à gauche du tableau des couples classés par potentiels) cède ses électrons à l'oxydant le plus fort (en bas à droite). La réaction spontanée va du couple de potentiel le plus bas vers le couple de potentiel le plus haut.
\end{aretenir}

\courssub{Conclusion}

Ce TP met en évidence deux grandes familles de réactions d'\cle{oxydoréduction} en solution aqueuse. D'une part, les acides chlorhydrique et sulfurique dilués attaquent les métaux suffisamment réducteurs (zinc, fer, aluminium) avec dégagement de dihydrogène et formation d'ions métalliques, tandis que les métaux nobles (cuivre, argent, or) restent inertes face aux acides non oxydants. D'autre part, un métal peut réduire les ions d'un autre métal moins noble : le zinc réduit \ce{Cu^2+} et \ce{Fe^2+}, le fer réduit \ce{Cu^2+}, mais jamais l'inverse. Ces observations conduisent à un classement des couples redox selon le pouvoir réducteur des métaux et permettent de \emph{prévoir le sens spontané} d'une réaction : le réducteur le plus fort cède ses électrons à l'oxydant le plus fort. C'est exactement ce principe de \cle{protection cathodique} (anode sacrificielle) qui explique pourquoi on protège les coques de bateaux, les pipelines enterrés (pétrole, eau au Cameroun) et les réservoirs en y fixant des blocs de zinc ou de magnésium, qui s'oxydent « à la place » du fer.

\newpage

\coursec{Méthodes et exercices résolus}

\coursec{Action d'un acide dilué sur un métal}

\begin{definition}[Oxydation, réduction, oxydant, réducteur — Point de vue électronique]
Une \cle{oxydation} est une perte d'électrons ; une \cle{réduction} est un gain d'électrons. L'espèce chimique qui perd des électrons est le \cle{réducteur} ; celle qui en gagne est l'\cle{oxydant}. Une réaction d'\cle{oxydoréduction} est un transfert d'électrons d'un réducteur vers un oxydant.
\end{definition}

\begin{definition}[Couple redox]
Un \cle{couple redox} est un couple d'espèces chimiques (oxydant/réducteur) lié par un échange d'électrons. On le note \ce{Ox/Red} (ex. : \ce{Zn^{2+}/Zn}, \ce{H+/H2}, \ce{Cu^{2+}/Cu}). L'oxydant est la forme oxydée, le réducteur la forme réduite.
\end{definition}

\begin{experience}[Action de l'acide chlorhydrique sur le zinc, le fer et le cuivre]
\begin{itemize}
\item \textbf{Dispositif :} Trois tubes à essai contenant chacun \SI{5}{mL} d'acide chlorhydrique \SI{1}{mol.L^{-1}}. On y introduit respectivement : de la grenaille de zinc, de la limaille de fer, un fil de cuivre. Un tube renversé au-dessus de chaque tube permet de recueillir un éventuel gaz.
\item \textbf{Protocole :} Observer l'évolution pendant quelques minutes. Tester le gaz recueilli par la flamme (détonation caractéristique). Prélever quelques gouttes de solution pour tester les cations par l'hydroxyde de sodium.
\item \textbf{Observations :}
\begin{itemize}
\item \textbf{Zinc :} Dégagement gazeux rapide, solution incolore. Flamme : « aboiement ». Test \ce{NaOH} : précipité blanc soluble en excès (\ce{Zn^{2+}}).
\item \textbf{Fer :} Dégagement gazeux plus lent, solution verdâtre. Flamme : « aboiement ». Test \ce{NaOH} : précipité vert \ce{Fe(OH)2} (s'oxyde en brun à l'air).
\item \textbf{Cuivre :} Aucun dégagement gazeux, solution inchangée, métal intact.
\end{itemize}
\item \textbf{Interprétation :} Zn et Fe sont oxydés par les ions \ce{H+} (réduits en \ce{H2}). Cu est moins réducteur que \ce{H2}, pas de réaction.
\end{itemize}
\end{experience}

\begin{methode}[Prévoir et identifier les produits de l'action d'un acide sur un métal]
\begin{enumerate}
\item \textbf{Identifier l'acide.} Les acides chlorhydrique \ce{HCl} et sulfurique dilué \ce{H2SO4} fournissent des ions \ce{H3O+} (souvent notés \ce{H+} pour simplifier). Seul le couple \ce{H+/H2} intervient.
\item \textbf{Classer le métal.} Un métal est attaqué par un acide dilué s'il est \cle{plus réducteur que le dihydrogène}. Au programme : le zinc, le fer et l'aluminium réagissent ; le cuivre, l'argent et l'or \textbf{ne réagissent pas}.
\item \textbf{Écrire les deux demi-équations} :
\[ \ce{M -> M^{n+} + n e^-} \qquad \ce{2H+ + 2e^- -> H2} \]
\item \textbf{Combiner} en multipliant pour égaliser les électrons échangés, puis simplifier.
\item \textbf{Identifier les produits expérimentalement} : dégagement de dihydrogène (détonation à la flamme : \og aboiement \fg) et formation d'un ion métallique (tests caractéristiques : \ce{Fe^{2+}} vert pâle, précipité vert avec la soude ; \ce{Zn^{2+}} précipité blanc soluble dans un excès de soude ; \ce{Al^{3+}} précipité blanc gélatineux).
\end{enumerate}
\textbf{Réflexe Bac :} si l'énoncé dit \og le métal n'est pas attaqué \fg, conclus que le métal est moins réducteur que \ce{H2} (cas de \ce{Cu}, \ce{Ag}, \ce{Au}).
\end{methode}

\begin{exoresolu}[Attaque du zinc par l'acide chlorhydrique]
On introduit une masse $m = \SI{6,54}{g}$ de grenaille de zinc dans un bécher contenant $V = \SI{100}{mL}$ d'acide chlorhydrique de concentration $C = \SI{1,0}{mol.L^{-1}}$. On observe un dégagement gazeux abondant.\par
1. Écrire les demi-équations électroniques puis l'équation-bilan de la réaction.\par
2. Déterminer le réactif limitant et le volume de gaz dégagé dans les conditions où $V_m = \SI{24,0}{L.mol^{-1}}$.\par
3. Comment identifier le gaz et le cation formé ?\par
\emph{Donnée :} $M(\ce{Zn}) = \SI{65,4}{g.mol^{-1}}$.
\tcblower
\textbf{1. Demi-équations et équation-bilan.}\par
Le zinc est un métal réducteur, il cède des électrons (oxydation) ; les ions \ce{H+} les captent (réduction) :
\[ \ce{Zn -> Zn^{2+} + 2e^-} \qquad \ce{2H+ + 2e^- -> H2} \]
Les électrons échangés sont déjà au nombre de deux de part et d'autre ; on additionne membre à membre :
\[ \boxed{\ce{Zn + 2H+ -> Zn^{2+} + H2}} \]
\textbf{2. Réactif limitant et volume de gaz.}\par
Quantités initiales :
\[ n_0(\ce{Zn}) = \frac{m}{M} = \frac{6,54}{65,4} = \SI{0,100}{mol} \qquad n_0(\ce{H+}) = C \times V = 1,0 \times 0{,}100 = \SI{0,100}{mol} \]
Dressons le tableau d'avancement :
\begin{center}
\begin{tabularx}{\linewidth}{|l|X|X|X|X|}
\hline
\rowcolor{popblueL} État & \ce{Zn} & \ce{2H+} & \ce{Zn^{2+}} & \ce{H2} \\
\hline
Initial (mol) & $0{,}100$ & $0{,}100$ & $0$ & $0$ \\
\hline
En cours (mol) & $0{,}100 - x$ & $0{,}100 - 2x$ & $x$ & $x$ \\
\hline
Final (mol) & $0{,}050$ & $0$ & $0{,}050$ & $0{,}050$ \\
\hline
\end{tabularx}
\end{center}
Hypothèses : si le zinc est limitant, $x_{\max} = \SI{0,100}{mol}$ ; si \ce{H+} est limitant, $x_{\max} = \dfrac{0{,}100}{2} = \SI{0,050}{mol}$. La plus petite valeur l'emporte : $x_{\max} = \SI{0,050}{mol}$, donc \textbf{l'acide est le réactif limitant}.\par
Volume de dihydrogène :
\[ V(\ce{H2}) = x_{\max} \times V_m = 0{,}050 \times 24{,}0 = \boxed{\SI{1,2}{L}} \]
\textbf{3. Identifications.}\par
\emph{Gaz :} on approche une flamme du tube recueillant le gaz ; une légère détonation (\og aboiement \fg) prouve la présence de \ce{H2}.\par
\emph{Cation :} on prélève quelques gouttes de la solution et on ajoute de la soude : un précipité blanc, soluble dans un excès de soude, caractérise les ions \ce{Zn^{2+}}.\par
\textbf{Conclusion :} le zinc, plus réducteur que \ce{H2}, est oxydé en \ce{Zn^{2+}} avec dégagement de \SI{1,2}{L} de dihydrogène ; l'acide est entièrement consommé.
\end{exoresolu}

\begin{exercice}[Attaque de l'aluminium par l'acide sulfurique dilué]
On verse \SI{50,0}{mL} d'une solution d'acide sulfurique \SI{0,50}{mol.L^{-1}} sur \SI{1,35}{g} d'aluminium en poudre.
1. Écrire l'équation-bilan de la réaction.
2. Déterminer le réactif limitant et le volume de dihydrogène formé ($V_m = \SI{24,0}{L.mol^{-1}}$).
3. Préciser l'aspect final du mélange.
\emph{Données :} $M(\ce{Al}) = \SI{27,0}{g.mol^{-1}}$ ; \ce{Al^{3+}} donne un précipité blanc gélatineux \ce{Al(OH)3} insoluble en excès de soude.
\end{exercice}

\coursec{Oxydation, réduction et couples redox : le point de vue électronique}

\begin{propriete}[Sens spontané d'une réaction d'oxydoréduction en solution aqueuse]
Une réaction d'oxydoréduction entre deux couples \ce{Ox1/Red1} et \ce{Ox2/Red2} est spontanée dans le sens où le \textbf{réducteur le plus fort} (celui qui donne le plus facilement ses électrons) réduit l'\textbf{oxydant du couple plus faible}.
\par Concrètement : un métal réduit les ions d'un métal \emph{moins réducteur} que lui. L'échelle des métaux (ordre croissant de caractère réducteur) : \ce{Au} < \ce{Ag} < \ce{Cu} < \ce{H2} < \ce{Fe} < \ce{Zn} < \ce{Al}.
\end{propriete}

\begin{aretenir}[Demi-équations et équation-bilan]
\begin{itemize}
\item Demi-équation d'oxydation : \ce{Red -> Ox + n e^-} (électrons à droite).
\item Demi-équation de réduction : \ce{Ox + n e^- -> Red} (électrons à gauche).
\item Équation-bilan : on multiplie les demi-équations pour avoir le même nombre d'électrons, on additionne, \textbf{les électrons disparaissent}.
\item Vérification systématique : conservation des éléments \textbf{et} des charges.
\end{itemize}
\end{aretenir}

\coursec{Action d'un ion métallique sur un métal}

\begin{experience}[Action du zinc et du fer sur des ions métalliques]
\begin{itemize}
\item \textbf{Dispositif :} Quatre tubes à essai : 1) Zn + \ce{Cu^{2+}} (bleu), 2) Zn + \ce{Fe^{2+}} (vert pâle), 3) Fe + \ce{Cu^{2+}} (bleu), 4) Fe + \ce{Zn^{2+}} (incolore). Concentrations \SI{0,1}{mol.L^{-1}}.
\item \textbf{Observations :}
\begin{itemize}
\item Tube 1 : Dépôt rouge sur le zinc, solution décolorée. \textbf{Réaction.}
\item Tube 2 : Dépôt grisâtre sur le zinc, solution décolorée. \textbf{Réaction.}
\item Tube 3 : Dépôt rouge sur le fer, solution pâlit (devient verte). \textbf{Réaction.}
\item Tube 4 : Aucun changement visible. \textbf{Pas de réaction.}
\end{itemize}
\item \textbf{Interprétation :} Zn réduit \ce{Cu^{2+}} et \ce{Fe^{2+}} (Zn plus réducteur que Cu et Fe). Fe réduit \ce{Cu^{2+}} (Fe plus réducteur que Cu). Fe ne réduit pas \ce{Zn^{2+}} (Fe moins réducteur que Zn).
\end{itemize}
\end{experience}

\begin{methode}[Réaction entre un ion métallique et un métal]
\begin{enumerate}
\item \textbf{Repérer les deux couples redox} mis en jeu : \ce{M1^{n+}/M1} et \ce{M2^{p+}/M2}.
\item \textbf{Prévoir le sens de la réaction} : un métal \cle{réduit} les ions d'un métal \emph{moins réducteur} que lui. Expérimentalement : le zinc et le fer réduisent \ce{Cu^{2+}} ; le zinc réduit \ce{Fe^{2+}} ; mais le fer ne réduit pas \ce{Zn^{2+}}.
\item \textbf{Écrire les demi-équations} dans le bon sens : le métal s'oxyde (perte d'électrons), l'ion se réduit (gain d'électrons).
\item \textbf{Multiplier} chaque demi-équation pour égaliser le nombre d'électrons, puis additionner : les électrons ne doivent \textbf{jamais apparaître} dans l'équation-bilan.
\item \textbf{Vérifier} la conservation des éléments et des charges, puis interpréter les observations (dépôt métallique, décoloration de la solution, attaque du métal).
\end{enumerate}
\end{methode}

\begin{exoresolu}[Le fer plongé dans une solution de sulfate de cuivre]
Un clou en fer de masse $m = \SI{5,60}{g}$ est plongé dans $V = \SI{200}{mL}$ d'une solution bleue de sulfate de cuivre(II) de concentration $C = \SI{0,50}{mol.L^{-1}}$. Après quelques minutes, le clou se recouvre d'un dépôt rougeâtre et la solution pâlit.\par
1. Interpréter les observations et écrire l'équation-bilan de la réaction.\par
2. Calculer la masse de cuivre déposée lorsque la réaction est totale.\par
\emph{Données :} $M(\ce{Fe}) = \SI{56,0}{g.mol^{-1}}$ ; $M(\ce{Cu}) = \SI{63,5}{g.mol^{-1}}$.
\tcblower
\textbf{1. Interprétation et équation-bilan.}\par
Le dépôt rougeâtre est du cuivre métallique : les ions \ce{Cu^{2+}} (responsables de la couleur bleue) ont été réduits. Le fer, plus réducteur que le cuivre, a été oxydé en ions \ce{Fe^{2+}} (la solution vire au vert pâle).\par
Demi-équations :
\[ \ce{Fe -> Fe^{2+} + 2e^-} \qquad \ce{Cu^{2+} + 2e^- -> Cu} \]
Équation-bilan (deux électrons de chaque côté, on additionne directement) :
\[ \boxed{\ce{Fe + Cu^{2+} -> Fe^{2+} + Cu}} \]
Vérification : éléments conservés (1 Fe, 1 Cu) ; charges : $+2$ de chaque côté.\par
\textbf{2. Masse de cuivre déposée.}\par
Quantités initiales :
\[ n_0(\ce{Fe}) = \frac{5{,}60}{56{,}0} = \SI{0,100}{mol} \qquad n_0(\ce{Cu^{2+}}) = 0{,}50 \times 0{,}200 = \SI{0,100}{mol} \]
L'équation montre des proportions 1 pour 1 : les réactifs sont en \textbf{proportions stœchiométriques}, $x_{\max} = \SI{0,100}{mol}$. Il se forme donc $n(\ce{Cu}) = \SI{0,100}{mol}$ :
\[ m(\ce{Cu}) = n \times M = 0{,}100 \times 63{,}5 = \boxed{\SI{6,35}{g}} \]
\textbf{Conclusion :} le fer réduit les ions \ce{Cu^{2+}} ; il se dépose \SI{6,35}{g} de cuivre sur le clou, qui est lui-même entièrement consommé. C'est le principe de la cémentation, utilisée en métallurgie.
\end{exoresolu}

\begin{exercice}[Zinc et ions fer(II)]
On plonge un morceau de zinc de masse $m = \SI{3,27}{g}$ dans $V = \SI{100}{mL}$ d'une solution de sulfate de fer(II) \SI{0,20}{mol.L^{-1}}.
1. Prévoir si une réaction a lieu. Si oui, écrire l'équation-bilan.
2. Calculer la masse de fer métallique déposée sur le zinc.
\emph{Données :} $M(\ce{Zn}) = \SI{65,4}{g.mol^{-1}}$ ; $M(\ce{Fe}) = \SI{56,0}{g.mol^{-1}}$.
\end{exercice}

\coursec{Autres réactions d'oxydoréduction en solution aqueuse}

\begin{methode}[Équilibrer une demi-équation électronique en milieu acide]
Pour un couple du type \ce{MnO4^{-}/Mn^{2+}} ou \ce{Cr2O7^{2-}/Cr^{3+}}, procéder dans l'ordre \textbf{E-O-H-e} :
\begin{enumerate}
\item \textbf{Élément} autre que O et H : équilibrer l'atome central (Mn, Cr...).
\item \textbf{Oxygène} : ajouter des molécules d'eau \ce{H2O} du côté qui manque d'oxygène.
\item \textbf{Hydrogène} : ajouter des ions \ce{H+} du côté qui manque d'hydrogène.
\item \textbf{Électrons} : ajouter des \ce{e-} pour équilibrer les charges.
\item Vérifier : éléments \emph{et} charges doivent être conservés.
\end{enumerate}
Pour obtenir l'équation-bilan : multiplier chaque demi-équation par le plus petit coefficient qui égalise les électrons, additionner, puis \textbf{simplifier} les \ce{H+} et \ce{H2O} apparaissant des deux côtés.
\end{methode}

\begin{exoresolu}[Oxydation des ions fer(II) par le permanganate en milieu acide]
On verse une solution violette de permanganate de potassium dans une solution acidifiée contenant des ions \ce{Fe^{2+}} : la solution se décolore. Couples mis en jeu : \ce{MnO4^{-}/Mn^{2+}} et \ce{Fe^{3+}/Fe^{2+}}.\par
1. Écrire les deux demi-équations électroniques.\par
2. En déduire l'équation-bilan de la réaction.\par
3. On dose $V_1 = \SI{20,0}{mL}$ d'une solution de \ce{Fe^{2+}} par une solution de \ce{KMnO4} à $C_2 = \SI{0,020}{mol.L^{-1}}$ ; l'équivalence est atteinte pour $V_{2E} = \SI{15,0}{mL}$. Calculer la concentration $C_1$ de la solution de \ce{Fe^{2+}}.
\tcblower
\textbf{1. Demi-équations.}\par
\emph{Couple \ce{MnO4^{-}/Mn^{2+}}} (méthode E-O-H-e) :
\begin{itemize}
\item Mn : déjà équilibré (1 de chaque côté) ;
\item O : 4 oxygènes à gauche, on ajoute 4 \ce{H2O} à droite : \ce{MnO4^{-} -> Mn^{2+} + 4H2O} ;
\item H : 8 hydrogènes à droite, on ajoute 8 \ce{H+} à gauche : \ce{MnO4^{-} + 8H+ -> Mn^{2+} + 4H2O} ;
\item Charges : à gauche $-1 + 8 = +7$, à droite $+2$ : on ajoute 5 électrons à gauche :
\[ \ce{MnO4^{-} + 8H+ + 5e^- -> Mn^{2+} + 4H2O} \]
\end{itemize}
\emph{Couple \ce{Fe^{3+}/Fe^{2+}}} (les ions \ce{Fe^{2+}} sont oxydés) :
\[ \ce{Fe^{2+} -> Fe^{3+} + e^-} \]
\textbf{2. Équation-bilan.}\par
Le permanganate consomme 5 électrons, le fer(II) n'en cède qu'un : on multiplie la seconde demi-équation par 5 :
\[ \ce{MnO4^{-} + 8H+ + 5e^- -> Mn^{2+} + 4H2O} \]
\[ \ce{5Fe^{2+} -> 5Fe^{3+} + 5e^-} \]
On additionne et les 5 électrons s'éliminent :
\[ \boxed{\ce{MnO4^{-} + 5Fe^{2+} + 8H+ -> Mn^{2+} + 5Fe^{3+} + 4H2O}} \]
Vérification des charges : à gauche $-1 + 10 + 8 = +17$ ; à droite $+2 + 15 = +17$. Correct.\par
La décoloration s'explique : \ce{MnO4^{-}} violet est transformé en \ce{Mn^{2+}} quasi incolore.\par
\textbf{3. Concentration de la solution de fer(II).}\par
À l'équivalence, les réactifs sont introduits dans les proportions stœchiométriques :
\[ \frac{n(\ce{Fe^{2+}})}{5} = \frac{n(\ce{MnO4^{-}})}{1} \quad \Longrightarrow \quad C_1 V_1 = 5\, C_2 V_{2E} \]
\[ C_1 = \frac{5 \times 0{,}020 \times 15{,}0}{20{,}0} = \boxed{\SI{0,075}{mol.L^{-1}}} \]
\textbf{Conclusion :} la solution de fer(II) a pour concentration \SI{7,5e-2}{mol.L^{-1}}. C'est le principe du dosage manganimétrique, très classique au Bac TI.
\end{exoresolu}

\begin{methode}[Éthanol/dichromate et diiode/thiosulfate : les deux classiques]
\textbf{A. Réaction entre l'éthanol et les ions dichromate \ce{Cr2O7^{2-}}} (milieu acide) :
\begin{enumerate}
\item Couples : \ce{Cr2O7^{2-}/Cr^{3+}} (orange $\to$ vert) et \ce{CH3COOH/CH3CH2OH}.
\item Demi-équation du dichromate (méthode E-O-H-e) : \ce{Cr2O7^{2-} + 14H+ + 6e^- -> 2Cr^{3+} + 7H2O}.
\item Demi-équation de l'éthanol : \ce{CH3CH2OH + H2O -> CH3COOH + 4H+ + 4e^-}.
\item Multiplier par 2 et 3 pour échanger 12 électrons, additionner, simplifier \ce{H+} et \ce{H2O}.
\end{enumerate}
\textbf{B. Réaction entre le diiode et les ions thiosulfate \ce{S2O3^{2-}}} :
\begin{enumerate}
\item Couples : \ce{I2/I-} et \ce{S4O6^{2-}/S2O3^{2-}}.
\item Demi-équations : \ce{I2 + 2e^- -> 2I-} et \ce{2S2O3^{2-} -> S4O6^{2-} + 2e^-}.
\item Équation-bilan : \ce{I2 + 2S2O3^{2-} -> 2I- + S4O6^{2-}} (décoloration du diiode brun).
\end{enumerate}
\textbf{Réflexe :} apprends par cœur les demi-équations des couples \og classiques \fg du programme ; au Bac, on gagne un temps précieux.
\end{methode}

\begin{exoresolu}[L'alcootest : oxydation de l'éthanol par le dichromate]
Les anciens alcootests contenaient des cristaux orangés de dichromate de potassium acidifiés. En présence d'éthanol dans l'air expiré, ils viraient au vert.\par
1. Écrire les demi-équations des couples \ce{Cr2O7^{2-}/Cr^{3+}} et \ce{CH3COOH/CH3CH2OH}.\par
2. Établir l'équation-bilan de la réaction.\par
3. Expliquer le changement de couleur observé.
\tcblower
\textbf{1. Demi-équations.}\par
\emph{Dichromate} (méthode E-O-H-e) : 2 Cr à droite ; 7 \ce{H2O} à droite pour l'oxygène ; 14 \ce{H+} à gauche pour l'hydrogène ; charges : $-2 + 14 = +12$ à gauche, $+6$ à droite, donc 6 électrons à gauche :
\[ \ce{Cr2O7^{2-} + 14H+ + 6e^- -> 2Cr^{3+} + 7H2O} \]
\emph{Éthanol} (oxydation en acide éthanoïque) : on équilibre O avec \ce{H2O}, H avec \ce{H+}, puis les charges avec \ce{e-} :
\[ \ce{CH3CH2OH + H2O -> CH3COOH + 4H+ + 4e^-} \]
\textbf{2. Équation-bilan.}\par
Le plus petit commun multiple de 6 et 4 est 12 : on multiplie la première par 2, la seconde par 3 :
\[ \ce{2Cr2O7^{2-} + 28H+ + 12e^- -> 4Cr^{3+} + 14H2O} \]
\[ \ce{3CH3CH2OH + 3H2O -> 3CH3COOH + 12H+ + 12e^-} \]
On additionne, puis on simplifie : $28\,\ce{H+} - 12\,\ce{H+} = 16\,\ce{H+}$ à gauche ; $14\,\ce{H2O} - 3\,\ce{H2O} = 11\,\ce{H2O}$ à droite :
\[ \boxed{\ce{2Cr2O7^{2-} + 3CH3CH2OH + 16H+ -> 4Cr^{3+} + 3CH3COOH + 11H2O}} \]
\textbf{3. Changement de couleur.}\par
Les ions dichromate \ce{Cr2O7^{2-}} sont \textbf{oranges} ; ils sont réduits en ions chrome(III) \ce{Cr^{3+}}, \textbf{verts}. L'éthanol de l'air expiré est oxydé en acide éthanoïque. Le virage de l'orange au vert trahit donc la présence d'alcool : c'est une application directe de l'oxydoréduction à la sécurité routière.
\end{exoresolu}

\begin{exercice}[Dosage iodométrique de l'eau de Javel]
L'eau de Javel contient des ions hypochlorite \ce{ClO-}. En milieu acide, ils oxydent les ions iodure \ce{I-} en diiode \ce{I2} : \ce{ClO- + 2I- + 2H+ -> Cl- + I2 + H2O}. Le diiode formé est dosé par des ions thiosulfate \ce{S2O3^{2-}}.
1. Rappeler l'équation-bilan de la réaction \ce{I2}/\ce{S2O3^{2-}}.
2. On a introduit \SI{10,0}{mL} d'eau de Javel dans un excès d'iodure de potassium acidifié. Le diiode libéré a nécessité \SI{12,5}{mL} de solution de thiosulfate \SI{0,100}{mol.L^{-1}} pour être dosé. Calculer la concentration molaire des ions \ce{ClO-} dans l'eau de Javel.
\end{exercice}

\coursec{Généralisation : les nombres d'oxydation}

\begin{definition}[Nombre d'oxydation (n.o.)]
Le \cle{nombre d'oxydation} d'un atome dans une espèce chimique est la charge qu'il porterait si toutes les liaisons étaient purement ioniques, les électrons étant attribués à l'atome le plus électronégatif. Il s'écrit en chiffres romains (ex. : $+\mathrm{II}$, $-\mathrm{III}$, $+\mathrm{VII}$).
\end{definition}

\begin{aretenir}[Règles de calcul des nombres d'oxydation]
\begin{enumerate}
\item Corps simple (élément pur) : n.o. = 0 (ex. : \ce{Fe}, \ce{O2}, \ce{I2}, \ce{S8}).
\item Ion monoatomique : n.o. = charge de l'ion (ex. : \ce{Fe^{2+}} : $+\mathrm{II}$ ; \ce{Cl-} : $-\mathrm{I}$).
\item Dans les composés : n.o.(O) = $-\mathrm{II}$ (sauf peroxydes $-\mathrm{I}$, superoxydes $-\frac{1}{2}$, \ce{OF2} $+\mathrm{II}$) ; n.o.(H) = $+\mathrm{I}$ (sauf hydrures métalliques $-\mathrm{I}$) ; n.o.(F) = $-\mathrm{I}$ (toujours).
\item \textbf{Règle de la somme :} la somme des n.o. est nulle pour une molécule neutre, égale à la charge de l'ion pour un ion polyatomique.
\end{enumerate}
\end{aretenir}

\begin{propriete}[Oxydation, réduction, oxydant, réducteur — Point de vue des nombres d'oxydation]
\begin{itemize}
\item \textbf{Oxydation :} \emph{augmentation} du nombre d'oxydation d'un élément.
\item \textbf{Réduction :} \emph{diminution} du nombre d'oxydation d'un élément.
\item \textbf{Oxydant :} espèce dont un élément voit son n.o. \emph{diminuer}.
\item \textbf{Réducteur :} espèce dont un élément voit son n.o. \emph{augmenter}.
\item \textbf{Équilibrage par les n.o. :} le gain total de n.o. (réduction) = perte totale de n.o. (oxydation). On en déduit les coefficients stœchiométriques des oxydant/réducteur, puis on équilibre les autres éléments (O, H, charges).
\end{itemize}
\end{propriete}

\begin{methode}[Déterminer un nombre d'oxydation (n.o.)]
\begin{enumerate}
\item Appliquer les règles de base (corps simple, ion monoatomique, O, H, F).
\item Poser $x$ le n.o. cherché pour l'élément variable.
\item Écrire l'équation de la somme des n.o. (0 pour molécule, charge pour ion).
\item Résoudre pour $x$. Le résultat peut être fractionnaire (n.o. \emph{moyen}).
\item \textbf{Exploiter :} variation du n.o. $\uparrow$ = oxydation ; $\downarrow$ = réduction.
\end{enumerate}
\end{methode}

\begin{exoresolu}[Nombres d'oxydation et équilibrage de la réaction diiode/thiosulfate]
On considère la réaction utilisée dans les dosages iodométriques :
\[ \ce{I2 + S2O3^{2-} -> I- + S4O6^{2-}} \]
1. Déterminer le nombre d'oxydation de l'iode dans \ce{I2} et \ce{I-}, puis celui du soufre dans \ce{S2O3^{2-}} et \ce{S4O6^{2-}}.\par
2. Identifier l'oxydant et le réducteur.\par
3. Équilibrer l'équation-bilan en utilisant les variations de n.o.
\tcblower
\textbf{1. Nombres d'oxydation.}\par
\emph{Iode :} dans \ce{I2} (corps simple), n.o.(I) $= 0$ ; dans \ce{I-} (ion monoatomique), n.o.(I) $= -\mathrm{I}$.\par
\emph{Soufre dans \ce{S2O3^{2-}} :} soit $x$ le n.o. moyen du soufre :
\[ 2x + 3 \times (-2) = -2 \quad \Longrightarrow \quad 2x = +4 \quad \Longrightarrow \quad x = +\mathrm{II} \]
\emph{Soufre dans \ce{S4O6^{2-}} :}
\[ 4x + 6 \times (-2) = -2 \quad \Longrightarrow \quad 4x = +10 \quad \Longrightarrow \quad x = +\tfrac{5}{2} \; (+\mathrm{II{,}5}) \]
Il s'agit d'un n.o. \emph{moyen} : les quatre atomes de soufre n'ont pas tous le même état d'oxydation.\par
\textbf{2. Oxydant et réducteur.}\par
Le n.o. de l'iode passe de $0$ à $-\mathrm{I}$ : il \emph{diminue}, donc \ce{I2} est \textbf{réduit} : c'est l'\textbf{oxydant}. Le n.o. du soufre passe de $+\mathrm{II}$ à $+\mathrm{II{,}5}$ : il \emph{augmente}, donc \ce{S2O3^{2-}} est \textbf{oxydé} : c'est le \textbf{réducteur}.\par
\textbf{3. Équilibrage par les n.o.}\par
Chaque atome d'iode gagne 1 degré d'oxydation ; une molécule \ce{I2} gagne donc $2 \times 1 = 2$ unités. Chaque ion \ce{S2O3^{2-}} contient 2 soufres dont le n.o. augmente de $0{,}5$ chacun, soit une perte totale de $2 \times 0{,}5 = 1$ unité par ion thiosulfate.\par
Pour compenser le gain de 2 unités par \ce{I2}, il faut donc 2 ions thiosulfate :
\[ \boxed{\ce{I2 + 2S2O3^{2-} -> 2I- + S4O6^{2-}}} \]
Vérification : iode (2/2), soufre (4/4), oxygène (6/6) ; charges : $2 \times (-2) = -4$ à gauche, $2 \times (-1) + (-2) = -4$ à droite.\par
\textbf{Conclusion :} la méthode des n.o. redonne exactement les coefficients obtenus par les demi-équations électroniques ; elle est particulièrement efficace quand les demi-équations sont délicates à écrire.
\end{exoresolu}

\begin{exercice}[Équilibrage par les nombres d'oxydation : Permanganate / Peroxyde d'hydrogène]
En milieu acide, les ions permanganate \ce{MnO4^{-}} oxydent le peroxyde d'hydrogène \ce{H2O2} en dioxygène \ce{O2}, en se réduisant en \ce{Mn^{2+}}.
1. Déterminer les n.o. du Mn dans \ce{MnO4^{-}} et \ce{Mn^{2+}}, de l'O dans \ce{H2O2} et \ce{O2}.
2. Identifier l'oxydant et le réducteur.
3. Équilibrer l'équation-bilan par la méthode des nombres d'oxydation.
\end{exercice}

\coursec{Réactions d'oxydoréduction par voie sèche et synthèse}

\begin{savaistu}[La rouille : une oxydoréduction lente au cœur de notre quotidien]
La \cle{rouille} (\ce{Fe2O3 . n H2O}, poudre friable brune) se forme quand le fer est au contact de l'air humide. C'est une réaction d'oxydoréduction par voie sèche (puisque l'eau est liquide mais non solvant majoritaire au début) ou humide selon les conditions.
\par \textbf{Mécanisme simplifié :}
\begin{itemize}
\item \textbf{Anode (oxydation) :} \ce{Fe -> Fe^{2+} + 2e^-} (le fer métal s'oxyde).
\item \textbf{Cathode (réduction) :} \ce{O2 + 2H2O + 4e^- -> 4OH-} (l'oxygène de l'air se réduit en présence d'eau).
\item \textbf{Global :} \ce{2Fe + O2 + 2H2O -> 2Fe^{2+} + 4OH-} puis précipitation/hydrolyse : \ce{2Fe^{2+} + 4OH- -> 2Fe(OH)2} (vert) $\xrightarrow{\ce{O2}}$ \ce{Fe2O3 . n H2O} (rouille brune).
\end{itemize}
\par \textbf{Conséquences au Cameroun :} corrosion des toitures en tôle, des armatures du béton (ponts, bâtiments), des canalisations. Protection : peinture, galvanisation (couche de Zn, plus réducteur que Fe, s'oxyde à sa place), traitement antirouille.
\end{savaistu}

\begin{savaistu}[Métallurgie extractive et CIMENCAM]
La production du fer (haut fourneau) et du ciment (CIMENCAM, Figuil) repose sur des oxydoréductions par voie sèche à haute température.
\par \textbf{Haut fourneau (fer) :} \ce{Fe2O3 + 3CO -> 2Fe + 3CO2} (le monoxyde de carbone, réducteur, réduit l'hématite). Le coke (\ce{C}) fournit la chaleur et le CO : \ce{C + O2 -> CO2} puis \ce{CO2 + C -> 2CO}.
\par \textbf{Four à ciment :} Décarbonatation du calcaire \ce{CaCO3 -> CaO + CO2} (pas d'oxydoréduction), puis combinaison \ce{CaO + SiO2 -> silicates} (clinker). La combustion du charbon/pétrole (\ce{C + O2 -> CO2}) fournit l'énergie (\SI{1450}{\celsius}).
\end{savaistu}

\begin{methode}[Reconnaître une oxydoréduction par voie sèche]
\begin{enumerate}
\item Une réaction \cle{par voie sèche} se déroule \textbf{sans solvant} (solides chauffés, gaz, combustion) : il n'y a ni ions \ce{H+} ni eau pour écrire des demi-équations électroniques classiques.
\item \textbf{Calculer les n.o.} des éléments avant et après la réaction.
\item Si au moins un n.o. \emph{augmente} (oxydation) pendant qu'un autre \emph{diminue} (réduction), c'est une réaction d'oxydoréduction.
\item Cas typiques au programme : combustion d'un métal dans le dioxygène (\ce{2Mg + O2 -> 2MgO}), réduction d'un oxyde métallique par le carbone ou le dihydrogène (\ce{CuO + H2 -> Cu + H2O}), réaction d'un halogène sur un métal (\ce{2Fe + 3Cl2 -> 2FeCl3}).
\item \textbf{Attention :} toute réaction sans variation de n.o. (acide-base, précipitation) n'est \emph{pas} une oxydoréduction, même à chaud.
\end{enumerate}
\end{methode}

\begin{exoresolu}[Réduction de l'oxyde de cuivre(II) par le dihydrogène]
On chauffe de l'oxyde de cuivre(II) \ce{CuO} noir dans un courant de dihydrogène : il apparaît un dépôt rouge de cuivre et de la buée sur les parois froides du tube.\par
1. Écrire l'équation-bilan de la réaction.\par
2. À l'aide des nombres d'oxydation, montrer qu'il s'agit d'une réaction d'oxydoréduction et préciser l'oxydant et le réducteur.\par
3. Calculer la masse de cuivre obtenue à partir de $m = \SI{7,95}{g}$ d'oxyde de cuivre(II) supposé pur.\par
\emph{Données :} $M(\ce{CuO}) = \SI{79,5}{g.mol^{-1}}$ ; $M(\ce{Cu}) = \SI{63,5}{g.mol^{-1}}$.
\tcblower
\textbf{1. Équation-bilan.}\par
La buée est de l'eau, le dépôt rouge est du cuivre :
\[ \boxed{\ce{CuO + H2 -> Cu + H2O}} \]
\textbf{2. Analyse par les nombres d'oxydation.}\par
\begin{center}
\begin{tabularx}{\linewidth}{|l|X|X|}
\hline
\rowcolor{popblueL} Élément & Avant la réaction & Après la réaction \\
\hline
Cuivre & $+\mathrm{II}$ dans \ce{CuO} & $0$ dans \ce{Cu} \\
\hline
Hydrogène & $0$ dans \ce{H2} & $+\mathrm{I}$ dans \ce{H2O} \\
\hline
Oxygène & $-\mathrm{II}$ & $-\mathrm{II}$ (inchangé) \\
\hline
\end{tabularx}
\end{center}
Le n.o. du cuivre \emph{diminue} de $+\mathrm{II}$ à $0$ : \ce{CuO} est \textbf{réduit}, c'est l'\textbf{oxydant}. Le n.o. de l'hydrogène \emph{augmente} de $0$ à $+\mathrm{I}$ : \ce{H2} est \textbf{oxydé}, c'est le \textbf{réducteur}. Il y a bien transfert d'électrons sans passage par une solution : c'est une \textbf{oxydoréduction par voie sèche}.\par
\textbf{3. Masse de cuivre obtenue.}\par
Quantité d'oxyde de cuivre :
\[ n(\ce{CuO}) = \frac{m}{M} = \frac{7{,}95}{79{,}5} = \SI{0,100}{mol} \]
D'après l'équation, $n(\ce{Cu}) = n(\ce{CuO}) = \SI{0,100}{mol}$ (proportions 1 pour 1), d'où :
\[ m(\ce{Cu}) = 0{,}100 \times 63{,}5 = \boxed{\SI{6,35}{g}} \]
\textbf{Conclusion :} le dihydrogène réduit \SI{7,95}{g} d'oxyde de cuivre(II) en \SI{6,35}{g} de cuivre métallique. Ce type de réaction est à la base de la métallurgie extractive (réduction des minerais).
\end{exoresolu}

\begin{exercice}[Combustion du magnésium et réduction de l'oxyde de fer(III)]
1. On brûle \SI{2,43}{g} de magnésium dans le dioxygène. Écrire l'équation-bilan, montrer que c'est une oxydoréduction par voie sèche (n.o.), et calculer la masse d'oxyde de magnésium formée.
2. On réduit \SI{15,96}{g} d'oxyde de fer(III) \ce{Fe2O3} par du monoxyde de carbone \ce{CO} (produisant \ce{CO2}). Écrire l'équation-bilan, identifier l'oxydant et le réducteur par les n.o., calculer le volume de \ce{CO2} formé ($V_m = \SI{24,0}{L.mol^{-1}}$).
\emph{Données :} $M(\ce{Mg}) = \SI{24,3}{g.mol^{-1}}$ ; $M(\ce{O}) = \SI{16,0}{g.mol^{-1}}$ ; $M(\ce{Fe2O3}) = \SI{159,6}{g.mol^{-1}}$.
\end{exercice}

\begin{attention}[Les 5 erreurs qui coûtent des points au Bac]
\begin{enumerate}
\item \textbf{Faire apparaître les électrons dans l'équation-bilan.} Les \ce{e-} n'existent que dans les demi-équations ; dans le bilan, ils doivent s'être éliminés après multiplication et addition. Un bilan avec des \ce{e-} est systématiquement sanctionné.
\item \textbf{Oublier de simplifier \ce{H+} et \ce{H2O}.} Quand \ce{H+} ou \ce{H2O} apparaissent des deux côtés après l'addition des demi-équations (cas éthanol/dichromate), il faut soustraire : écrire \ce{2Cr2O7^{2-} + 3CH3CH2OH + 28H+ -> ... + 14H2O} sans simplifier est faux.
\item \textbf{Inverser oxydant et réducteur.} Réflexe : le \emph{réducteur} \emph{donne} des électrons, donc il est \emph{oxydé} (son n.o. augmente). L'oxydant les reçoit et est réduit. Une inversion entraîne la fausseté de tout l'exercice.
\item \textbf{Négliger les proportions stœchiométriques dans les dosages.} À l'équivalence, $n(\ce{Fe^{2+}}) = 5\,n(\ce{MnO4^{-}})$ et non $n(\ce{Fe^{2+}}) = n(\ce{MnO4^{-}})$ : relis toujours les coefficients de l'équation-bilan avant d'écrire la relation d'équivalence.
\item \textbf{Mal calculer les nombres d'oxydation des ions polyatomiques.} La somme des n.o. est égale à la \emph{charge de l'ion}, pas à zéro : pour \ce{S2O3^{2-}}, $2x - 6 = -2$ et non $2x - 6 = 0$. Et n'oublie pas que le n.o. peut être fractionnaire (soufre : $+\mathrm{II{,}5}$ dans \ce{S4O6^{2-}}).
\end{enumerate}
\end{attention}

\begin{corrige}[Exercice : Attaque de l'aluminium par l'acide sulfurique dilué]
1. \ce{2Al + 6H+ -> 2Al^{3+} + 3H2} (ou \ce{2Al + 3H2SO4 -> Al2(SO4)3 + 3H2}).
2. $n_0(\ce{Al}) = 1,35/27,0 = 0,0500$ mol. $n_0(\ce{H+}) = 2 \times 0,50 \times 0,0500 = 0,0500$ mol. Stœchiométrie : 2 Al pour 6 H+ $\Rightarrow$ 1 Al pour 3 H+. $x_{\max}(\ce{Al}) = 0,0500$ ; $x_{\max}(\ce{H+}) = 0,0500/3 = 0,0167$ mol. \textbf{H+ limitant}. $n(\ce{H2}) = x_{\max} = 0,0167$ mol. $V(\ce{H2}) = 0,0167 \times 24,0 = \SI{0,40}{L}$.
3. Solution incolore (\ce{Al^{3+}}), pas de métal visible (Al consommé), gaz \ce{H2} parti. Test \ce{NaOH} : précipité blanc gélatineux \ce{Al(OH)3} insoluble en excès.
\end{corrige}

\begin{corrige}[Exercice : Zinc et ions fer(II)]
1. Zn plus réducteur que Fe $\Rightarrow$ réaction. \ce{Zn + Fe^{2+} -> Zn^{2+} + Fe}. Dépôt gris (Fe) sur Zn, solution décolorée.
2. $n_0(\ce{Zn}) = 3,27/65,4 = 0,0500$ mol. $n_0(\ce{Fe^{2+}}) = 0,20 \times 0,100 = 0,0200$ mol. Proportions 1:1 $\Rightarrow$ \ce{Fe^{2+}} limitant. $n(\ce{Fe}) = 0,0200$ mol. $m(\ce{Fe}) = 0,0200 \times 56,0 = \SI{1,12}{g}$.
\end{corrige}

\begin{corrige}[Exercice : Dosage iodométrique de l'eau de Javel]
1. \ce{I2 + 2S2O3^{2-} -> 2I- + S4O6^{2-}}.
2. $n(\ce{S2O3^{2-}}) = 0,100 \times 0,0125 = 1,25 \times 10^{-3}$ mol. $n(\ce{I2}) = n(\ce{S2O3^{2-}})/2 = 6,25 \times 10^{-4}$ mol. Équation \ce{ClO-}/\ce{I2} : 1 \ce{ClO-} $\to$ 1 \ce{I2}. $n(\ce{ClO-}) = n(\ce{I2}) = 6,25 \times 10^{-4}$ mol. $C(\ce{ClO-}) = 6,25 \times 10^{-4} / 0,0100 = \SI{0,0625}{mol.L^{-1}}$.
\end{corrige}

\begin{corrige}[Exercice : Équilibrage par n.o. Permanganate / Peroxyde]
1. Mn : $+\mathrm{VII}$ dans \ce{MnO4^{-}} $\to$ $+\mathrm{II}$ dans \ce{Mn^{2+}} (gain 5). O : $-\mathrm{I}$ dans \ce{H2O2} $\to$ $0$ dans \ce{O2} (perte 1 par atome O, soit 2 par \ce{H2O2}).
2. Mn diminue : \ce{MnO4^{-}} oxydant. O augmente : \ce{H2O2} réducteur.
3. Gain 5 = Perte 2 $\times$ coeff. PPCM 10. 2 \ce{MnO4^{-}} (gain 10) et 5 \ce{H2O2} (perte 10).
   \ce{2MnO4^{-} + 5H2O2 + 6H+ -> 2Mn^{2+} + 5O2 + 8H2O}. (Vérif charges : -2 + 6 = +4 à g ; +4 à d. OK).
\end{corrige}

\begin{corrige}[Exercice : Combustion Mg et Réduction Fe2O3]
1. \ce{2Mg + O2 -> 2MgO}. Mg : $0 \to +\mathrm{II}$ (oxydation). O : $0 \to -\mathrm{II}$ (réduction). Oxydoréduction voie sèche. $n(\ce{Mg}) = 2,43/24,3 = 0,100$ mol. $n(\ce{MgO}) = 0,100$ mol. $M(\ce{MgO}) = 40,3$. $m = 0,100 \times 40,3 = \SI{4,03}{g}$.
2. \ce{Fe2O3 + 3CO -> 2Fe + 3CO2}. Fe : $+\mathrm{III} \to 0$ (réduction, oxydant \ce{Fe2O3}). C : $+\mathrm{II} \to +\mathrm{IV}$ (oxydation, réducteur \ce{CO}). $n(\ce{Fe2O3}) = 15,96/159,6 = 0,100$ mol. $n(\ce{CO2}) = 3 \times 0,100 = 0,300$ mol. $V(\ce{CO2}) = 0,300 \times 24,0 = \SI{7,20}{L}$.
\end{corrige}

\newpage

\coursec{Exercices}

\courssub{Je vérifie mes connaissances}

\begin{exercice}[QCM : les bons réflexes]
Pour chaque question, une seule réponse est exacte. Entoure-la et justifie brièvement ton choix.
\begin{enumerate}
\item Au cours d'une oxydation, une espèce chimique :
\begin{enumerate}
\item gagne des électrons ;
\item perd des électrons ;
\item gagne des protons \ce{H+} ;
\item voit son nombre d'oxydation diminuer.
\end{enumerate}
\item Dans le couple \ce{Cu^2+}/\ce{Cu}, le réducteur est :
\begin{enumerate}
\item \ce{Cu^2+} ; \item \ce{Cu} ; \item \ce{H+} ; \item l'électron.
\end{enumerate}
\item Une lame de cuivre plongée dans une solution d'acide chlorhydrique dilué :
\begin{enumerate}
\item est attaquée avec dégagement de \ce{H2} ;
\item est attaquée avec dégagement de \ce{Cl2} ;
\item n'est pas attaquée ;
\item se recouvre de rouille.
\end{enumerate}
\item Le nombre d'oxydation du chlore dans l'ion \ce{ClO3-} est :
\begin{enumerate}
\item $-1$ ; \item $+3$ ; \item $+5$ ; \item $+7$.
\end{enumerate}
\end{enumerate}
\end{exercice}

\begin{exercice}[Vrai ou faux ?]
Indique si chacune des affirmations suivantes est vraie ou fausse, puis justifie ta réponse en une ou deux phrases.
\begin{enumerate}
\item Toute réaction d'oxydoréduction met en jeu un transfert d'électrons.
\item L'ion \ce{MnO4-} en milieu acide est un réducteur.
\item Dans la réaction \ce{Zn + Cu^2+ -> Zn^2+ + Cu}, le zinc est oxydé.
\item Une réaction par voie sèche ne peut pas être une oxydoréduction car il n'y a pas de solution aqueuse.
\item Le nombre d'oxydation d'un élément dans un corps simple est toujours nul.
\end{enumerate}
\end{exercice}

\begin{exercice}[Texte à trous]
Complète le texte suivant avec les mots ou expressions de la liste : \emph{oxydation ; réducteur ; électrons ; augmente ; réduction ; oxydant ; diminue ; couple redox}.

Lorsqu'une lame de fer est plongée dans une solution contenant des ions \ce{Cu^2+}, le fer perd des \ldots\ldots\ldots : il subit une \ldots\ldots\ldots On dit que le fer est le \ldots\ldots\ldots de la réaction. Les ions \ce{Cu^2+} captent ces électrons : ils subissent une \ldots\ldots\ldots ; \ce{Cu^2+} joue le rôle d'\ldots\ldots\ldots Les deux espèces \ce{Fe^2+} et \ce{Fe} forment un \ldots\ldots\ldots Du point de vue des nombres d'oxydation, celui du fer \ldots\ldots\ldots au cours de la réaction, tandis que celui du cuivre \ldots\ldots\ldots
\end{exercice}

\begin{exercice}[Définitions essentielles]
\begin{enumerate}
\item Définis, du point de vue électronique : oxydant, réducteur, oxydation, réduction.
\item Définis, du point de vue des nombres d'oxydation : oxydation et réduction.
\item Qu'appelle-t-on couple oxydant/réducteur ? Donne deux exemples rencontrés dans cette leçon.
\end{enumerate}
\end{exercice}

\courssub{Je m'entraîne}

\begin{exercice}[Demi-équations et équations-bilans]
On considère les couples redox suivants : \ce{Fe^3+}/\ce{Fe^2+} ; \ce{I2}/\ce{I-} ; \ce{Zn^2+}/\ce{Zn} ; \ce{Al^3+}/\ce{Al}.
\begin{enumerate}
\item Écris la demi-équation électronique de chaque couple (sens de la réduction).
\item On mélange une solution de chlorure de fer(III) avec du diiode en présence d'ions iodure : précise, en justifiant par la force relative des couples, quelle réaction peut se produire entre \ce{Fe^3+} et \ce{I-}, puis écris son équation-bilan.
\item Une lame d'aluminium est plongée dans une solution contenant des ions \ce{Zn^2+}. Écris l'équation-bilan de la réaction observée en précisant l'oxydant et le réducteur.
\end{enumerate}
\end{exercice}

\begin{exercice}[Nombres d'oxydation du soufre]
\begin{enumerate}
\item Détermine le nombre d'oxydation du soufre dans chacune des espèces suivantes : \ce{H2S}, \ce{SO2}, \ce{H2SO4}, \ce{S2O3^2-}, \ce{S4O6^2-}. (Pour les ions polythionates, donne le nombre d'oxydation moyen).
\item On considère la réaction utilisée lors du dosage du diiode par le thiosulfate :
\[ \ce{I2 + 2 S2O3^2- -> 2 I- + S4O6^2-} \]
À l'aide des nombres d'oxydation, montre qu'il s'agit bien d'une oxydoréduction et identifie l'oxydant et le réducteur.
\end{enumerate}
\end{exercice}

\begin{exercice}[Oxydoréduction ou pas ?]
On considère les réactions suivantes :
\begin{enumerate}
\item \ce{CH4 + 2 O2 -> CO2 + 2 H2O} (combustion du méthane dans le gaz de cuisine) ;
\item \ce{NaCl(s) -> Na+(aq) + Cl-(aq)} (dissolution du sel de cuisine) ;
\item \ce{CuO + H2 -> Cu + H2O} (réduction de l'oxyde de cuivre par le dihydrogène) ;
\item \ce{Ag+ + Cl- -> AgCl(s)} (précipitation du chlorure d'argent).
\end{enumerate}
Pour chacune d'elles, calcule les nombres d'oxydation des éléments concernés avant et après la réaction, puis conclus : s'agit-il d'une réaction d'oxydoréduction ? Justifie.
\end{exercice}

\begin{exercice}[L'éthylotest]
Lors d'un contrôle d'alcoolémie, l'éthanol \ce{C2H5OH} contenu dans l'air expiré réduit les ions dichromate \ce{Cr2O7^2-} (orange) en ions chrome(III) \ce{Cr^3+} (verts) ; l'éthanol est oxydé en acide éthanoïque \ce{CH3COOH}. Les couples mis en jeu sont \ce{Cr2O7^2-}/\ce{Cr^3+} et \ce{CH3COOH}/\ce{C2H5OH}.
\begin{enumerate}
\item Écris les deux demi-équations électroniques en milieu acide.
\item Établis l'équation-bilan de la réaction.
\item Explique pourquoi le changement de couleur observé dans le tube de l'éthylotest prouve la présence d'éthanol dans l'air expiré.
\end{enumerate}
\end{exercice}

\courssub{Je me prépare au Bac}

\begin{exercice}[Type Bac : le zinc et le sulfate de cuivre]
Au laboratoire du lycée, on plonge une lame de zinc de masse $m = \SI{10,0}{g}$ dans un volume $V = \SI{200}{mL}$ d'une solution de sulfate de cuivre(II) de concentration $C = \SI{0,50}{mol.L^{-1}}$. On observe un dépôt rougeâtre sur la lame et la décoloration progressive de la solution.

\textbf{Données :} $M(\ce{Zn}) = \SI{65,4}{g.mol^{-1}}$ ; $M(\ce{Cu}) = \SI{63,5}{g.mol^{-1}}$.

\begin{enumerate}
\item Précise les couples redox mis en jeu et écris les demi-équations électroniques correspondantes (sens de la réduction).
\item Écris l'équation-bilan de la réaction. Identifie l'oxydant et le réducteur.
\item Calcule les quantités de matière initiales de zinc et d'ions \ce{Cu^2+}. À l'aide d'un tableau d'avancement, détermine le réactif limitant.
\item Calcule la masse de cuivre déposée sur la lame et la masse de zinc restant à la fin de la réaction.
\item Calcule la concentration finale en ions \ce{Zn^2+} dans la solution (on suppose le volume inchangé).
\item Cite une application industrielle de ce type de réaction dans la protection des métaux.
\end{enumerate}
\end{exercice}

\begin{exercice}[Type Bac : dosage des ions fer(II) par le permanganate]
Une eau de source de la région de l'Ouest est soupçonnée de contenir des ions fer(II). Pour vérifier sa qualité, on dose un volume $V_1 = \SI{20,0}{mL}$ d'une solution contenant des ions \ce{Fe^2+}, acidifiée à l'acide sulfurique, par une solution violette de permanganate de potassium de concentration $C_2 = \SI{0,020}{mol.L^{-1}}$. Le changement de teinte persistant (rose pâle) est obtenu pour un volume versé $V_{2E} = \SI{16,0}{mL}$.

\textbf{Données :} couples redox : \ce{MnO4-}/\ce{Mn^2+} et \ce{Fe^3+}/\ce{Fe^2+} ; $M(\ce{Fe}) = \SI{55,8}{g.mol^{-1}}$.

\begin{enumerate}
\item Fais le schéma légendé du dispositif de dosage et indique comment on repère l'équivalence.
\item Écris les demi-équations électroniques des deux couples en milieu acide.
\item Établis l'équation-bilan de la réaction de dosage.
\item Définis l'équivalence et écris la relation entre les quantités de matière à l'équivalence.
\item Calcule la concentration $C_1$ en ions \ce{Fe^2+} de la solution, puis la concentration massique en fer (en \si{g.L^{-1}}).
\item La norme admise pour l'eau de boisson est de \SI{0,20}{g.L^{-1}} de fer. Cette eau est-elle conforme ? Justifie.
\end{enumerate}
\end{exercice}

\begin{exercice}[Type Bac : la tôle galvanisée et la rouille]
Dans les quartiers de Douala, les tôles de toiture sont en fer galvanisé, c'est-à-dire recouvertes d'une fine couche de zinc. On constate qu'une tôle galvanisée, même rayée, résiste beaucoup mieux à la rouille qu'une tôle de fer nue exposée à l'air humide.

\textbf{Données :} couples redox : \ce{Zn^2+}/\ce{Zn}, \ce{Fe^2+}/\ce{Fe}, \ce{O2}/\ce{H2O}. Le zinc est un réducteur plus fort que le fer. Nombres d'oxydation : dans \ce{Fe2O3}, le fer est au degré $+3$.

\begin{enumerate}
\item La rouille est essentiellement de l'oxyde de fer(III) hydraté. Calcule le nombre d'oxydation du fer dans \ce{Fe} métallique et dans \ce{Fe2O3}. La formation de la rouille est-elle une oxydation ou une réduction du fer ? Justifie.
\item Écris la demi-équation électronique d'oxydation du fer en ions \ce{Fe^2+}, puis celle de réduction du dioxygène en milieu humide (couple \ce{O2}/\ce{H2O}, milieu neutre/basique). Établis l'équation-bilan de la corrosion du fer (formation d'hydroxyde de fer(II)).
\item Une rayure met le fer à nu, en contact direct avec le zinc. En comparant le pouvoir réducteur des deux métaux, explique pourquoi c'est le zinc qui est oxydé préférentiellement au fer. Écris l'équation-bilan de la réaction entre le zinc et les ions \ce{Fe^2+} éventuellement formés.
\item Explique pourquoi on qualifie le zinc d'\emph{anode sacrificielle} et pourquoi la tôle galvanisée rayée reste protégée, contrairement à la tôle de fer nue.
\item Propose une autre méthode simple de protection des objets en fer contre la rouille, utilisée au quotidien au Cameroun, et explique son principe.
\end{enumerate}
\end{exercice}

\newpage

\coursec{Corrigés des exercices}

\begin{corrige}[Exercice 1 — QCM : les bons réflexes]
\begin{enumerate}
\item \textbf{Réponse b : perd des électrons.} Par définition, l'\cle{oxydation} est une perte d'électrons. La réponse a décrit une réduction, la réponse c une réaction acido-basique, et la réponse d est fausse car au cours d'une oxydation le nombre d'oxydation \emph{augmente}.
\item \textbf{Réponse b : \ce{Cu}.} Dans un couple Ox/Red, le \cle{réducteur} est l'espèce capable de céder des électrons : c'est le métal \ce{Cu} (\ce{Cu -> Cu^2+ + 2e-}). \ce{Cu^2+} est l'oxydant du couple.
\item \textbf{Réponse c : n'est pas attaquée.} Le cuivre est un réducteur trop faible pour réduire les ions \ce{H+} : il fait partie, avec l'argent et l'or, des métaux non attaqués par les acides chlorhydrique et sulfurique dilués. Seuls les métaux comme \ce{Zn}, \ce{Fe}, \ce{Al} dégagent \ce{H2}.
\item \textbf{Réponse c : $+5$.} Dans \ce{ClO3-}, avec n.o.(O) $= -2$ : $x + 3\times(-2) = -1$, d'où $x = +5$.
\end{enumerate}
\textbf{Point de vigilance :} ne pas confondre la charge de l'ion ($-1$) avec le nombre d'oxydation du chlore ($+5$).
\end{corrige}

\begin{corrige}[Exercice 2 — Vrai ou faux ?]
\begin{enumerate}
\item \textbf{Vrai.} Une oxydoréduction est par définition une réaction au cours de laquelle il y a transfert d'électrons du réducteur vers l'oxydant (directement ou via variation des nombres d'oxydation).
\item \textbf{Faux.} L'ion \ce{MnO4-} en milieu acide est un \emph{oxydant} puissant : il capte des électrons pour être réduit en \ce{Mn^2+} (couple \ce{MnO4-}/\ce{Mn^2+}).
\item \textbf{Vrai.} Le zinc passe de \ce{Zn} (n.o. $= 0$) à \ce{Zn^2+} (n.o. $= +2$) : il perd deux électrons, il est donc oxydé (c'est le réducteur de la réaction).
\item \textbf{Faux.} La combustion du méthane \ce{CH4 + 2O2 -> CO2 + 2H2O} ou la réduction de \ce{CuO} par \ce{H2} sont des oxydoréductions par voie sèche : le critère est la variation des nombres d'oxydation, pas la présence d'une solution.
\item \textbf{Vrai.} Dans un corps simple (\ce{Fe}, \ce{O2}, \ce{Cu}, \ce{S8}\ldots), le nombre d'oxydation de l'élément est par convention nul.
\end{enumerate}
\end{corrige}

\begin{corrige}[Exercice 3 — Texte à trous]
Lorsqu'une lame de fer est plongée dans une solution contenant des ions \ce{Cu^2+}, le fer perd des \textbf{électrons} : il subit une \textbf{oxydation}. On dit que le fer est le \textbf{réducteur} de la réaction. Les ions \ce{Cu^2+} captent ces électrons : ils subissent une \textbf{réduction} ; \ce{Cu^2+} joue le rôle d'\textbf{oxydant}. Les deux espèces \ce{Fe^2+} et \ce{Fe} forment un \textbf{couple redox}. Du point de vue des nombres d'oxydation, celui du fer \textbf{augmente} au cours de la réaction (de $0$ à $+2$), tandis que celui du cuivre \textbf{diminue} (de $+2$ à $0$).
\end{corrige}

\begin{corrige}[Exercice 4 — Définitions essentielles]
\begin{enumerate}
\item \textbf{Point de vue électronique :}
\begin{itemize}
\item \cle{Oxydant} : espèce chimique capable de \emph{capter} un ou plusieurs électrons.
\item \cle{Réducteur} : espèce chimique capable de \emph{céder} un ou plusieurs électrons.
\item \cle{Oxydation} : perte d'électrons par une espèce chimique.
\item \cle{Réduction} : gain d'électrons par une espèce chimique.
\end{itemize}
\item \textbf{Point de vue des nombres d'oxydation :} une oxydation est une transformation au cours de laquelle le nombre d'oxydation d'un élément \emph{augmente} ; une réduction est une transformation au cours de laquelle le nombre d'oxydation d'un élément \emph{diminue}.
\item Un \cle{couple oxydant/réducteur} (couple redox) est l'ensemble de deux espèces conjuguées, noté Ox/Red, reliées par une demi-équation électronique : $\text{Ox} + n\,e^- \ce{<=>} \text{Red}$. Exemples : \ce{Cu^2+}/\ce{Cu}, \ce{Fe^3+}/\ce{Fe^2+}, \ce{MnO4-}/\ce{Mn^2+}, \ce{I2}/\ce{I-}.
\end{enumerate}
\end{corrige}

\begin{corrige}[Exercice 5 — Demi-équations et équations-bilans]
\begin{enumerate}
\item Demi-équations (sens de la réduction) :
\begin{align*}
\ce{Fe^3+ + e- &<=> Fe^2+} & \ce{I2 + 2e- &<=> 2 I-}\\
\ce{Zn^2+ + 2e- &<=> Zn} & \ce{Al^3+ + 3e- &<=> Al}
\end{align*}
\item \ce{Fe^3+} est un oxydant plus fort que \ce{I2} ; \ce{I-} est un réducteur plus fort que \ce{Fe^2+}. La réaction spontanée se fait donc entre l'oxydant du couple le plus « haut » et le réducteur de l'autre couple :
\[ \ce{Fe^3+ + e- -> Fe^2+} \qquad \ce{2 I- -> I2 + 2e-} \]
En multipliant la première par 2 pour égaliser les électrons échangés :
\[ \ce{2 Fe^3+ + 2 I- -> 2 Fe^2+ + I2} \]
\item L'aluminium est un réducteur plus fort que le zinc : \ce{Al} est oxydé et \ce{Zn^2+} est réduit.
\[ \ce{Al -> Al^3+ + 3e-} \;(\times 2) \qquad \ce{Zn^2+ + 2e- -> Zn} \;(\times 3) \]
\[ \ce{2 Al + 3 Zn^2+ -> 2 Al^3+ + 3 Zn} \]
Oxydant : \ce{Zn^2+} ; réducteur : \ce{Al}.
\end{enumerate}
\textbf{Point de vigilance :} vérifier toujours la conservation des éléments \emph{et} de la charge dans l'équation-bilan finale.
\end{corrige}

\begin{corrige}[Exercice 6 — Nombres d'oxydation du soufre]
\begin{enumerate}
\item Avec n.o.(H) $= +1$ et n.o.(O) $= -2$ :
\begin{itemize}
\item \ce{H2S} : $2(+1) + x = 0 \Rightarrow x = -2$ ;
\item \ce{SO2} : $x + 2(-2) = 0 \Rightarrow x = +4$ ;
\item \ce{H2SO4} : $2(+1) + x + 4(-2) = 0 \Rightarrow x = +6$ ;
\item \ce{S2O3^2-} : $2x + 3(-2) = -2 \Rightarrow x = +2$ (moyenne) ;
\item \ce{S4O6^2-} : $4x + 6(-2) = -2 \Rightarrow x = +2{,}5$ (moyenne).
\end{itemize}
\item Dans \ce{I2}, n.o.(I) $= 0$ ; dans \ce{I-}, n.o.(I) $= -1$ : l'iode \emph{diminue} son nombre d'oxydation, donc \ce{I2} est \emph{réduit} : c'est l'\textbf{oxydant}. Le soufre passe de $+2$ (dans \ce{S2O3^2-}) à $+2{,}5$ (dans \ce{S4O6^2-}) : son nombre d'oxydation \emph{augmente}, donc \ce{S2O3^2-} est \emph{oxydé} : c'est le \textbf{réducteur}. Il y a bien variation des nombres d'oxydation : c'est une \cle{oxydoréduction}.
\end{enumerate}
\textbf{Point de vigilance :} un nombre d'oxydation \emph{moyen} peut être fractionnaire ($+2{,}5$) : il s'agit d'une moyenne sur des atomes de soufre qui n'ont pas tous le même environnement.
\end{corrige}

\begin{corrige}[Exercice 7 — Oxydoréduction ou pas ?]
\begin{enumerate}
\item \ce{CH4 + 2 O2 -> CO2 + 2 H2O} : dans \ce{CH4}, n.o.(C) $= -4$ ; dans \ce{CO2}, n.o.(C) $= +4$ (augmente). Dans \ce{O2}, n.o.(O) $= 0$ ; dans \ce{CO2} et \ce{H2O}, n.o.(O) $= -2$ (diminue). \textbf{Oxydoréduction} (combustion par voie sèche) : \ce{CH4} est le réducteur, \ce{O2} l'oxydant.
\item \ce{NaCl -> Na+ + Cl-} : n.o.(Na) $= +1$ avant et après ; n.o.(Cl) $= -1$ avant et après. Aucune variation : \textbf{ce n'est pas une oxydoréduction} (simple dissolution).
\item \ce{CuO + H2 -> Cu + H2O} : Cu passe de $+2$ à $0$ (réduction), H passe de $0$ à $+1$ (oxydation). \textbf{Oxydoréduction} par voie sèche : \ce{H2} est le réducteur, \ce{CuO} l'oxydant.
\item \ce{Ag+ + Cl- -> AgCl} : n.o.(Ag) $= +1$ et n.o.(Cl) $= -1$ avant et après. Aucune variation : \textbf{ce n'est pas une oxydoréduction} (réaction de précipitation).
\end{enumerate}
\textbf{Méthode :} pour conclure, il suffit de repérer \emph{au moins un élément} dont le nombre d'oxydation varie entre réactifs et produits.
\end{corrige}

\begin{corrige}[Exercice 8 — L'éthylotest]
\begin{enumerate}
\item Demi-équations en milieu acide :
\[ \ce{Cr2O7^2- + 14 H+ + 6 e- <=> 2 Cr^3+ + 7 H2O} \]
\[ \ce{C2H5OH + H2O <=> CH3COOH + 4 H+ + 4 e-} \]
\item On multiplie la première par 2 (12 $e^-$ captés) et la seconde par 3 (12 $e^-$ cédés), puis on additionne et simplifie les \ce{H+} et \ce{H2O} :
\[ \ce{2 Cr2O7^2- + 3 C2H5OH + 16 H+ -> 4 Cr^3+ + 3 CH3COOH + 11 H2O} \]
Vérification des charges : $2\times(-2) + 16 = +12$ à gauche ; $4\times(+3) = +12$ à droite. L'équation est correcte.
\item Les ions dichromate \ce{Cr2O7^2-} sont orange et les ions chrome(III) \ce{Cr^3+} sont verts. Si l'air expiré contient de l'éthanol, celui-ci réduit les ions dichromate : la réaction ci-dessus se produit et le réactif du tube vire \emph{de l'orange au vert}. Ce changement de couleur net et visible à l'œil nu constitue donc un test qualitatif de la présence d'éthanol.
\end{enumerate}
\textbf{Point de vigilance :} ne pas oublier d'équilibrer d'abord les éléments autres que O et H, puis l'oxygène avec \ce{H2O}, puis l'hydrogène avec \ce{H+}, et enfin la charge avec les électrons.
\end{corrige}

\begin{corrige}[Exercice 9 — Type Bac : le zinc et le sulfate de cuivre]
\begin{enumerate}
\item Couples mis en jeu : \ce{Cu^2+}/\ce{Cu} et \ce{Zn^2+}/\ce{Zn}. Demi-équations (sens de la réduction) :
\[ \ce{Cu^2+ + 2e- <=> Cu} \qquad \ce{Zn^2+ + 2e- <=> Zn} \]
\item Le zinc, réducteur plus fort que le cuivre, est oxydé ; les ions \ce{Cu^2+} sont réduits :
\[ \ce{Zn + Cu^2+ -> Zn^2+ + Cu} \]
Oxydant : \ce{Cu^2+} ; réducteur : \ce{Zn}. Le dépôt rougeâtre est le cuivre métallique formé ; la décoloration traduit la disparition des ions \ce{Cu^2+} bleus.
\item Quantités initiales :
\[ n_i(\ce{Zn}) = \frac{m}{M} = \frac{10{,}0}{65{,}4} \approx \SI{0,153}{mol} \qquad n_i(\ce{Cu^2+}) = C \times V = 0{,}50 \times 0{,}200 = \SI{0,100}{mol} \]
\begin{center}
\begin{tabular}{|l|c|c|c|c|}
\hline
\rowcolor{popblueL} Équation & \ce{Zn} & \ce{+ Cu^2+} & \ce{-> Zn^2+} & \ce{+ Cu} \\
\hline
État initial (mol) & 0{,}153 & 0{,}100 & 0 & 0 \\
\hline
État final (mol) & $0{,}153 - x_{max}$ & $0{,}100 - x_{max}$ & $x_{max}$ & $x_{max}$ \\
\hline
\end{tabular}
\end{center}
Hypothèses : $x_{max} = 0{,}153$ mol (Zn limitant) ou $x_{max} = 0{,}100$ mol (\ce{Cu^2+} limitant). La plus petite valeur l'emporte : $x_{max} = \SI{0,100}{mol}$, le \textbf{réactif limitant est \ce{Cu^2+}}.
\item Masse de cuivre déposée :
\[ m(\ce{Cu}) = x_{max} \times M(\ce{Cu}) = 0{,}100 \times 63{,}5 = \SI{6,35}{g} \]
Zinc consommé : $0{,}100 \times 65{,}4 = \SI{6,54}{g}$ ; masse de zinc restant :
\[ m_f(\ce{Zn}) = 10{,}0 - 6{,}54 = \SI{3,46}{g} \]
\item Il s'est formé $n(\ce{Zn^2+}) = x_{max} = \SI{0,100}{mol}$ dans $V = \SI{0,200}{L}$ :
\[ [\ce{Zn^2+}]_f = \frac{0{,}100}{0{,}200} = \SI{0,50}{mol.L^{-1}} \]
\item Application industrielle : la \textbf{galvanisation} (revêtement du fer par du zinc pour protéger tôles et ferronnerie de la corrosion) ou la métallisation par cémentation.
\end{enumerate}
\textbf{Barème indicatif (sur 10) :} couples et demi-équations 2 pts ; équation-bilan + oxydant/réducteur 2 pts ; quantités initiales et réactif limitant 2 pts ; masses 2 pts ; concentration finale 1 pt ; application 1 pt.\\
\textbf{Points de vigilance :} convertir le volume en litres ; ne pas oublier que la lame de zinc \emph{perd} de la masse tandis que le cuivre se dépose ; vérifier l'homogénéité des unités à chaque étape.
\end{corrige}

\begin{corrige}[Exercice 10 — Type Bac : dosage des ions fer(II) par le permanganate]
\begin{enumerate}
\item Dispositif : la solution de permanganate de potassium (violette) est placée dans la \textbf{burette} graduée ; le volume $V_1 = \SI{20,0}{mL}$ de solution de \ce{Fe^2+} acidifiée est placé dans l'\textbf{erlenmeyer} (ou bécher) avec un barreau aimanté et un agitateur magnétique.
\begin{popfigure}
\begin{center}
\begin{tikzpicture}[scale=0.9]
% burette
\draw[thick] (0,4.2) -- (0,0.8);
\draw[thick] (0.5,4.2) -- (0.5,0.8);
\draw[thick] (0.15,0.8) -- (0.15,0.3) -- (0.35,0.3) -- (0.35,0.8);
\fill[popink!40] (0.02,4.1) rectangle (0.48,2.6);
\draw[thick] (0.25,0.3) -- (0.25,0.05);
\foreach \y in {3.0,3.4,3.8} \draw (0,\y) -- (0.12,\y);
% erlenmeyer
\draw[thick] (-0.9,-2.2) -- (-0.35,-0.9) -- (-0.35,-0.55);
\draw[thick] (1.4,-2.2) -- (0.85,-0.9) -- (0.85,-0.55);
\draw[thick] (-0.35,-0.55) -- (0.85,-0.55);
\draw[thick] (-0.9,-2.2) -- (1.4,-2.2);
\fill[popgold!30] (-0.78,-2.15) -- (1.28,-2.15) -- (1.05,-1.35) -- (-0.55,-1.35) -- cycle;
\fill[popdark] (0.25,-2.0) ellipse (0.22 and 0.07);
% goutte
\fill[popink] (0.25,-0.35) circle (0.05);
\node[right] at (0.7,3.4) {\small Burette : \ce{KMnO4} ($C_2$)};
\node[right] at (1.5,-1.6) {\small Erlenmeyer : \ce{Fe^2+} acidifié ($V_1$)};
\node[below] at (0.25,-2.35) {\small Barreau aimanté};
\end{tikzpicture}
\end{center}
\legende{Schéma du dosage des ions \ce{Fe^2+} par le permanganate de potassium.}
\end{popfigure}
L'équivalence est repérée par la \textbf{première teinte rose pâle persistante} : tant que des ions \ce{Fe^2+} restent, le permanganate versé est décoloré ; dès qu'il est en léger excès, sa couleur violette-rosée persiste (auto-indication, sans indicateur coloré).
\item Demi-équations en milieu acide :
\[ \ce{MnO4- + 8 H+ + 5 e- <=> Mn^2+ + 4 H2O} \qquad \ce{Fe^2+ <=> Fe^3+ + e-} \]
\item On multiplie la seconde par 5 et on additionne :
\[ \ce{MnO4- + 5 Fe^2+ + 8 H+ -> Mn^2+ + 5 Fe^3+ + 4 H2O} \]
\item L'\cle{équivalence} est atteinte lorsque les réactifs ont été mélangés dans les proportions stœchiométriques de la réaction de dosage :
\[ \frac{n_E(\ce{MnO4-})}{1} = \frac{n_i(\ce{Fe^2+})}{5} \quad \Longleftrightarrow \quad n_i(\ce{Fe^2+}) = 5\, C_2 V_{2E} \]
\item Calculs :
\[ n_E(\ce{MnO4-}) = C_2 V_{2E} = 0{,}020 \times 16{,}0\times 10^{-3} = \SI{3,2e-4}{mol} \]
\[ n_i(\ce{Fe^2+}) = 5 \times \SI{3,2e-4}{mol} = \SI{1,6e-3}{mol} \]
\[ C_1 = \frac{n_i(\ce{Fe^2+})}{V_1} = \frac{\SI{1,6e-3}{mol}}{\SI{20,0e-3}{L}} = \SI{0,080}{mol.L^{-1}} \]
Concentration massique en fer :
\[ t = C_1 \times M(\ce{Fe}) = 0{,}080 \times 55{,}8 = \SI{4,46}{g.L^{-1}} \]
\item $t = \SI{4,46}{g.L^{-1}} \gg \SI{0,20}{g.L^{-1}}$ : la teneur en fer dépasse largement la norme. \textbf{Cette eau n'est pas conforme} pour la boisson ; elle devra être traitée (aération puis filtration, par exemple) avant consommation.
\end{enumerate}
\textbf{Barème indicatif (sur 10) :} schéma légendé et repérage de l'équivalence 2 pts ; demi-équations 2 pts ; équation-bilan 1,5 pt ; définition et relation à l'équivalence 1,5 pt ; calculs de $C_1$ et $t$ 2 pts ; conclusion 1 pt.\\
\textbf{Points de vigilance :} ne jamais ajouter d'acide chlorhydrique (les ions \ce{Cl-} seraient oxydés par \ce{MnO4-}) : on acidifie à l'acide sulfurique ; bien convertir les volumes en litres ; ne pas oublier le facteur 5 de la stœchiométrie.
\end{corrige}

\begin{corrige}[Exercice 11 — Type Bac : la tôle galvanisée et la rouille]
\begin{enumerate}
\item Dans le fer métallique \ce{Fe} (corps simple), n.o.(Fe) $= 0$. Dans \ce{Fe2O3} : $2x + 3(-2) = 0 \Rightarrow x = +3$. Le nombre d'oxydation du fer \emph{augmente} de $0$ à $+3$ : la formation de la rouille est donc une \textbf{oxydation} du fer (le fer joue le rôle de réducteur face au dioxygène de l'air).
\item Demi-équations (milieu neutre/basique, corrosion humide) :
\[ \ce{Fe -> Fe^2+ + 2 e-} \qquad \ce{O2 + 2 H2O + 4 e- -> 4 OH-} \]
En multipliant la première par 2 et en additionnant, les ions \ce{Fe^2+} et \ce{OH-} forment l'hydroxyde de fer(II) :
\[ \ce{2 Fe + O2 + 2 H2O -> 2 Fe(OH)2} \]
Cet hydroxyde est ensuite oxydé et déshydraté en rouille \ce{Fe2O3} hydratée.
\item Le zinc est un \textbf{réducteur plus fort que le fer} : c'est donc lui qui cède préférentiellement ses électrons et qui est oxydé à la place du fer. Si des ions \ce{Fe^2+} se forment localement au niveau de la rayure, le zinc les réduit immédiatement :
\[ \ce{Zn + Fe^2+ -> Zn^2+ + Fe} \]
Le fer est ainsi régénéré et reste protégé.
\item Le zinc est qualifié d'\cle{anode sacrificielle} car il s'oxyde (se « sacrifie ») à la place du fer : tant qu'il reste du zinc en contact électrique avec le fer, c'est le zinc qui est consommé par la corrosion et le fer demeure intact, même à nu au niveau d'une rayure. Sur une tôle de fer \emph{nue}, rien ne s'oppose à l'oxydation du fer par \ce{O2} en présence d'eau : la rouille se développe directement.
\item Méthodes de protection utilisées au quotidien au Cameroun :
\begin{itemize}
\item la \textbf{peinture} des portails, lits et ferronneries : elle forme une barrière étanche qui empêche le contact du fer avec l'air humide et l'eau ;
\item le \textbf{graissage ou huilage} des outils, houes et machettes : le film gras isole le métal de l'humidité ;
\item le \textbf{séchage et le stockage à l'abri de la pluie} : supprimer l'eau, élément indispensable à la corrosion.
\end{itemize}
Principe commun : empêcher le contact simultané du fer avec le dioxygène et l'eau, ou interposer un métal plus réducteur qui s'oxyde à sa place.
\end{enumerate}
\textbf{Barème indicatif (sur 10) :} nombres d'oxydation et conclusion 2 pts ; demi-équations et bilan de la corrosion 2,5 pts ; comparaison des pouvoirs réducteurs et équation 2 pts ; notion d'anode sacrificielle 2 pts ; méthode de protection expliquée 1,5 pt.\\
\textbf{Points de vigilance :} bien distinguer la corrosion (oxydation du fer) de la protection cathodique (le métal protégé devient le siège de la réduction) ; dans le bilan de corrosion en milieu humide, utiliser \ce{OH-} et non \ce{H+} car le milieu n'est pas acide.
\end{corrige}

\newpage

\coursec{Fiche bilan}

\begin{popfigure}
\begin{center}
\begin{tikzpicture}[
  mindmap,
  grow cyclic,
  every node/.style={concept, font=\small\bfseries, text=white},
  root concept/.append style={concept color=popdark, font=\large\bfseries},
  level 1/.append style={level distance=3.6cm, sibling angle=60, font=\footnotesize},
  level 2/.append style={level distance=2.4cm, sibling angle=45, font=\scriptsize},
  concept connection/.append style={color=popmuted}
]
\node[root concept]{Oxydoréduction spontanée}
  child[concept color=popblue]{ node {Acide + métal}
    child { node {\ce{Zn + 2H+ -> Zn^2+ + H2}} }
    child { node {Métaux nobles inertes} }
  }
  child[concept color=poporange]{ node {Ion métallique + métal}
    child { node {\ce{Cu^2+ + Zn -> Cu + Zn^2+}} }
    child { node {Règle du gamma} }
  }
  child[concept color=popgreen]{ node {Point de vue électronique}
    child { node {Oxydation = perte de \ce{e-}} }
    child { node {Réduction = gain de \ce{e-}} }
    child { node {Couple Ox/Red} }
  }
  child[concept color=poppurple]{ node {Réactions en solution}
    child { node {\ce{MnO4-}/\ce{Mn^2+}} }
    child { node {\ce{Cr2O7^2-}/\ce{Cr^3+}} }
    child { node {\ce{I2}/\ce{I-} et \ce{S4O6^2-}/\ce{S2O3^2-}} }
  }
  child[concept color=poppink]{ node {Nombres d'oxydation}
    child { node {Règles de calcul} }
    child { node {Oxydation : n.o. augmente} }
    child { node {Équilibrage} }
  }
  child[concept color=popgold]{ node {Voie sèche}
    child { node {\ce{2Fe + 3Cl2 -> 2FeCl3}} }
    child { node {Rouille du fer} }
  };
\end{tikzpicture}
\end{center}
\legende{Carte mentale de la leçon : les six branches de l'oxydoréduction spontanée.}
\end{popfigure}

\begin{bilan}[L'essentiel de la leçon]
\textbf{Définitions clés (point de vue électronique)}
\begin{itemize}
  \item \cle{Oxydation} : perte d'électron(s) ; \cle{réduction} : gain d'électron(s).
  \item \cle{Oxydant} : espèce capable de capter des électrons ; \cle{réducteur} : espèce capable de céder des électrons.
  \item \cle{Couple redox} Ox/Red : $\text{Ox} + n\,\ce{e-} \ce{<=>} \text{Red}$. Toute oxydoréduction met en jeu \textbf{deux} couples : le réducteur d'un couple cède ses électrons à l'oxydant de l'autre.
\end{itemize}

\textbf{Définitions clés (point de vue des nombres d'oxydation)}
\begin{itemize}
  \item Oxydation : le \cle{nombre d'oxydation} (n.o.) \textbf{augmente} ; réduction : le n.o. \textbf{diminue}.
  \item Oxydant : espèce dont un élément voit son n.o. diminuer ; réducteur : espèce dont un élément voit son n.o. augmenter.
  \item Une réaction est une oxydoréduction si le n.o. d'au moins un élément varie.
\end{itemize}

\textbf{Règles de calcul du n.o.}
\begin{itemize}
  \item Corps simple : n.o. $= 0$ ; ion monoatomique : n.o. $=$ charge de l'ion.
  \item H : $+1$ (sauf hydrures, $-1$) ; O : $-2$ (sauf peroxydes, $-1$).
  \item Molécule neutre : somme des n.o. $= 0$ ; ion polyatomique : somme $=$ charge de l'ion.
\end{itemize}

\textbf{Équations-types à connaître par c\oe ur}
\begin{center}
\rowcolors{2}{popblueL}{white}
\begin{tabularx}{\linewidth}{|l|X|}
\hline
\rowcolor{popblue!40} \textbf{Situation} & \textbf{Équation-bilan} \\
\hline
Acide + zinc & \ce{Zn + 2H+ -> Zn^2+ + H2} \\
\hline
Acide + fer & \ce{Fe + 2H+ -> Fe^2+ + H2} \\
\hline
Ion cuivre + zinc & \ce{Cu^2+ + Zn -> Cu + Zn^2+} \\
\hline
Permanganate + fer(II) (milieu acide) & \ce{MnO4- + 5Fe^2+ + 8H+ -> Mn^2+ + 5Fe^3+ + 4H2O} \\
\hline
Dichromate + éthanol & \ce{Cr2O7^2- + 3C2H5OH + 16H+ -> 2Cr^3+ + 3CH3COOH + 11H2O} \\
\hline
Diiode + thiosulfate & \ce{I2 + 2S2O3^2- -> 2I- + S4O6^2-} \\
\hline
Voie sèche (fer + dichlore) & \ce{2Fe + 3Cl2 -> 2FeCl3} \\
\hline
\end{tabularx}
\end{center}

\textbf{Demi-équations de référence}
\begin{itemize}
  \item \ce{MnO4- + 8H+ + 5e- <=> Mn^2+ + 4H2O} ; \ce{Cr2O7^2- + 14H+ + 6e- <=> 2Cr^3+ + 7H2O}
  \item \ce{I2 + 2e- <=> 2I-} ; \ce{S4O6^2- + 2e- <=> 2S2O3^2-} ; \ce{Fe^3+ + e- <=> Fe^2+}
\end{itemize}

\textbf{Méthodes express}
\begin{itemize}
  \item \emph{Équilibrer une demi-équation} : (1) équilibrer l'élément autre que O et H ; (2) équilibrer O avec \ce{H2O} ; (3) équilibrer H avec \ce{H+} ; (4) équilibrer les charges avec \ce{e-}.
  \item \emph{Équation-bilan} : multiplier les demi-équations pour égaliser les électrons échangés, additionner, simplifier. Les électrons ne doivent \textbf{jamais} apparaître dans le bilan.
  \item \emph{Prévoir une réaction} (règle du gamma) : la réaction spontanée a lieu entre l'oxydant le plus fort et le réducteur le plus fort (réactifs placés en « gamma » sur l'échelle des couples).
  \item \emph{Métaux et acides} : seuls les métaux dont le couple est au-dessus de \ce{H+}/\ce{H2} (Zn, Fe, Al...) sont attaqués avec dégagement de \ce{H2} ; Cu, Ag, Au sont inertes.
\end{itemize}
\end{bilan}

\begin{aretenir}[Checklist avant le Bac]
\begin{itemize}
  \item $\square$ Je sais définir oxydation, réduction, oxydant, réducteur et couple redox des deux points de vue (électronique et n.o.).
  \item $\square$ Je sais écrire la demi-équation électronique d'un couple Ox/Red en milieu acide.
  \item $\square$ Je sais équilibrer \ce{MnO4-}/\ce{Mn^2+} et \ce{Cr2O7^2-}/\ce{Cr^3+} sans hésiter.
  \item $\square$ Je sais combiner deux demi-équations pour obtenir l'équation-bilan sans électrons.
  \item $\square$ Je connais par c\oe ur les réactions \ce{MnO4-} + \ce{Fe^2+}, \ce{Cr2O7^2-} + éthanol et \ce{I2} + \ce{S2O3^2-}.
  \item $\square$ Je sais prévoir si un acide dilué attaque un métal donné (Zn, Fe, Al oui ; Cu, Ag, Au non).
  \item $\square$ Je sais prévoir la réaction entre un ion métallique et un métal grâce à la règle du gamma.
  \item $\square$ Je sais calculer le nombre d'oxydation d'un élément dans une molécule ou un ion.
  \item $\square$ Je sais repérer une oxydoréduction en suivant la variation des n.o. (y compris par voie sèche).
  \item $\square$ Je sais identifier les produits d'une réaction : dégagement de \ce{H2} (test à la flamme), changements de couleur (\ce{MnO4-} violet $\to$ incolore, \ce{Cr2O7^2-} orange $\to$ \ce{Cr^3+} vert).
  \item $\square$ Je vérifie toujours la conservation des atomes \textbf{et} des charges dans mes équations finales.
  \item $\square$ Je sais expliquer la formation de la rouille comme une oxydation du fer par le dioxygène en milieu humide.
\end{itemize}
\end{aretenir}

\findecours

\end{document}
